\documentclass[../main.tex]{subfiles}
\begin{document}
%==========================================TIÊU ĐỀ TẬP========================================
\begin{table}[h]
    \begin{adjustwidth}{-1cm}{}
        \begin{tabular}{>{\centering\arraybackslash}p{6cm}|>{\raggedright\arraybackslash}p{14cm}}
            \multirow{5}{6cm}
            % Cột bên trái: ÉPISODE và số tập
            {\\[-40pt]
                \flaregothic\fontsize{20pt}{20pt}\selectfont \centering
                \color{eptype}HISTOIRE\color{black}\\
                \gangofthree\fontsize{72pt}{72pt}\selectfont
                \color{epnum}3
                \\[18pt]
                % \flaregothic\fontsize{14pt}{14pt}\selectfont
                % \color{epattr}BIENVENUE, \uline{\color{Banpire} Mel} !
                % \flaregothic\fontsize{18pt}{18pt}\selectfont
                % \color{epattr}ÉPISODE PERDU
                \color{black}
            }
             &
            % Dòng thứ nhất: tên tập tiếng Nhật
            {\fotskip\fontsize{18pt}{18pt}\selectfont \ruby{大}{だい}\ruby{掃}{そう}\ruby{除}{じ}だ！\ruby{新}{しん}\ruby{年}{ねん}を\ruby{迎}{むか}えよう！}
            \\[6pt]
            % Dòng thứ hai: tên tập tiếng Anh
             & {\cafeteria\fontsize{18pt}{18pt}\selectfont Cleaning Day! Preparing for the New Year!} \\[4pt]
            % Dòng thứ ba: tên tập tiếng Việt
             & {\vnmsans\fontsize{18pt}{18pt}\selectfont Dọn dẹp nhà cửa, chuẩn bị đón Tết thôi!}     \\[8pt]
            % Dòng thứ tư: Nhân vật xuất hiện
             & {\brian{\fontsize{14pt}{14pt}\selectfont Track used: Lunar Zoo Year - Loonboon}}
        \end{tabular}
    \end{adjustwidth}
\end{table}
%===========================================NỘI DUNG==========================================
\color{black}

Để thuận tiện hơn cho việc dọn dẹp nhà cửa, hơn hai mươi người trong gia đình Hikawa cùng các thành viên của \hll\ sẽ tự phân chia nhau để làm việc. Theo đó, từ bảng dưới đây, ta có:

\begin{table}[h!]
    \centering
    \begin{tabular}{|c|l|l|}
        \hline
        \textbf{STT} & \textbf{Thành viên làm}             & \textbf{Công việc phải làm}    \\ \hline
        1            & Roboco, Noel, Haato                 & Quét, lau nhà phòng của Kuon   \\ \hline
        2            & Shion, Subaru                       & Dọn phòng của Hijaki           \\ \hline
        3            & Matsuri, Sora, Kuon, Ryujin, A-chan & Đi mua đồ                      \\ \hline
        4            & Okayu, Korone                       & Sơn lại phòng các tầng 3, 4, 5 \\ \hline
        5            & Mio, Choco, Ayame                   & Lau dọn nhà bếp                \\ \hline
        6            & Miko, Mel                           & Lau dọn bàn thờ                \\ \hline
        7            & Kaien, Ryuuhou, Aki                 & Giặt kiêm phơi quần áo         \\ \hline
        8            & Fubuki, Rushia, Aqua                & Lau dọn tầng 3                 \\ \hline
        9            & Tất cả mọi người                    & Dọn kho, lau cầu thang         \\ \hline
    \end{tabular}
    \caption{Bảng phân công nhiệm vụ dọn dẹp nhà Hikawa}
\end{table}

\vspace{-40pt}
\subsubsection*{\phantomsection\label{ep-003.IVa}}
\addcontentsline{toc}{subsubsection}{Roboco, Noel, Haato}

\begin{center}
    {\color{T1} Nhóm 1 - Roboco, Noel, Haato}\\Phòng của Kuon.
\end{center}

Roboco và Noel đứng trước căn phòng của Kuon, hai cô ngốc cùng nhau suy nghĩ về cách làm sao để dọn dẹp hiệu quả. Căn phòng đầy những thiết bị điện tử, từ dàn âm thanh đến các thiết bị máy tính và TV, khiến mọi thứ trông có vẻ bừa bộn dù thực tế nó khá ngăn nắp.

\begin{dialogue}
    \speak{Roboco} Hmm\dots\ Làm gì bây giờ nhỉ?
    \speak{Noel} Xem chừng phòng của họ kha khá nhiều đồ đây\dots
\end{dialogue}

Haato - người canh chừng hai cô gái ngố này - đứng gần đó, nở một nụ cười tự tin: ``Có gì đâu, chỉ cần quét lau nhà đến khi sạch là ổn mà. Hai người chưa bao giờ quét dọn nhà mình bao giờ à?''

Hai cô nàng nhìn nàng Song tính, rồi nhìn nhau, lắc đầu.

\begin{dialogue}
    \speak{Roboco} Bình thường chị không dọn dẹp cặn kẽ như thế này\dots
    \speak{Noel} Em cùng lắm cũng chỉ bê vác đồ xung quanh thôi, chứ dọn nhà thì chưa bao giờ làm\dots
\end{dialogue}

Nghe câu trả lời của hai người, Haato thở dài, rồi chống tay vào hông. ``Thôi được rồi, bây giờ thế này, một người lau nhà, một người quét nhà, được chưa? Tôi cũng sẽ phụ giúp hai người mà.''

Cả hai cô nàng cười khì, tỏ vẻ đồng tình với kế hoạch này từ cô nàng Thế hệ thứ Nhất: ``Ý tưởng hay đó!''

Haato mỉm cười, khẽ lắc đầu khiêm nhường: ``Cái này hai người có thể tự làm được mà nhỉ.''

Nàng Hiệp sĩ chỉ tay về phía mấy đồ lỉnh kỉnh ở nhà Kuon. Chúng bao gồm hai tủ quần áo, tất cả các đồ chất đống trên đó, một tủ kính đựng đồ lưu niệm, kệ vô tuyến và dàn loa. ``Em có cần vác hết đống đồ này không?''

``Làm gì?'' nàng Song tính nhíu mày. Roboco cười khẽ, đáp với tông giọng lém lỉnh: ``Để chị quét nhà hút bụi chứ sao nữa\~{}''

Haato nghe lý do từ nàng Người máy, tỏ vẻ không chắc chắn: ``Em cũng không biết nữa\dots\ Chắc Kuon-senpai không cho phép đâu\dots''

Nhưng Noel vẫn khá cứng đầu, cô tiếp tục thuyết phục bề trên của mình bằng một lý lẽ không thể cãi vào đâu được: ``Có khi thấy căn phòng không có một hạt bụi thì chắc chắn anh ấy sẽ vui lắm đấy.''

Nàng Người máy gật đầu: ``Với cả anh ấy muốn chúng ta dọn sạch còn gì nữa.''

Sau khi được thu phục bởi hai cô gái, cô cảm thấy họ nói đúng. Cô nhìn họ với vẻ mặt tràn đầy quyết tâm: ``Hai người nói cũng có lý. Thôi được rồi, triển đi!''

Noel và Roboco lập tức đức nghiêm, chào kiểu quân đội: ``Yes, Madam! (Vâng, thưa bà!)''

Vừa dứt lời, nàng Hiệp sĩ liền vào nhà vệ sinh (hình b) tìm chai nước lau nhà và cái chổi lau nhà. Thấy có cái chổi đặt trên một cái chậu có cối giũ, cô nghĩ bụng: ``Chắc cái chổi là cái này,'' và xách nó ra ngoài.

``Kể ra em cung đoán được cái chổi lau nhà ha,'' Haato đứng nhìn cô với ánh mắt nghiêm túc. Noel nghe vậy, liền cười gượng: ``C-Chị cứ quá khen\dots''

Nàng Song tính hỏi tiếp, "Vậy em có biết pha nước để lau nhà như thế nào không?" Nghe đến đây, nàng Hiệp sĩ ngơ ngác ra, không biết phải trả lời sao.

Haato mỉm cười, chỉ vào chậu nước gần đó và hướng dẫn: ``Cứ lấy một chậu nước gần đầy vào, rồi lấy khoảng một, hai nắp nước lau nhà là được.''

``À, ra thế. Vậy thì dễ mà,'' Noel gật đầu, chạy thẳng vào phòng vệ sinh tìm nước lau nhà. Không lâu sau, cô tìm thấy chai nước lau nhà trên kệ.

Cô với tay lên, lấy được chai nước nhưng không may, chỉ mới chạm nhẹ vào cái kệ, lập tức tất cả đồ vật trên đó cùng với cái kệ rơi bẹp xuống đất, đồ đạc rơi loảng xoảng, khiến cho cô giật mình.

Roboco từ phòng khách lao vào: ``Có chuyện gì thế!?'' Haato thì ngỡ ngàng nhìn cảnh phòng tắm còn gọn gàng ban nãy, giờ đã trở nên lộn xộn. Cô nhìn Noel, hỏi: ``Em vừa làm cái gì vậy!?''

Nàng Hiệp sĩ run run: ``E-Em chỉ mới lấy cái chai nước lau nhà thôi mà\dots'' Nhưng nàng Song tính lại tỏ vẻ không tin. Cô chỉ vào kệ đồ bị đổ ụp xuống và gặng hỏi: ``Thế tại sao đồ vật rơi xuống đất thế!?''

Cô nàng ngây ra, tỏ vẻ khó hiểu không biết chuyện gì vừa xảy ra: ``Em có biết đâu, mới đụng nhẹ vào cái kệ thôi mà\dots''

Nàng Người máy tiến tới trấn an hậu bối của mình đang rơm rớm nước mắt kia: ``Thôi, không sao, tí nữa chúng ta sẽ lắp lại sau. Bây giờ cứ tập trung lau dọn phòng khách đã.''

\ct

Noel khiêng chiếc chậu nước và cây chổi lau nhà ra giữa căn phòng, trong khi Haato và Roboco phụ khiêng những đồ đạc lớn để lấy chỗ lau nhà.

Nàng Hiệp sĩ bắt đầu pha nước lau nhà theo đúng tỷ lệ mà nàng Song tính đã hướng dẫn. Mọi chuyện sẽ không có gì đáng nói\dots

\dots cho đến khi cô sử dụng cây lau nhà và\dots\ nhúng liên tiếp vào cái chậu nước với một lực cực mạnh, khiến cho nước bắn tung tóe khắp nơi.

Nàng ``Quản lý lâm thời '' của họ định lên tiếng, nhưng Roboco đã giải thích ngay cho cô bằng một nụ cười trừ: ``Thì, đằng nào ta chẳng phải lau nhà mà.''

Nàng Song tính tỏ vẻ khó chịu: ``Biết là thế, nhưng ta phải quét qua đã chứ!'' Nàng Người máy tỏ vẻ ngạc nhiên: ``Ể\~{}? Tại sao lại như thế?''

Cô nhìn hai ``tên ngốc'', rồi nói với giọng nghiêm túc: ``Nhỡ đâu trong phòng còn rác thì sao? Chưa kể còn bụi bặm các thứ nữa chứ.''

Nàng Hiệp sĩ gật gù: ``Chị nói cũng phải\dots'' nói rồi, cô dựng cây lau nhà và dựng nước ở đúng vị trí mà cô làm đổ nước lau nhà trước đó.

Haato thúc giục hai người: ``Đi nào, mau dọn bớt đồ lên trên cao, đệm thì gập lên rồi dựng tường.''

Dứt lời, nàng tóc vàng bắt đầu cầm chổi quét qua cái sàn, trong đó có sót vài mảnh giấy và gói bánh ở tận bên trong ngóc ngách phòng. Cô quét đi, nhưng ánh mắt không quên hướng về hai người bạn đang làm việc của mình.

Cô ngước nhìn hai cô nàng vẫn đang đứng đó, hầu như không làm gì sau khi họ chất xong đống đồ đạc lên cao, rồi nói trong cơn thở dài: ``Không định phụ tôi một tay à hai người?''

Ngay tức khắc, Roboco thiết lập một cái máy hút bụi từ bàn tay trái của mình. Về cơ bản là giống với bàn tay bình thường của cô, chỉ khác là có một cái lỗ hút ngay trên lòng bàn tay. Cả Noel và Haato nhìn món ``bảo bối'' của cô trong sự trầm trồ.

\begin{dialogue}
    \speak{Noel} Oa\~{}!
    \speak{Roboco} Thế nào? Thấy chiếc máy hút bụi mini này của chị chứ? Nó có thể hút bụi từ mọi ngóc ngách đấy.
\end{dialogue}

Như không muốn để nàng Người máy dài dòng, Haato đã vội lên tiếng: ``Gì cũng được! Miễn là nó có thể giúp bọn mình công việc này! Nhanh lên đi cái!''

Nói là làm, cô hướng lòng bàn tay lên trần nhà và kích hoạt nó. Một luồng gió siêu mạnh được tạo ra từ bàn tay và hút toàn bộ mọi thứ, từ mạng nhện, bụi bẩn đến cả\dots\ vữa tường đang lủng lẳng trên đó.

``Nhưng\dots\ chắc sẽ phải làm lại mấy chỗ trên tường đấy,'' nàng Người máy chớp mắt, nhìn những mảng trần đang rơi xuống, có lẽ là do lực hút của chiếc máy.

Chỉ trong chốc lát, cả trần nhà đã sạch như mới, nhưng đánh đổi lại là\dots\ vài chỗ lóc chóc trên tường. Mảng tường bị hút một phần nhỏ làm cho không gian có vẻ hơi khác biệt, nhưng không ai để ý.

Nhưng, chưa kịp được Haato và Noel khen, bản tính ``PON'' của cô lại ``trỗi dậy,'' khi cô vô thức hướng bàn tay còn đang bật máy hút về phía hai người bạn.

Và, dĩ nhiên, hậu quả là\dots\ nó hút luôn cả quần áo của hai người, chỉ còn lại đúng đồ lót trên người họ. Cả hai cô gái đỏ mặt như quả gấc chín, vội vàng che đi cả thân mình.

\begin{dialogue}
    \speak{Noel} \textbf{Không\~{}!!}
    \speak{Haato} \textbf{C-Chị làm cái gì thế hả!?}
\end{dialogue}

Nhìn thấy hai cô bạn đang ngượng ngùng, Roboco tỏ vẻ bối rối, vội vàng tắt máy hút đi, nhưng vẫn cười khì: ``Xin lỗi hai đứa nha\~{}, chị quên chưa tắt máy\dots''

Nhưng lời xin lỗi của cô không làm cho nàng Song tính nguôi ngoai cơn giận. ``\textbf{Cười cái gì chứ! Mau trả lại quần áo cho bọn em đây!}''

Trong khi đó, nàng Hiệp sĩ vội vàng nhìn ngó xung quanh, tay vẫn cố gắng che đi ngực. Rất may là xung quanh họ không có ai, ngoài bản thân họ ra.

``Chị bảo đảm đi!'' Haato quấn quanh người bên chiếc rèm, mặt của cô vẫn còn đỏ như quả cà chua. ``Nếu còn cái trò này lần nữa thì đừng trách em đấy!''

Nàng Người máy chỉ biết gãi đầu. Cô đưa quần áo vẫn còn nguyên vẹn từ trong bụng của mình ra, rồi gãi đầu cười trừ: ``Lần sau chị sẽ cẩn thận hơn. Thật đấy!''

Nàng Song tính bực bội, trong đầu cô thầm nghĩ: ``\textit{Kuon-senpai nói đúng\dots\ Chị ấy đúng là PON thật\dots}''

\ct

Sau màn ``hút sạch'' quần áo bất đắc dĩ, ba cô gái tiếp tục công việc dọn dẹp. Họ chuyển sang lau dọn đằng sau mấy cái tủ quần áo và bộ dàn điện tử.

Nàng Hiệp sĩ hăng hái xung phong bê hết đống đồ ra chỗ khác. Với các tủ quần áo và tủ kính, mọi thứ diễn ra dễ dàng, nhưng đến lượt kệ đồ điện tử thì\dots

``Cẩn thận nha, đằng sau có nhiều dây điện lắm đấy,'' Haato đứng ở cuối phòng cảnh báo cho cô.

Cô nhìn đống dây điện chằng chịt ở đằng sau, liền gật gù: ``Có khi chúng ta dỡ từng món đồ trên kệ xuống trước rồi mới bê cả cái tủ đi nhỉ?'' Nàng Rô-bốt cũng phân tích tình hình, rồi tỏ vẻ tán thành với cô: ``Cũng hợp lý. Để chị phụ nhé.''

Họ bắt tay tháo dỡ từng món: thùng loa, ti-vi, và dàn âm thanh. Từng món được đặt ngay dưới chân tủ quần áo họ đã kê tạm ở góc phòng. Nhưng khi vừa dọn đến phía sau, bụi lập tức bay mù mịt, làm cả nàng Hiệp sĩ lẫn nàng Song tính hắt hơi liên tục.

\begin{dialogue}
    \speak{Noel} Nhiều bụi qua\dots\ *khụ khụ*!
    \speak{Haato} *khụ khụ* Lại còn nhiều rác nữa chứ! Chỗ này chắc phải cả năm chưa dọn ấy!
    \speak{Noel} Để em kiếm cái hót rác\dots
\end{dialogue}

Trông thấy mớ hổ lốn ở đằng sau kệ tủ, Roboco khoanh tay, nghi ngợi một lúc. Sau đó, mắt cô lại lóe lên, bất giác chìa tay ra trước: ``Thôi để chị tạo ra cái máy hút to hơn, cho nhanh.''

Chưa kịp để hai cô gái nói thêm một lời nào, cánh tay của nàng Người máy bỗng chốc biến đổi, tạo ra một chiếc máy hút bụi khổng lồ, to cỡ bằng nửa cái bàn tay.

``\textbf{Hút!}''

Chiếc máy hút từ tay cô khởi động, tạo ra một luồng gió cực mạnh. Chẳng mấy chốc, toàn bộ bụi và rác vụn vặt đã bị hút sạch nhẵn. Rút kinh nghiệm từ màn ``hút quá đà'' ban nãy, cô nhanh chóng tắt máy trước khi đưa cánh tay trở lại bình thường.

Với phần kệ tủ không còn bụi rác, cô chống tay đắc chí như một đứa trẻ muốn được khen: ``Thế nào? Chị cũng giúp ích chứ?''

Noel há hốc mồm, tỏ vẻ ngạc nhiên tột độ: ``\dots Thật tuyệt vời! Cái cánh tay này\dots\ đúng là lợi hại quá đi!'' Trong khi đó, Haato thì vừa thở dài nhẹ nhõm vì cô tưởng Roboco sẽ lại gây chuyện, vừa chống hông nhìn: ``Thôi được. Xem như lần này chị giúp được một việc\dots\ Nhưng!''

``Nhưng sao ạ?'' nàng Hiệp sĩ vội tiếp lời. Nàng Song tính nghiêm mặt, tay chỉ vào đống kệ: ``Chúng ta vẫn còn cái đống dây điện loằng ngoằng kia kìa! Ai tháo được thì tháo giùm tôi cái đã!''

\ct

Noel cầm chổi lau nhà, bắt đầu dọn dẹp sàn nhà, cũng như lau dọn bên trong tủ tồ và kệ điện tử. Hai người còn lại cũng phụ giúp, mỗi người cầm một chiếc chổi xuất hiện từ\dots\ hư không.

Thấy nàng Người máy đang lau nhà một cách bình thường đến ngạc nhiên, nàng Hiệp sĩ hỏi: ``Ủa, Roboco-senpai? Chị không định tạo cái chổi lau nhà từ cánh tay chị ạ?''

Roboco lắc đầu và cười khì: ``Chị vẫn chưa có ý định làm thứ đó\dots\ Xin lỗi nha.''

Đên cả Haato cũng tỏ vẻ bất ngờ, nhất là với cô nàng có những ý tưởng điên rồ. Cô nghiêng đầu, vừa hỏi vừa tưởng tượng: ``Mà nếu chị tạo ra chổi lau nhà, chị định lau kiểu gì? Sử dụng tay còn lại của mình à?''

Nàng Người máy tỏ vẻ nghiêm túc hơn hẳn: ``Không, chị sẽ dùng chính cái tay đó làm chổi lau luôn.''

Nàng Song tính dần hình dung ra ``bản thiết kế'' mà tiền bối của cô nói, rồi khẽ rùng mình lắc đầu: ``Nghe không khả thi chút nào\dots''

Một khoảng thời gian vắng lặng hiếm hoi, khi chỉ có xoèn xoẹt của chổi lau nhà, tiếng lạch cạch của chậu nước và tiếng xùm xụp của chậu rửa.

\ct

Lau nhà được một lúc, họ gần như hoàn thành công việc. Duy chỉ còn một vết bẩn cứng đầu ở góc phòng mà ba người cố gắng lau mãi không sạch.

Cả ba cô gái đưa ra ý kiến để giải quyết triệt để. Bắt đầu với nàng Song tính, ``quản lý lâm thời'' của hai người: ``Hay thử dùng khăn ẩm lau đi? Biết đâu lại được?''

Tuy nhiên, hai người kia có những ý tưởng\dots\ táo bạo hơn. Noel thì giơ cây chùy giắt ở hông lên, mắt sáng quắc: ``Hay em dùng cây chùy này đập vết bẩn đi nhỉ?''

Lập tức, cô bị một Haato hoảng hốt can ngăn: ``Em định phá nhà người ta hả!?''

Đáp lại, Noel cười một cách vô tư: ``Nếu không rửa sạch được thì cứ đập bỏ đi thôi\~{}'' Nàng Song tính bức xúc, vội ôm chặt lấy tay của cô nàng: ``Không ai làm thế cả! Em bị ấm đầu rồi à!?''

Trong khi hai người tranh luận, Roboco âm thầm rời đi và mang vòi nước từ nhà vệ sinh ra, mạnh dạn xịt thẳng vào vết bẩn. Với áp lực nước mạnh, vết bẩn cứng đầu cuối cùng cũng chịu ``bay màu''.

``Xong rồi nè\~{}''

Nhìn vết bẩn ấy đã không cánh mà bay, cả hai cô gái liền đứng ngây người ra, như thể không tin vào mắt họ.

Nàng Song tính thở dài: ``\dots Hoặc là cách này cũng được.'' Nàng Hiệp sĩ thì gãi đầu cười trừ: ``Sao em không nghĩ ra nhỉ?''

Thấy hai người không có ý kiến gì, nàng Người máy tỏ ra hào hứng: ``Đã thế thì chị sẽ xịt nước khắp phòng luôn, cho sạch sẽ hẳn!''

\dots và tất nhiên, bản chất ``PON'' của cô tiếp tục được trỗi dậy.

Noel vỗ tay cổ vũ tiền bối, còn Haachama giật mình, vội ngăn cản: ``Đừng làm vậy! Chị mà xịt khắp nơi thì nước sẽ bắn vào đồ điện tử mất!''

Nhưng, như thể không nghe thấy lời cảnh báo, Roboco liền xịt nước thẳng lên trần nhà, để dòng nước nhỏ xuống dàn điện tử bên dưới. Đỉnh điểm, cô còn xịt vào ổ cắm điện, khiến tia lửa lập lòe rồi phát ra một tiếng ``bép'' lớn. Khói đen từ ổ điện bốc lên, mùi khét lẹt lan khắp phòng.

Haato nổi đóa, tay chỉ về ổ cắm điện: ``Đó, em đã bảo rồi mà! Chị xem giờ phải làm sao đây!?'' Roboco thì tỏ vẻ ngơ ngác, nghiêng đầu nhìn ổ điện đang xì khói đen: ``Ủ-Ủa\~{} Sao lại thế nhỉ?''

``Sao với trăng gì nữa! Kuon-senpai về mà nhìn thấy cảnh này sẽ giết chúng ta mất!'' mặt của Haachama xanh hẳn đi.

Dường như mọi thứ vẫn chưa kết thúc, Noel sau đó đã đổ thêm dầu vào lửa, khi cô bước lên, cầm lấy cây chùy trên tay và nói với giọng điệu tự tin: ``Cứ để đó cho em! Chuyện nhỏ\~{}'' và lập tức đập thẳng vào ổ điện mà không đợi ai lên tiếng. Chiếc ổ điện gắn tường thì vỡ tanh bành, vỡ luôn cả một mảng tường.

``\textbf{Má trẻ này nữa! Đừng làm mọi chuyện rối hơn chứ!!}'' cô ``quản lý'' ôm đầu, càng lúc càng tỏ vẻ hoảng loạn hơn.

\ct

Cuối cùng, sau khi gây họa, cả Roboco và Noel phải rút tiền túi để sửa ổ điện cho gia chủ. Họ cùng đồng ý sẽ chỉ sử dụng cây lau nhà như cách thông thường, không dám áp dụng bất kỳ ``chiêu trò sáng tạo'' nào nữa.

Cả căn phòng giờ đã sáng sủa hơn hẳn, không còn lấy một hạt bụi, cũng không còn rác, như thể căn phòng này vừa được xây mới.

\dots Ngoại trừ trên bức trường trông cũ kỹ kia còn những vết bong tróc và một ổ điện bị vỡ vụn.

\begin{dialogue}
    \speak{Haato} *lau trán, thở phào* Vậy là xong rồi đấy! Giờ ta chỉ cần kê nốt đống đồ này đi, rồi nghỉ!
    \speak{Noel} *gãi đầu* Vâng\dots Lần sau em sẽ bớt sáng tạo hơn\dots
    \speak{Roboco} Vậy hẹn lần sau chị nghĩ ra cái gì thú vị hơn nhé\~{}
\end{dialogue}

Thế nhưng, ngay khi ba cô gái chuẩn bị bắt tay vào việc, họ nghe tiếng bước chân rầm rập từ cầu thang. Hai thành viên \textit{GAMERS}, Korone và Okayu, xuất hiện với hai xô sơn trên tay, khuôn mặt của họ lấm lem, thậm chí còn chẳng mặc một cái tạp dề.

Haato nghiêng đầu thắc mắc: ``Hai đứa làm cái gì vậy?'' Cặp đôi Chó-Mèo bình thản đáp, tay vẫn cần xô sơn, với tính cách có phần đối nghịch nhau:

\begin{dialogue}
    \speak{Okayu} *bình thản* Bọn em được phân công sơn lại cả nhà\~{}
    \speak{Korone} *hào hứng* Để tụi này giải quyết nha\~{} Lên nào, Ogayu!
\end{dialogue}

Không chờ thêm lời nào, nàng Khuyển cầm chổi cuộn sơn và nhanh chóng bắt tay vào việc. Cô lấy một ít sơn từ xô, nhúng chổi rồi quét lên bức tường với những đường nét mạnh mẽ và đầy năng lượng. Mỗi lần quét, sơn bắn tung tóe, nhưng Korone dường như không hề để ý, miệng còn cất tiếng hát vui tai. Trong khi đó, nàng Miêu dùng máy sấy để làm khô lớp sơn ngay tức thì. Máy sấy phát ra tiếng rít nhẹ, làm cho lớp sơn mới khô nhanh như thổi, bề mặt bóng mịn và đều màu.

Noel trố mắt, reo lên: ``Oa\dots\ Nhanh quá!'' Cô không giấu được sự ngưỡng mộ trước tốc độ làm việc của cặp đôi này.

Roboco cũng gật gù, mắt sáng lên đầy thích thú. Cô nàng bắt đầu tính toán trong đầu, có lẽ đang nghĩ đến việc tích hợp một máy sấy vào cánh tay mình.

Haato thì ôm trán, thở dài: ``Hai đứa này\dots\ làm việc nhanh thật, nhưng mà\dots'' Cô nhìn quanh, thấy vài vết sơn loang lổ trên sàn và trần nhà, lắc đầu: ``\dots làm xong rồi mà còn nhem hơn cả lúc chưa sơn!''

Chẳng mấy chốc, cả bức tường trắng tinh khôi như mới. Korone vẫy đuôi tươi cười, còn Okayu gật đầu hài lòng.

\begin{dialogue}
    \speak{Korone} Chúng em xong việc rồi\~{}
    \speak{Okayu} Tiếp nào, Koro-san! Xuống tầng dưới thôi\~{}
\end{dialogue}

Hai cô gái thoăn thoắt xách xô sơn rời đi, để lại ba người còn lại đứng trầm trồ. Ngay cả Noel và Roboco cũng phải ngạc nhiên trước tốc độ làm việc như chớp của cặp đôi này.

\begin{center}
    \uline{15 phút sau.}
\end{center}

Sau một hồi miệt mài hoàn tất việc quét dọn, lau sàn, rồi kê đồ đạc về vị trí cũ, và thậm chí sửa cả kệ đồ trong nhà vệ sinh, cuối cùng họ cũng hoàn thành công việc.

Haato ngồi bệt xuống dựa lưng vào tường, vừa thở sâu vừa lau mồ hôi trên trán, tỏ vẻ mệt mỏi: ``Cuối cùng cũng xong\dots\ Quản với hai đứa này mệt thật đấy! \textit{Giờ mình đã hiểu được cảm giác của Kuon-senpai thì anh ấy phải quản bọn mình rồi\dots}''

Noel thì vẫn còn sung sức, cô vẫn có thể làm việc tiếp. Cô chỉ đống dây máy móc mà họ tháo ra để lau dọn vệ sinh, hỏi: ``Nhưng với mớ dây này thì làm sao đây?''

Nàng Người máy hào hứng xung phong: ``Để chị đấu lại cho\~{}'' nhưng bị hậu bối cô cản lần nữa: ``Roboco-senpai, em nghĩ chị đừng nên động vào thì hơn. Vụng tay như chị có khi làm hỏng đồ nhà anh ấy đấy.''

Roboco xị mặt, làu bàu ``Ể\~{}?'' tiếc nuối, vừa hơi tức khi cô ấy ám chỉ rằng mình là ``PON''. Nàng Hiệp sĩ liền nhan nhảu ngắt lời: ``Thế bây giờ chúng ta làm gì tiếp theo nhỉ?''

Haato từ từ đứng dậy, chỉnh lại quần áo, rồi hắng giọng: ``Kuon-senpai bảo chúng ta lau cầu thang đấy. Chuẩn bị nào mọi người.''

Hai cô nàng ``Ngố'' liền đồng thanh, đưa tay lên chào như quân nhân: ``Yes sir! (Rõ, thưa chỉ huy!)''

Haato toát mồ hôi, xua tay: ``Mấy người không cần phải thưa lệnh như quân đội đâu\dots''

Cả ba người họ rời khỏi căn phòng của cậu Tổng, vừa đi vừa trò chuyện, coi như cho mình một khoảng giải lao ngắn.

Nàng Người máy vừa đi vừa cười khúc khích: ``Nhưng mà chị thấy thú vị mà\~{} Thưa chỉ huy Haato\~{}''

Nàng Song tính lườm nhẹ và thở dài: ``Cả chị nữa, đừng có trêu em.'' Nàng Hiệp sĩ cười phì: ``Thôi nào, Haachama-chan. Vui vẻ một chút đi. Chúng ta còn nhiều việc phải làm lắm.''

Ba người tiếp tục bước xuống cầu thang, vừa đi vừa trò chuyện rôm rả. Tiếng cười của họ vang vọng khắp ngôi nhà, xua tan bầu không khí mệt mỏi của công việc.

Ai cũng biết, dù vất vả thế nào, chỉ cần có bạn đồng hành, mọi chuyện đều trở nên dễ dàng hơn.

\vspace{-40pt}
\subsubsection*{\phantomsection\label{ep-003.IVb}}
\addcontentsline{toc}{subsubsection}{Shion và Subaru}

\begin{center}
    {\color{T2} Nhóm 2 - Shion và Subaru}\\Phòng của Hijaki.
\end{center}

``X-Xin phép ạ\dots''

Nàng Vịt mở cửa kéo của căn phòng cùng tầng với phòng của Kuon ra. Đó là một căn phòng về cơ bản giống như phòng bên cạnh, nhưng bên trong có hai bàn học với những kệ tủ đầy sách, một cái đệm nằm sát với góc phòng, và hai tủ quần áo.

Vừa bước vào, Subaru đảo mắt một lượt và thở dài: ``Trông như bãi chiến trường vậy. Xem chừng khoai rồi đấy, Shion-chan. Hay là chúng ta đợi\dots''

Nhưng nàng Phù thủy nhanh chóng cắt lời: ``Khỏi cần, cứ để đó cho tôi.''

Nàng ``Tỏi'' lẩm nhẩm vài câu thần chú, một vòng tròn phép màu xanh lá hiện ra trước mặt cô. hững ký tự kỳ lạ phát sáng lơ lửng, rồi cô giơ tay lên cao và hô to: ``\textbf{Tạo phép!}''

Một tia sáng chói lòa bao phủ cả căn phòng, khiến Subaru phải che mắt lại. Khi ánh sáng dịu dần, nàng Tomboy mở mắt ra và\dots

\dots trước mặt họ chỉ còn là một căn phòng rỗng không, ``bay sạch'' gần như mọi thứ, kể cả vết bẩn, bụi bặm hay rác rưởi. Nhìn sơ qua, người ta sẽ tưởng đây là một căn phòng mới được xây, chỉ trừ mỗi bức tường vẫn chưa được sơn lại.

Nàng Vịt há hốc mồm trước khung cảnh đó: ``S-Shion!? Bà vừa làm cái quái gì thế!?''

Shion bình thản đáp: ``Thì tôi vừa mối dọn phòng rồi còn gì. Đúng như lời Kuon-senpai nói đấy.''

``\textbf{Dọn phòng không phải như thế!!}'' Subaru gắt lên, tỏ vẻ tức giận trước một căn phòng trống trơn.

Shion nhìn căn phòng, rồi ngơ ngác nghiêng đầu nhìn nàng Tomboy đang tức giận: ``Ể? Tôi tưởng dọn phòng là dọn bụi với đổ rác chứ?''

``Nhưng không có nghĩa là làm \emph{mọi thứ} biến mất sạch sẽ như vậy!'' nàng Vịt nhấn mạnh. Nhưng trước khi cô kịp cằn nhằn gì thêm, từ cầu thang lại vang lên tiếng hô lớn:

\begin{dialogue}
    \speak{Korone} Đội Sơn tường, xuất kích! Gâu, gâu!
    \speak{Okayu} Xin thứ lỗi mọi người nha, meo\~{}
\end{dialogue}

Subaru thảng thốt: ``Đây có phải là cuộc đua đâu! Mà này, hai người làm cái trò gì vậy!?''

Cặp đôi Khuyển-Miêu lao vào phòng với tốc độ nhanh như gió, tay xách theo dụng cụ sơn và máy sấy. Không để ai kịp phản ứng, Korone bắt đầu quét sơn lên tường, vừa làm vừa hát nghêu ngao.

Cô nàng Chó sơn nhanh đến mức còn\dots\ vẽ cả hoa văn ngẫu hứng lên cửa kính, thậm chí còn bay lên hẳn trần nhà để sơn. Okayu lặng lẽ theo sau, cầm máy sấy làm khô từng lớp sơn. Chỉ trong vài phút, bức tường đã được phủ một lớp màu hồng tươi sáng như mới.

\begin{dialogue}
    \speak{Korone} *vẫy đuôi phấn khởi* Xong rồi nha\~{} Sang phòng khác thôi, Ogayu!
    \speak{Okayu} *gật đầu* Rồi\~{}
\end{dialogue}

Hai nàng lập tức rời đi, để lại Subaru và Shion đứng ngỡ ngàng với tốc độ làm việc thần tốc.

``Giờ bức tường cũng được sơn xong rồi, ta có thể xuống nhà được chưa?'' nàng Phù thủy hỏi với sự bình thản.

Subaru, dù vẫn còn sốc với khung cảnh Korone và Okayu làm với tốc độ siêu thanh, liền lớn giọng đáp: ``Ít nhất bà phải trả lại nội thất cho người ta chứ!''

Shion thở dài: ``Đúng là phiền hà thật.'' Cô búng tay một cái, lập tức toàn bộ đồ đạc xuất hiện trở lại, đúng nguyên vị trí cũ, nhưng vẫn giữ sạch sẽ đến mức bóng loáng, không một hạt bụi hay rác thải gì. ``Được rồi, thế này được chưa? Vẫn không một hạt bụi luôn nhé.''

Subaru ôm mặt ngán ngẩm, tới mức tưởng như sắp khóc thét đến nơi: ``Đáng ra bà phải làm thế từ đầu chứ!'' Đầu cô như thể đang xì khói vì vẫn đang nén cơn tức giận.

Dù thế, nàng Phù thủy mặc kệ nàng Vịt đang ôm đầu đó, mở cánh cửa trượt ra lần nữa rồi bảo: ``Thôi, đừng inh ỏi nữa. Xuống chuẩn bị lau cầu thang nào.''

Nàng Tomboy vội ngẩng đầu lên, nhìn cô nàng với vẻ bối rối: ``Ể? Cầu thang à?'' Nàng ``Tỏi'' nhướng mày: ``Não bà để đi đâu vậy? Kuon-senpai bảo làm xong việc thì cùng mọi người xuống lau cầu thang mà.''

Không để Subaru phản ứng, nàng Phù thủy hất tóc và bước ra khỏi phòng. Cô đứng ngẩn ra một lúc rồi vội vàng chạy theo, reo lớn: ``Ơ, từ từ đã, đợi tôi với! Mồ, cái bà này, lúc thì lười chảy thây, lúc thì nhanh như chớp\dots\ đúng là không hiểu nổi mà!''

Nghe lời bình của người đồng Thế hệ, Shion ngoái lại, cười khẩy: ``Bà còn định lắm mồm tới bao giờ nữa! Nhanh lên nào, không là chậm bây giờ!''

Nàng Vịt vừa chạy vừa tiếp tục cằn nhằn: ``Tôi đã bảo rồi, làm cái gì cũng phải từ từ và cẩn thận! Bà cứ như vậy thì có ngày gây hoạ lớn cho mà xem!''

Shion nhún vai, lè lưỡi như muốn trêu tức cô: ``Miễn là xong việc là được thôi, ai quan tâm chứ\~{}''

Hai người vừa đi xuống cầu thang vừa tiếp tục đôi co, làm náo động cả tầng dưới. Tiếng bước chân rầm rập và tiếng tranh cãi rộn ràng của họ khiến những đội khác cũng phải bật cười.

Dẫu có trái tính trái nết, họ vẫn phối hợp ăn ý theo cách rất riêng của mình.

\vspace{-40pt}
\subsubsection*{\phantomsection\label{ep-003.IVc}}
\addcontentsline{toc}{subsubsection}{Matsuri, Sora, Kuon, Ryujin, A-chan}

\begin{center}
    {\color{T3} Nhóm 3 - Matsuri, Sora, Kuon, Ryujin, A-chan}\\Siêu thị GO! - Kyuukami.
\end{center}

Ngày 29 Tết, siêu thị náo nhiệt đến mức chẳng khác gì một lễ hội. Hàng trăm người chen chúc nhau giữa những kệ hàng, những chiếc xe đẩy chất đầy đồ đạc lách qua lách lại, và âm thanh từ loa thông báo liên tục vang lên làm không khí thêm phần sôi động. Nhóm mua sắm của \hll\ cũng hoà mình vào dòng người đông đúc, với đội hình gồm nàng Thần tượng Sora, nàng Cổ động Matsuri, nàng Thỏ Pekora, dưới sự giám sát của A-chan, cùng sự hướng dẫn của hai bố con nhà Hikawa - Kuon và ông Ryujin.

Nàng Cổ động nhìn biển người hối hả mua đồ, đổ mồ hôi hột: ``Thế khó rồi đấy\dots'' Pekora thì tỏ ra bất ngờ: ``Em chưa bao giờ thấy nhiều người cùng lúc đi mua sắm như thế này, peko!''

Kuon thì bày tỏ sự bình thản, như thể đã quá quen với khung cảnh này: ``Cũng chẳng trách được. Mấy hôm giáp Tết người ta đổ xô đi mua đồ mà.'' Sora bày tỏ lo lắng: ``Thế này có ổn không nhỉ? Em sợ là có khi hết hàng mất\dots''

Matsuri gật đầu: ``Cũng phải\dots\ Mọi người có cao kiến gì không?''

Ông Ryujin, với phong thái điềm tĩnh thường thấy, đứng suy nghĩ một hồi rồi đưa ra ý kiến: ``Chú nghĩ chúng ta nên chia nhóm ra. Mỗi nhóm vào một khu riêng để mua đồ theo danh sách. Làm như vậy sẽ tiết kiệm thời gian và tăng cơ hội tìm được hàng hơn.''

``Con thấy cách này có vẻ khả dĩ hơn đấy,'' Kuon tỏ vẻ tán thành với ý kiến của ông ấy, rồi nói tiếp: ``Có khi nhiều người mua cùng lúc mới có cơ may tìm được hàng.''

Những người còn lại nhìn nhau gật đầu đồng tình. Họ nhanh chóng chia thành hai nhóm: Kuon, Pekora và Ryujin nhận nhiệm vụ mua đồ tươi như thịt, cá, và hải sản; trong khi nhóm của Sora, Matsuri và A-chan phụ trách mua đồ khô như bánh kẹo, mứt Tết và gia vị.

\ct

Là những người bản địa, Kuon và Ryujin nắm rất rõ sơ đồ siêu thị. Họ dẫn nàng Thỏ Pekora đi qua từng gian hàng một cách bài bản. Dù dòng người đông đúc đôi lúc khiến việc di chuyển trở nên khó khăn, nhưng với sự hướng dẫn nhanh nhẹn của hai bố con nhà Hikawa, họ vẫn kịp thu gom đủ đồ tươi cần thiết: một tảng thịt bò, một tảng thịt lợn, vài cánh gà, và một số hải sản gồm mực tươi cùng tôm.

Tuy nhiên, với Pekora, đây lại là trải nghiệm đầu tiên của cô nàng khi đi mua đồ ở một nơi đông nghịt đến thế. Ban đầu, cô còn khá bỡ ngỡ khi nhìn cảnh tượng hàng người đang nườm nượp đi mua sắm, nhưng rồi khi được cậu Tổng nói\dots

``Cảm giác đi mua hàng cứ như thể là đi đánh trận vậy. Chẳng biết thế này chúng ta có tranh được đồ không nữa.''

\dots bỗng chốc, nàng Thỏ lại tỏ vẻ phấn khích hơn hẳn. Đôi tai thỏ của cô dựng đứng, ánh mắt sáng lên đầy quyết tâm.

Cả người cô như thể bùng lên một ngọn lửa, trông cô như đang chuẩn bị đi đánh trận thật. Cô cười lớn: ``Đi đánh trận thì phải để em, peko! Để xem ai mới là người giành được chiến thắng hôm nay!''

Nhìn cô bùng cháy lên, cả hai bố con nhà Hikawa liền đổ mồ hôi hột. Ryujin cười khì: ``Nhân viên của con đúng là thú vị thật,'' còn bản thân Kuon thì thở dài: ``Làm ơn đừng có làm loạn lên đó, Pekora-chan.''

\il

Công cuộc đi mua sắm của ba người họ vẫn diễn ra rất bình thường, cho đến khi nàng Thỏ nhắm đến một gian hàng gần đó, nơi đang treo biển ``Giảm giá 50\% - Chỉ còn 10 phần!'' (viết bằng tiếng Etsunan, nhưng bằng cách nào đó cô nàng có thể đoán được) với đám đông chen chúc chật kín. Gian hàng này bán các hộp trứng gà, vốn là món hàng không thể thiếu trong mỗi bữa ăn Tết.

Nàng Thỏ nhìn thấy vậy, liền chạy tới gian hàng kia, không kịp để Kuon ngăn cản. ``P-Pekora-chan! Từ từ đã\dots!''

Nhưng chưa kịp ngăn, nàng Thỏ đã lao thẳng về phía đám đông như một cơn lốc. Cô len lỏi qua từng người với tốc độ nhanh không tưởng, vừa chạy vừa hô lớn: ``\textbf{Trứng gà giảm giá là của ta, peko!!}''

Dòng người chợt sững lại trong giây lát vì bất ngờ, nhưng nhanh chóng trở nên náo nhiệt hơn, bất chấp việc cô nàng vừa hét lên bằng thứ ngôn ngữ mà không ai hiểu. Ai nấy đều lao vào cuộc chiến tranh giành những hộp trứng cuối cùng.

``{\color{red} Cái quái gì đang xảy ra vậy trời\dots}'' Kuon nhìn khung cảnh tranh giành những món hàng cuối cùng, không biết nên phản ứng như thế nào.

Ở giữa khung cảnh hỗn loạn, Pekora bật nhảy cao lên, vươn tay chộp lấy ba hộp trứng ở tầng trên cùng của kệ hàng, rồi tung tất cả lên không trung. ``\textbf{Kuon-senpai, Ryujin-san! Bắt lấy này, peko!!}'' cô hô vang, hướng tới hai bố con.

Cô tung hai hộp trứng lên cao với tốc độ rất nhanh. Kuon thảng thốt hét lên: ``Thế này thì nhanh quá!!'' Trái ngược với người con đang hoảng hốt, người bố nhảy dựng lên, vững vàng đón được hai hộp trứng.

``B-Bố ơi\dots\ T-Thế nào mà lại\dots'' cậu nhìn ông ấy đón được hai hộp trứng mà vẫn còn nguyên vẹn với một cặp mắt ngỡ ngàng.

Ông bố từ tốn đáp: ``Bố từng làm thủ môn cho đội bóng của lớp hồi còn đi học mà.'' Kuon nghe vậy, đổ mồ hôi hột: ``Bảo sao\dots''

``\textbf{Kuon-senpai! Còn tiếp này!}'' nàng Thỏ hét lên làm cậu giật mình. Ngay khi cậu ngoảnh mặt về phía nàng Thỏ, hộp trứng còn lại liền được tung lên làm cho cậu không phản ứng kịp.

``\textbf{Em đùa anh à!!!}''

Nhưng may thay, cậu kịp chạy về phía sau, xoay người, đưa tay chụp lấy hộp còn lại như một vận động viên bóng chày chuyên nghiệp.

Cậu nhìn hộp trứng vừa bắt được vẫn còn nguyên vẹn, liền thở dài nhẹ nhõm. ``Làm ơn đừng có làm mấy trò hú tim như vậy nữa được không\dots''

Mọi người xung quanh đứng sững lại, mắt chữ O, mồm chữ A, không tin nổi cảnh tượng vừa diễn ra trước mắt. Có người thậm chí còn vỗ tay cổ vũ, không rõ là ngưỡng mộ kỹ năng của họ hay chỉ đơn giản vì thấy màn trình diễn quá thú vị.

\begin{dialogue}
    \speak{Cus1} {\color{red} Cô gái đó\dots\ là vận động viên Olympic à?}
    \speak{Cus1} {\color{red} Còn hai người kia thì chắc là vận động viên bắt bóng! Đỉnh ghê!}
\end{dialogue}

Pekora, vẫn đang đứng trên kệ hàng, cười lớn và giơ ngón cái về phía hai bố con nhà Hikawa: ``Chiến thắng thuộc về chúng ta, peko!''

Kuon lồm cồm bò dậy, tay đưa hộp trứng vào xe đẩy hàng, thở dài ngán ngẩm: ``Đây là siêu thị chứ có phải là đấu trường đâu\dots\ Làm ơn đừng có làm anh phát điên ở đây chứ\dots''

Trong khi đó, người bố thì bật cười: ``Kể ra thì cũng thú vị đấy chứ.'' Kuon nghe người bố của mình nói vậy, cũng chẳng biết làm gì hơn, không đưa ra thêm một ý kiến nào nữa.

Dù có chút ồn ào, cuối cùng họ cũng rời khỏi gian hàng trong sự ngưỡng mộ và lời bàn tán của những người xung quanh. Kuon nhanh chóng kéo Pekora lại gần và nhắc nhở: ``Lần sau đừng làm vậy nữa, không là anh chẳng biết phải xin lỗi người ta thế nào đâu!''

Nàng Thỏ xanh vẫn tỏ ra đắc ý: ``Nhưng em đã mang lại trứng gà cho mọi người rồi mà, peko!'' Kuon ôm đầu bất lực: ``Nhưng mà\dots\ *thở dài* Đúng là hết cách.''

Cả ba tiếp tục hành trình mua sắm, dù Kuon vẫn giữ ánh mắt cảnh giác, chỉ sợ nàng Thỏ lại gây thêm rắc rối gì nữa. Nhưng trong lòng cậu, cậu cũng phải thừa nhận rằng sự nhiệt tình của Pekora khiến không khí mua sắm hôm nay trở nên vui nhộn hơn hẳn.

\ct

Việc mua hàng cũng đã xong xuôi, với một chiếc xe đẩy ú ụ đồ, nhiều đồ hơn hẳn so với mọi năm nhà Hikawa hay tổ chức.

Pekora hứng khởi nói: ``Làm việc với Kuon-senpai đúng là dễ dàng quá, peko! Mọi thứ đều ở đúng chỗ em cần, peko!''

Kuon mỉm cười đáp: ``Thì, bọn anh sống ở đây mà. Chỉ cần biết chỗ thì mọi thứ đơn giản thôi.'' Nhưng rồi nụ cười của cậu tắt hẳn: ``Cơ mà đây chưa phải là cái khó nhất đâu.''

``Ể?''

``Đây này. Nhìn đằng trước đi.''

Quả đúng như cậu nói, thử thách lớn nhất của họ vẫn là\dots\ xếp hàng thanh toán. Từng quầy thu ngân đều chật kín hàng người, khiến cho việc tìm một quầy để thanh toán nhanh gọn trở nên cực kỳ khó khăn.

\begin{dialogue}
    \speak{Ryujin} Đông thật đấy\dots
    \speak{Kuon} Cũng phải thôi, hôm nay 29 Tết mà.
    \speak{Pekora} Đi đâu cũng toàn hàng người là người vậy, peko!
\end{dialogue}

Cứ tưởng rằng họ sẽ phải xếp hàng dài đằng đẵng tại một trong những quầy thu ngân đông đúc, thì bất ngờ nàng Thỏ tinh mắt phát hiện ra một quầy gần như trống trơn, chỉ có ba đến bốn khách hàng đang đứng chờ.

``Quầy kia đang vắng người kìa, peko!'' nàng Thỏ la lên, chỉ tay về phía quầy thu ngân đó, gương mặt rạng rỡ như vừa tìm được kho báu.

``Nó cũng khá gần hơn với cửa ra đấy\dots'' cậu Tổng nhận định, khi quầy thu ngân đó chỉ cách họ có đúng một quầy hàng.

Không ai bảo ai, cả ba đẩy xe hàng tiến về quầy thu ngân với tốc độ nhanh nhất có thể trong điều kiện đông đúc. Khi đến nơi, họ thở phào nhẹ nhõm, thấy hàng chờ trước mặt không quá dài, và tự nhủ rằng việc thanh toán sẽ diễn ra suôn sẻ.

Thế nhưng, khi lượt của họ gần tới, một người đàn ông cao lớn với dáng vẻ bặm trợn bất ngờ chen ngang. Gã ta chẳng thèm để ý đến hàng người đang xếp trật tự phía sau, cứ thế đẩy xe của mình lên trước mặt ba người. Hắn thét lên một giọng hống hách như khạc ra lửa, nói bằng tiếng bản xứ: ``{\color{red} Tránh ra! Tao bận, không có thời gian xếp hàng đâu!}''

Trước sự vô ý thực của hắn ta, Pekora lên tiếng bất bình: ``{\color{red} Này anh kia! Anh đang làm cái trò gì thế hả, peko!?}''

Nếu bạn đang định hỏi bằng cách nào mà Pekora có thể nói tiếng bản xứ ở đây, thì Kuon đã đưa cho các cô gái \hll\ bánh mì chuyển ngữ do chính cậu làm, nên việc giao tiếp với con người tại đây không phải là vấn đề quá lớn.

Gã đàn ông quay lại, nheo mắt nhìn nàng Thỏ, giọng khàn đặc đầy mỉa mai: ``{\color{red} Mày biết bố mày là ai không mà dám lên tiếng với tao hả, con thỏ nhãi kia?}''

Những lời lẽ thô lỗ của gã khiến mọi người xung quanh nhíu mày khó chịu, nhưng không ai dám can thiệp. Ông Ryujin, vốn là người điềm đạm, định bước lên khuyên nhủ thì nàng Thỏ đã nhanh hơn. ``Việc này cứ để cháu lo, peko.''

``Nhưng mà, Pekora-chan\dots'' cậu định lên tiếng tỏ vẻ lo lắng cho cô, thì cô ấy đã nháy mắt với cậu, nói với sự tự tin: ``Rồi em sẽ cho mọi người thấy thế nào là sức mạnh của beta carotene\footnote{Một loại tiền chất của vitamin A, có trong cà rốt.}!''

Chưa kịp để cậu mở lời thêm câu nào, Pekora đã sồn sồn tiến về phía trước và đụng lại vào thân gã.

``{\color{red} Lại là mày à, tên oắt con kia?}'' gã nói. Nàng Thỏ đáp lại, cười một cách lạnh lùng: ``{\color{red} Tôi không biết bố anh là ai, nhưng tôi biết luật là phải xếp hàng, peko. Anh muốn làm khác thì phải qua tôi trước đã!}''

``{\color{red} Mày cũng gớm đấy nhỉ? Một con thỏ đớ èo oặt như mày thì làm cái trò trống gì-}'' nhưng chưa để hắn ta nói hết câu, Pekora đã nhanh như chớp lao tới, khéo léo dùng sức bật và tốc độ của mình để áp sát.

Một cú xoay người chuẩn xác khiến gã mất thăng bằng, rồi cô quét chân, làm gã ngã dúi dụi xuống sàn, trước sự ngỡ ngàng cùng với sự cười đùa của mọi người xung quanh.

``{\color{red} M-Mày dám\dots! Thằng chó này!}'' hắn ta thét lên khi bị mọi người chế giễu hắn ta.

``{\color{red} Tôi là Thỏ, chứ không phải là Chó đâu, peko!}'' Pekora đáp lại, cười hắn ta một cách khinh bỉ.

Tiếng xì xào bắt đầu vang lên từ đám đông, nhiều người lấy điện thoại ra quay lại cảnh tượng hiếm thấy này.

Hắn ta đứng bò dậy, mắt nghiền lại. ``{\color{red} Bà tiên sư thỏ nhà mày, nhận lấy!}'' Hắn giơ nắm đấm lên, nhằm vào nàng Thỏ đang múa máy tay chân. ``{\color{red} \textbf{Chết đi!!}}''

Thế nhưng, ngay khi nắm đấm của hắn ta đến nơi, Pekora nhảy bật lên cao, khiến cho nắm đấm của gã trúng vào thanh sắt ngăn cách của quầy thu ngân.

Lập tức, gã đã bị nàng Thỏ khống chế hoàn toàn bằng một cú khóa tay gọn gàng, mặc cho hắn ta kêu lên như một con lợn bị cắt tiết. ``{\color{red} Đừng giãy nảy nữa, kẻo tay ông đau đấy, peko!}''

Chỉ ít phút sau sự lộn xộn đó, nhân viên siêu thị cùng một đội bảo vệ xuất hiện. ``{\color{red} Chuyện quái gì đang xảy ra ở đây vậy?}'' quản lý của siêu thị nghiêm giọng hỏi, cùng với đó là các bảo vệ nhanh chóng tiếp cận tên ``đầu trâu mặt ngựa'' đó.

Kuon nhanh chóng giải thích sự việc, chỉ rõ gã đàn ông đã chen ngang và gây rối. Sau khi kiểm tra camera an ninh và lấy lời khai từ những người xung quanh, bảo vệ quyết định gọi cảnh sát.

``{\color{red} Nào! Bây giờ mời anh cùng chúng tôi lên phường!}'' một trong những bảo vệ `áp giải' hắn ra khỏi siêu thị. Trong lúc bị họ đưa đi, miệng của hắn ta vẫn không ngừng mắng chửi, thậm chí còn đòi thề sẽ trả thù, nhưng rõ ràng hắn không thể làm gì thêm. Thậm chí, hắn còn bị một trong những bảo vệ khác đấm thẳng vào bụng khiến cho hắn ta câm nín.

Sau khi tên đồ tể kia đi rồi, nàng Thỏ phẩy tay: ``Xong việc rồi, peko! Chúng ta thanh toán thôi.''

Cậu con trai sau khi chứng kiến cảnh vừa rồi cũng chỉ có thể vỗ trán ngán ngẩm: ``Lần nào đi với các em cũng có chuyện\dots\ Nhưng mà, cảm ơn em, Peko-chan.''

Người bố thì cười một cách sảng khoái: ``Nhưng không thể phủ nhận là Pekora-san làm rất tốt. Dù sao thì mọi chuyện cũng êm đẹp rồi.''

Cả ba đẩy chiếc xe hàng tới thanh toán hàng hóa. Dù có chút náo động, nhưng rõ ràng họ đã trở thành tâm điểm chú ý và còn nhận được những ánh nhìn ngưỡng mộ từ đám đông.

Nữ nhân viên thu ngân nhìn ba người họ, nở một nụ cười thân thiện: ``{\color{red} Của mọi người hết 4.230.000 ạ.}''

Pekora bày tỏ ngạc nhiên: ``Tại sao giá ở đây lại cao thế, peko!?'' Kuon bật cười, giải thích: ``Bốn triệu ở đây là bốn triệu đồng, không phải bốn triệu yên đâu. Đổi ra có hơn hai mươi nghìn yên thôi à.''

Ngay khi được cậu giải thích, nàng Thỏ lập tức chuyển đổi thái độ sang vui vẻ hơn hẳn: ``Tưởng thế nào chứ! Giá ở đây rẻ thật á!''

Cả ba người bật cười với câu bình của nàng Thỏ đó. Sau khi họ trả tiền cho cô thu ngân, cô ấy nói: ``{\color{red} Quý khách đã trả đủ tiền rồi ạ. Thật sự cảm ơn các vị rất nhiều! Chúc mừng năm mới các vị nhé!}''

Pekora cười lớn: ``Cảm ơn cô, peko!'', Ryujin nhìn đồng hồ trên điện thoại của mình, rồi chỉ về khu vực bên ngoài: ``Xong rồi, giờ ra ngoài chờ nhóm kia thôi.''

Cả ba người đẩy chiếc xe hàng ngồi nghỉ trên băng ghế bên ngoài siêu thị, tận hưởng chút không khí mát lạnh của buổi chiều cuối năm. Hai bố con ngồi lại nhìn khung cảnh hối hả của chiều cuối năm, trong khi Pekora chạy ra đi vệ sinh.

Nhưng với Kuon, tâm trí cậu vẫn còn để ở nhóm còn lại. Cậu khẽ thở dài, quay sang người bố: ``{\color{red} Con có cảm giác\dots\ không ổn với nhóm \hll\ kia.}''

Ryujin bật cười: ``{\color{red} Cứ chờ xem. Nếu có chuyện gì, chắc chắn họ sẽ tự xoay sở được. Con cứ yên tâm đi.}''

``{\color{red} \dots Con cảm thấy nghi lắm.}''

\il

Đúng như linh cảm của Kuon, nhóm của Sora, Matsuri và A-chan lại không được may mắn như vậy. Ngay từ lúc bước vào khu vực đồ khô, sự chật chội cùng với những chiếc kệ gần như trống trơn đã khiến họ lúng túng. Và rồi, mọi chuyện nhanh chóng trở nên lộn xộn\dots

\begin{center}
    \uline{Quầy bánh mì.}
\end{center}

Trong danh sách mua hàng, họ cần mua hai ổ bánh mì dài để dùng trong các bữa ăn ngày Tết. Tuy nhiên, như một quy luật bất thành văn, quầy bánh mì luôn là một trong những  nơi tập trung đông đúc nhất trong siêu thị vào dịp cuối năm. Một dòng người chen chúc, ai cũng muốn nhanh tay lấy được ổ bánh nóng hổi vừa ra lò, khiến nhóm ba người chần chừ trước khi tiến lại gần.

Sora nhìn dòng người tấp nập đang xếp hàng lấy bánh mì, bày tỏ vừa bất ngờ vừa e sợ: ``Đông thật đấy!''

A-chan thở dài: ``Cứ thế này thì chẳng biết bao giờ mới tới lượt mình\dots\ Mà càng muộn thì người lại càng đông hơn.''

Thấy vậy, nàng Cổ động hào hứng giơ tay: ``Vậy để em xử lý cho!'' Trước khi hai người kia kịp nói gì, cô nàng đã lanh lẹ chen vào dòng người. Matsuri len lỏi một cách chuyên nghiệp, thậm chí còn nhanh chóng tìm được một chỗ đứng tốt ở gần quầy bánh mì.

``Liệu có ổn không nhỉ? Matsuri-chan mà vào đấy thì không biết có chuyện gì xảy ra không nữa\dots'' nàng Tăng ca nhìn cô ấy, bày tỏ sự lo lắng. Sora nhìn về cô nàng đó, rồi dịu dàng đáp: ``Em nghĩ Matsuri-chan sẽ xoay xở được thôi. Hay là trong lúc chờ, chúng ta đi mua đồ ở các quầy khác trước rồi quay lại?''

A-chan đành phải gật đầu miễn cưỡng: ``Chắc là vậy\dots\ Nhưng nhớ quay lại nhanh nhé, không Natsuiro-san lại làm loạn lên mất.''

Vậy là hai người quyết định đẩy xe hàng rời khỏi khu vực đông đúc, tiến về phía gian hàng đồ khô. Cả hai nghĩ rằng sẽ mất kha khá thời gian để Matsuri lấy được bánh mì, nên tranh thủ hoàn thành nốt các mục trong danh sách mua sắm.

Tuy nhiên, điều mà họ không ngờ tới là ngay khi họ vừa khuất bóng, quầy bánh mì đã bắt đầu bày ra những ổ bánh mì dài, vàng ruộm, thơm phức. Nhân viên vừa đặt bánh lên các khay gỗ, cả đám đông lập tức lao tới như một cơn sóng. Matsuri với bản năng nhanh nhẹn, không chỉ lấy được hai ổ bánh mì như kế hoạch ban đầu mà còn gom thêm ba ổ nữa.

Nàng Lễ hội hào hứng cầm năm ổ bánh mì nóng hổi trên tay, thầm nghĩ: ``\textit{Lấy thêm vài ổ mang về, chia cho mọi người ăn vặt, chắc sẽ không sao đâu nhỉ!}'' Cô hí hửng quay lại chỗ mà A-chan và Sora đáng lẽ đang đứng, nhưng\dots\ không thấy bóng dáng ai cả.

``À-À-rế!? Họ đi đâu mất rồi?'' cô hoảng hốt. Nhìn quanh mà không thấy, nàng tóc nâu bắt đầu chạy khắp khu vực quầy bánh mì, vừa cố tránh dòng người đông đúc, vừa giữ chặt túi bánh mì, vừa cất giọng gọi lớn: ``\textbf{A-chan! Sora-chan! Hai người ở đâu vậy!?}''

Tiếng gọi khẩn thiết của cô vang lên trong không gian ồn ào của siêu thị, thu hút ánh nhìn tò mò của những người xung quanh.

Một vài người lớn tuổi lắc đầu, khẽ cười trước cảnh tượng một cô gái trẻ, tay xách túi bánh mì, miệng gọi bạn như thể đang tìm kiếm một đứa trẻ lạc. Một số khách hàng thì tò mò nhìn cô như muốn hỏi: ``Người lớn rồi mà còn phải gọi bạn như thế à?''

Sẵn tiện cũng nói, ở Kawachi không mấy nhiều người biết về \hll, nên đa phần mọi người trong siêu thị cũng chỉ nhìn cô như là một người bình thường thay vì là một nàng hề- à không, idol; với Sora và A-chan cùng như vậy. Cùng lắm sẽ có những người thấy họ là các cosplayer\footnote{Những người hóa trang thành các nhân vật trong các tác phẩm (có thể là phim ảnh, văn học, hoặc trò chơi điện tử).}, chứ không coi họ như những người đặc biệt hay người nổi tiếng gì cả.

Matsuri, không hay biết mình đang trở thành tâm điểm chú ý, tiếp tục chạy lòng vòng giữa các quầy hàng, vừa chạy vừa la to: ``\textbf{Sora-senpai! A-chan! Em tìm hai người khắp nơi mà không thấy, ở đâu thế!?}''

Dẫu vậy, tiếng huyên náo của khách hàng xung quanh quá lớn, tiếng nói của cô chỉ như ``vọng vào núi rừng'', chẳng ai nghe thấy hồi đáp\footnote{Nguyên văn trong phiên bản Word: ``Dẫu thế, tiếng nói chuyện của khách hàng quá lớn khiến cho \textit{tiếng tôi vang rừng núi, sao không ai trả lời}.'' Đây là hai câu hát từ bài \textit{Chiếc khăn Piêu}, sáng tác bởi Doãn Nho.}.

\begin{center}
    \uline{Khu đồ khô.}
\end{center}

Sora cầm chiếc điện thoại ghi danh sách hàng cần mua, lẩm nhẩm: ``Được rồi, tiếp theo trong danh sách là\dots''

A-chan liếc qua danh sách, rồi cũng nhẩm theo: ``Mít sấy khô, kẹo dẻo, socola, mứt\dots\ Có vẻ như nhà Hikawa-san ăn nhiều đồ ngọt nhỉ. Ta nên mua thế nào đây?''

Nàng Thần tượng nhìn vào danh sách một lần nữa, rồi cười trừ với người bạn của mình: ``Anh ấy không định rõ số lượng hoặc loại gì ở trong này. Hay là chúng ta cứ lấy mỗi thứ một ít đi?''

Nàng Tăng ca gật đầu đồng tình: ``Tốt nhất là thế. Dù gì, không phải chỉ có nhà anh ấy ăn Tết mà.''

Tuy nhiên, khi Sora nói ``mỗi thứ một ít'', ý cô nàng lại khác hẳn. Với mỗi loại bánh kẹo, thậm chí là từng hương vị, nàng Thần tượng đều chọn một gói. Dẫu vậy, cô vẫn khéo léo chọn những món vừa túi tiền và được bạn đồng hành góp ý thêm, tránh trường hợp mua quá nhiều hoặc không hợp khẩu vị.

Chẳng mấy chốc, xe đẩy của họ đã đầy ú ụ đồ ngọt, từ kẹo, socola cho đến các loại mít sấy. Với mứt Tết, họ chọn đủ loại phổ biến như mứt dừa, mứt gừng, mứt hạt sen, mỗi loại hai hoặc ba hộp, càng khiến cho xe đẩy của họ đầy hơn, với hàng hóa chất cao như núi, khiến cho những người mua hàng không khỏi ngoái nhìn với số lượng lớn như vậy.

Nàng Đeo kính nhìn đống hàng đang lủng lẳng trên xe, thở dài: ``Chúng ta đã đi được một nửa danh sách, nhưng nhìn xe đẩy này thì sắp hết chỗ rồi.''

Nàng Thần tượng cười gượng, đồng tình với cô bạn: ``Được rồi, để em đi lấy thêm một chiếc xe nữa vậy.''

Nói xong, cô chạy nhanh về phía khu vực gần quầy thu ngân, kéo thêm một chiếc xe đẩy mới. Quay lại, cô nàng nhanh chóng tiếp tục cùng bạn mình sắp bớt đồ sang xe thứ hai, rồi tiếp tục kiểm tra danh sách: ``Tiếp theo sẽ là\dots''

\begin{center}
    \uline{Khu đồ uống.}
\end{center}

``\dots nước giải khát.''

Hai người đứng trước gian hàng đồ uống. Sora cầm danh sách, mắt nhìn lướt qua dòng ghi chú: ``Đồ uống: cái gì cũng được. Trà bí đao hoặc nước quả là tốt nhất.''

Nàng Thần tượng bày tỏ ngạc nhiên: ``Kuon-senpai thích uống trà bí đao sao?'' A-chan cười nhẹ: ``Có thể. Bình thường nếu không uống cà phê hay trà thì cũng uống nước ngọt có ga mà.''

Sora nở một nụ cười, dường như đọc vị được ý định của cậu Tổng: ``Hay là anh ấy đang muốn giảm cân?''

Nàng Đeo kính khẽ lắc đầu: ``Chuyện đó thì tôi không chắc. Nhưng thôi, cứ tìm thử xem, nếu có thì mua, còn không thì lấy đại loại nào khác cũng được.''

Cả hai cô gái bắt đầu đi quanh qua khu đồ uống. Sau một hồi tìm kiếm, họ phát hiện thùng trà bí đao nằm ở vị trí cao trên kệ. Cả hai ngước nhìn, rồi đồng loạt thở dài:

\begin{dialogue}
    \speak{Sora} Đằng kia kìa! Nhưng mà\dots
    \speak{A-chan} Để ở vị trí cao quá nhỉ?
\end{dialogue}

Họ tính tới chuyện trèo lên trên kia để lấy, nhưng có vẻ như cả hai người họ đều không thích ý tưởng này. Thêm nữa, hỏi một vị khách bất kỳ lấy đồ hàng cho họ cũng không khả thi, khi người người nhà nhà đều đang bận rộn với công việc mua sắm của họ. Hoặc ít nhất, là họ nghĩ vậy.

``Để em gọi nhân viên xem.'' Cô tiến tới một nhân viên siêu thị gần đó, nhẹ nhàng nhờ vả lấy hộ thùng bí đao ở kệ trên cùng kia.

Có vẻ như anh chàng nhân viên này bị thu hút bởi vẻ ngoài đáng yêu của nàng Thần tượng, vì anh lập tức trèo lên lấy thùng trà bí đao với sự nhanh nhẹn hiếm có.

Sau khi được anh nhân viên để hai thùng bí đao vào trong xe đẩy hàng, Sora lịch sự cúi đầu: ``Cảm ơn anh nhiều nhé!''

Anh nhân viên ngượng ngùng, rồi cắp chân tiếp tục công việc với vẻ mặt đầy hứng khởi. A-chan cười với sự đáng yêu của anh ấy: ``Nhìn anh ta có vẻ vui nhỉ?''

Bản thân nàng Thần tượng khẽ đỏ mặt: ``Chắc tại đây là lần đầu họ gặp người mua hàng như chúng ta. Mà, đừng trêu em nữa!''

Sau khi đặt thùng trà bí đao vào xe, Sora kiểm tra danh sách thêm lần nữa. ``Chúng ta chỉ còn mỗi dầu ăn, nước tương, sốt trứng mayonnaise với hạt nêm thôi.''

``Chắc mấy thứ đó nằm trong khu gia vị nhỉ?''

``Có thể lắm.'' Sora gật đầu. ``Đi thôi.'' Hai người đẩy xe tiến về khu gia vị, không quên trò chuyện rôm rả.

Trong lúc tiến tới khu tiếp theo, hai người họ không giấu nổi vẻ tò mò.

\begin{dialogue}
    \speak{A-chan} Không biết giờ mọi người đang làm gì nhỉ?
    \speak{Sora} Chắc là đang hối hả không kém chúng ta đâu.
\end{dialogue}

Họ tiếp tục cuộc hành trình mua sắm, còn trong lòng thì thầm nghĩ liệu nhóm Matsuri đã xoay sở xong chưa.

\il

\begin{center}
    \uline{Sảnh chờ.}
\end{center}

Nàng thỏ với mái tóc bạc, hai tai vểnh lên đầy vẻ bồn chồn, ngồi không yên trên ghế dài. Cô lẩm bẩm với chất giọng quen thuộc: ``Mọi người đi mua lâu thật đó, peko\dots''

Kuon ngước mắt lên từ chiếc điện thoại của mình, rồi cất giọng an ủi, dù không dấu nổi sự sốt ruột: ``Cứ kiên nhẫn thêm chút nữa đi. Siêu thị lớn thế này mà, mua đồ cho hơn hai mươi người thì đâu có nhanh được.''

Cậu bất giác nhìn sang xe đẩy chất đầy các túi rau củ mà Pekora mua từ trước, chép miệng: ``Mà anh để ý em mua nhiều rau củ lắm nhá. Nhất là\dots\ cà rốt.''

Nghe vậy, Pekora vùng vằng, mặt đỏ bừng: ``Thì sao chứ? Anh biết em yêu thích món tủ này như thế nào rồi mà sao còn nói, peko!''

Thấy phản ứng có phần hơi mãnh liệt của cô, cậu cười nhẹ, tiếp tục trêu: ``Ừ, anh chỉ thắc mắc thôi. Chứ nhìn cách em lấy cà rốt, chắc đến nửa năm không cần mua thêm đâu nhỉ?''

``M-Mồ! Đừng có mà chọc em nữa, peko,'' nàng Thỏ phồng má lên.

Cậu con trai chỉ nhún vai cười xòa, rồi lại đưa mắt nhìn xe đẩy trước mặt cả ba người họ. Sau một lúc ngắm nghía, cậu khẽ lẩm bẩm, giọng có chút ngạc nhiên: ``{\color{red} \dots Nhà mình mua thế này là cực nhiều đấy, so với mấy năm trước.}''

Bố cậu, Ryujin, đang ngồi ung dung uống ly trà từ bình giữ nhiệt, mỉm cười đáp: ``{\color{red} Nhà mình có tận hơn hai mươi người ăn Tết với nhau mà, nên bố nghĩ thế là vừa đủ đấy.}''

Cậu gật gù: ``{\color{red} Cũng có lý.} Mà công nhận, năm nay nhà mình ăn đông người thật.''

Pekora vừa cúi đầu nghịch điện thoại, vừa xen vào với giọng điệu như than vãn: ``Mất tận hơn hai mươi nghìn yên chứ chả ít gì, peko\dots''

Câu nói của cô khiến cả Kuon và Ryujin bật cười. Cậu nhìn tờ hóa đơn, rồi khẽ thở dài: ``Cũng may là bố con anh góp tiền mua, nên mới đủ. Chứ cứ như mọi năm thì thôi, viêm màng túi luôn.''

Pekora ngẩng đầu lên, đôi mắt cam nhìn cậu  ngơ ngác: ``Viêm màng túi? Ý anh là sao, peko?'' Cậu cười, giải thích với giọng dí dỏm: ``Là không đủ tiền mua ấy. Cách nói đùa ở xứ này là thế mà.''

Pekora khẽ ``à'' một tiếng, rồi lẩm bẩm ``mấy từ khó hiểu của người Etsunan'' trong khi tiếp tục nghịch điện thoại.

Không khí trong sảnh chờ rơi vào trạng thái yên lặng tương đối. Mỗi người cầm trên tay một chiếc điện thoại, làm việc riêng. Ryujin thỉnh thoảng kiểm tra tin nhắn từ người thân trong gia đình, Kuon thì tranh thủ lướt tin nhắn và đọc báo, còn Pekora lại mải mê với một trò chơi, đôi lúc bật cười khúc khích mà chẳng ai hiểu nổi lý do.

Thời gian dường như trôi qua chậm rãi hơn, và Kuon, trong lúc mắt vẫn dán vào màn hình, không quên liếc nhìn đồng hồ đeo tay, tự nhủ: ``Không biết nhóm kia còn lâu thế nào nữa nhỉ?''

\il

\begin{center}
    \uline{Một góc của siêu thị.}
\end{center}

``\textbf{A-chan! Sora-chan! Hai người đâu rồi!?}''

Tiếng gọi thất thanh của cô nàng Thế hệ thứ Hai vang lên giữa không gian đông đúc của siêu thị. Matsuri, với vẻ mặt gần như muốn khóc, đi lòng vòng khắp nơi để tìm hai người đồng đội. Trong mắt những người mua sắm gần đó, cô trông chẳng khác nào một đứa trẻ bị lạc mẹ, đặc biệt khi giọng nói của cô bắt đầu run rẩy và những giọt nước mắt lăn dài trên má.

Tệ hơn nữa, cô đang phải đối mặt với một ``tình huống khẩn cấp'': nhu cầu vệ sinh đã đạt đến mức báo động đỏ. Cô thầm trách mình vì đã không đi vệ sinh trước khi rời nhà và giờ đây, giữa biển người đông đúc, tìm được nhà vệ sinh còn khó hơn cả việc tìm lại Sora và A-chan.

Matsuri lẩm bẩm, mặt đỏ bừng trước tình cảnh oái oăm mà chính cô đang phải trải qua: ``Làm sao bây giờ\dots\ Buồn tè quá\dots'' Một tay cô vẫn nắm chặt túi bánh mì, giờ đã nguội ngắt, tay còn lại thì khẩn trương giữ lấy hạ bộ trong nỗ lực cuối cùng để cầm cự.

Cô đi loanh quanh, hết rẽ phải rồi rẽ trái, suốt hơn nửa tiếng đồng hồ nhưng vẫn chẳng tìm được hai người kia, cũng như chẳng thấy được bóng dáng của nhà vệ sinh nào.

Giờ đây, nàng Cổ động trở nên tuyệt vọng: ``Tệ rồi đây\dots\ Mình còn không mang điện thoại đi mới đau chứ!''

Cô ngửa mặt lên trời, như thể đang trách cứ số phận: ``\textbf{Hai người kia đi đâu rồi không biết!?}''

Cuối cùng, nhận ra rằng việc đi vòng vòng không mang lại kết quả gì, Matsuri quyết định chạy thẳng về quầy thu ngân. Cô lục túi lấy vài đồng yên lẻ, dùng tiền túi của mình trả cho túi bánh mì để có thể rời khỏi siêu thị mà không vi phạm quy định.

Nhưng ngay khi cô đang định tiến tới quầy thu ngân\dots

\il

\begin{center}
    \uline{Quay lại một khoảng thời gian trước đó.\\Khu gia vị.}
\end{center}

Sora và A-chan, với chiếc xe đẩy thứ hai đã đầy ắp hàng hóa, đang hoàn thành nốt danh sách mua sắm. Các món đồ như dầu ăn, nước tương, và hạt nêm được đặt gần nhau trên cùng một kệ, giúp hai người tiết kiệm được khá nhiều thời gian.

Nàng Tăng ca cầm lên một gói hạt nêm, cười nhẹ: ``Nhìn này, mua cái này còn được tặng kèm một chai nước mắm nhỏ nữa. Quá tiện lợi!''

Nàng Thần tượng gật đầu, đặt món đồ vào chiếc xe đẩy: ``Một công đôi việc luôn nhỉ. Như thế là xong rồi đấy\dots''

Cả hai người nhìn quanh, vẻ mặt có chút băn khoăn. Trong suốt khoảng thời gian mua đồ, họ vừa hoàn thành danh sách mua hàng, vừa đi tìm Matsuri, nhưng đến khi họ đã xong xuôi thì cô nàng vẫn chưa được tìm thấy.

\begin{dialogue}
    \speak{Sora} Lạ thật đấy. Matsuri-chan đâu nhỉ?
    \speak{A-chan} Không biết nữa. Lần cuối tôi thấy cô ấy là ở quầy bánh mì. Hay là cô ấy vẫn còn ở đó?
\end{dialogue}

Sora mỉm cười: ``Vậy chúng ta thử quay lại quầy bánh mì xem nào.'' Với chiếc xe đẩy nặng trĩu, hai người nhanh chóng di chuyển về phía quầy bánh mì, hy vọng rằng cô nàng Lễ hội vẫn đang chờ họ ở đó.

Bất chợt, họ nghe thấy tiếng vang lên thất thanh với giọng điệu quen thuộc: ``\textbf{A-chan! Sora-chan! Hai người đâu rồi!?}''

Cả hai cô nàng liền cau mày. ``Cậu có nghe thấy gì không? Hình như ai đó đang gọi tên chúng ta sao?'' A-chan tỏ rõ sự thắc mắc.

Sora nghiêng đầu: ``Hử? Em không để ý lắm\dots\ Hay là Matsuri-chan?''

Họ dừng lại một lúc lâu, cố gắng đoán xem đó có phải là giọng của cô nàng không. Một lát sau, một giọng hét thất thanh khác vang lên, cũng với giọng điệu đó: ``\textbf{Hai người kia đi đâu rồi không biết!?}''

Đến lúc này, họ mới nhận định chính xác rằng giọng nói đó là Matsuri. Nàng Thần tượng thì tỏ vẻ lo lắng cho đàn em của cô, còn nàng Đeo kính thì thở dài: ``Chắc chắn là Natsuiro-san rồi. Được rồi, quay lại tìm cô ấy thôi, không lại làm ầm cả siêu thị lên bây giờ.''

Cả hai nhanh chóng đẩy xe về khu vực quầy bánh mì, và cảnh tượng đầu tiên họ thấy khiến họ không khỏi sững sờ: Matsuri đứng giữa dòng người đông đúc đang chờ thanh toán, tay vung vẩy túi bánh mì như đang biểu diễn, còn cả người thì nhảy nhót tại một quầy thu ngân gần đó.

Thậm chí họ còn có thể nghe thấy cô nàng kia đang lẩm bẩm, hoàn toàn bẵng đi sự xấu hổ của một người con gái: ``Mót tè, mót tè, mót tè\dots'' Hành động kỳ lạ của cô thu hút sự chú ý của tất cả mọi người xung quanh. Một số người nhìn cô với ánh mắt ngơ ngác, một số khác lại tỏ vẻ vừa buồn cười vừa khó hiểu.

Sora cất tiếng gọi lớn, vẫy tay gây sự chú ý của cô nàng: ``A, Matsuri-chan!''

Nghe tiếng gọi quen thuộc, cô nàng tóc nâu lập tức quay đầu lại. Vừa nhìn thấy hai người bạn của mình, cô liền hét lên với vẻ trách móc: ``Mồ! Hai người đi đâu vậy, để em đi loanh quanh tìm hai người suốt!''

A-chan đẩy kính, thở dài bất lực. Cô chẹp miệng: ``Thấy chưa, tôi đã bảo sẽ có chuyện không ổn mà.'' Sora thì mỉm cười, tỏ vẻ hối lỗi với đàn em mình: ``Chị xin lỗi em nha. Em đứng chờ ở đây lâu chưa?''

Matsuri vẫn xách túi bánh mì, vừa nhảy nhót tại chỗ: ``Em cũng vừa mới đến thôi- Mà, hai chị ở đây rồi, để em đi tìm nhà vệ sinh\dots\ Sắp tè ra quần rồi!''

A-chan và Sora nhìn nhau, thầm hiểu rõ tình cảnh ``ngàn cân treo sợi tóc'' của cô bạn mình. Họ cũng muốn giúp, nhưng vấn đề lại nằm ở chỗ: thứ nhất, họ không biết nhà vệ sinh ở đâu, và thứ hai, hàng người chờ thanh toán trước mặt họ vẫn còn dài dằng dặc. Chưa kể rằng đống đồ mà họ muốn thanh toán cũng không phải ít ỏi gì.

Nàng Đeo kính chỉ có thể thở dài mà trấn an: ``Chắc cô phải nhịn thêm tí nữa thôi. Tôi không nghĩ ở đây có nhà vệ sinh đâu.'' Nàng Cổ động nhìn chị quản lý, giọng nghẹn ngào, mếu máo nói: ``Nhịn\dots\ nhịn sao\dots?''

Bất lực trước tình cảnh hiện tại, cô nàng đành phải nghe theo lời cấp trên, tay ôm túi bánh mì và gương mặt đầy khổ sở.

\ct

Sau khoảng gần mười phút đứng xếp hàng, cuối cùng cũng đến lượt họ. Matsuri, lúc này gần như mất hết kiên nhẫn, liên tục thúc giục nhân viên thu ngân làm thật nhanh để cô có thể đi vệ sinh. A-chan và Sora đứng bên cạnh, vừa phải giúp cô giữ bình tĩnh, vừa cố gắng xử lý tình huống.

Thỉnh thoảng, nàng Lễ hội rên rỉ: ``Nhanh\dots\ Nhanh lên đi mà\dots'', để rồi Sora và A-chan vỗ vai giúp cô bình tĩnh lại.

Thế nhưng, vì số hàng A-chan và Sora mua quá nhiều, cùng với việc có rất nhiều loại mặt hàng khác nhau, đã ba phút trôi qua mà chỉ mới quét được hơn một phần ba.

Nhân viên thu ngân nhìn số hàng khổng lồ mà ba người họ mua, lắc đầu rồi nói với vẻ bất lực: ``Hai người mua nhiều loại hàng quá\dots''

Sora gãi đầu, tỏ vẻ sự áy náy: ``Em xin lỗi chị nha. Tại quản lý của chúng em không nói rõ nên mua loại hàng nào\dots'' A-chan cũng cúi đầu xin lỗi cô ấy: ``Đã làm phiền cô nhiều rồi.''

Nàng Đeo kính quay sang Matsuri, tiếp tục động viên cô nàng đang sắp ``vỡ đê'' ấy: ``Natsuiro-san, cô có nghĩ cô chịu được nữa không?''

Matsuri nghiến răng, nước mắt tuôn rơi, cất lên giọng như sắp khóc: ``Em sắp chịu hết nổi rồi\dots''

A-chan nhìn khung cảnh xung quanh họ. Có quá nhiều người cũng tới mua hàng giống như cô, ai ai cũng bày tỏ sự sốt ruột. Cô chỉ có thể lắc đầu với nàng Lễ hội đó: ``Tôi cũng muốn cho cô xả lắm, nhưng ở đây đông người quá\dots''

Sora ngoảnh đầu lại, cúi người xuống, tiếp tục vỗ vai: ``Chịu khó thêm một tí nữa nha, Matsuri-chan!''

Matsuri chỉ có thể thở dài, nói với giọng tiu nghỉu: ``Rồi\dots''

\ct

Sau khi nhân viên siêu thị quét xong món cuối cùng, cô quay sang thông báo: ``{\color{red} Của ba người hết 2.140.000 đồng ạ.}''

Cả Sora, A-chan và Matsuri (vẫn đang tuyệt vọng giữ chặt hạ bộ) đều tỏ vẻ ngạc nhiên.

    {\color{red}
        \begin{dialogue}
            \speak{Sora} Ể? Đồng ạ?
            \speak{A-chan} Chứ không phải là yên à?
        \end{dialogue}}

Không để mọi người cũng như những khách hàng tiếp theo chờ đợi lâu, vị nhân viên thu ngân giải thích thêm: ``Nếu quy đổi ra yên thì sẽ là 10.080 yên ạ. Siêu thị này chấp nhận mọi loại tiền giao dịch, nên chị không phải lo đâu.''

Matsuri, trong tình trạng không thể đợi thêm, vội vàng nhắng nhít: ``Gì cũng được, hai người trả tiền đi nhanh lên!''

Lời nói đi kèm với hành động khẩn cấp: hai tay cô giữ chặt hạ bộ, mặt nhăn nhó như muốn khóc, trong khi quần lót bắt đầu rỉ nước. Cô biết mình sẽ không thể chịu đựng thêm dù chỉ một phút nữa.

Nàng Đeo kính nhanh chóng lấy ví, rút tiền mặt và đưa cho nhân viên đủ số thanh toán. Sau đó, nhân viên thu ngân cũng đưa lại cho cô tờ thanh toán dài như tờ sớ mà các Táo Quân hay mang lên trình cho Ngọc Hoàng khiến cả nhóm ngỡ ngàng trong giây lát.

A-chan ngạc nhiên: ``Tờ hóa đơn gì mà dài thế này!?'' Sora cũng sửng sốt không kém: ``Dài hơn cả danh sách mua hàng luôn ấy!''

Nữ nhân viên thở dài: ``Thì bởi vì quý khách mua rất nhiều loại hàng, nên tờ hóa đơn mới dài như vậy đó ạ.''

Sora cười trừ. Cả ba người họ đẩy xe hàng đi, trước khi không quên gửi lời chúc mừng năm mới cho cô nhân viên. Sau khi đưa tờ hóa đơn cho bảo vệ để đóng dấu ``đã thanh toán,'' cả ba người cuối cùng cũng rời khỏi khu mua sắm.

Khi vừa thoát khỏi biển người, họ thấy Kuon và Pekora cùng ông Ryujin đang ngồi ở ghế chờ, cắm mặt vào điện thoại.

A-chan vẫy tay gọi: ``Bọn tôi xong rồi đây.'' Sora cũng chiêm vào, tỏ vẻ ngại thay cho ba người: ``Hai người chờ bọn em có lâu không ạ?''

Cậu Tổng ngẩng mặt lên nhìn, rồi khẽ thở dài: ``Cũng khá lâu đấy. Nhưng mà biết sao được, Tết mà, đông là phải.'' Cậu quay sang người bố vẫn đang cặm cụi vào điện thoại: ``{\color{red} Bố ơi, họ xong rồi đấy.}''

Ông Ryujin gật đầu, nhìn mọi người tập trung đầy đủ liền lên tiếng: ``Vậy chuẩn bị đi về thôi, kẻo lại-''

Nhưng câu nói của ông bị Matsuri cắt ngang, giọng đầy tuyệt vọng: ``Nhà vệ sinh ở đâu vậy!?'' Nàng Thỏ chỉ tay về bên phải, nói: ``Kia kìa, peko. Tôi cũng vừa đi xong.''

Không cần đợi thêm, Matsuri chạy như bay về phía chỉ dẫn, miệng không ngừng hét: ``Sắp đái ra quần rồi! Ra mất thôi\~{}\dots''

Nhìn dáng vẻ hoảng loạn của cô nàng, Kuon chỉ tủm tỉm cười. Sau đó, cậu quay sang xem xét xe đẩy đầy ắp hàng hóa mà A-chan và Sora đã mua, không khỏi cảm thán: ``Hai người mua nhiều đồ quá!''

Nàng Đeo kính tỏ vẻ thản nhiên, dù có một chút bất bình: ``Bởi cậu có chỉ cho chúng tôi biết cụ thể là mua đồ nào với số lượng bao nhiêu đâu.'' Nàng Thần tượng chiêm vào: ``Vậy nên bọn em mới mua mỗi loại một cái đó.''

Nhìn đống hàng từ hai chiếc xe, Kuon tỏ vẻ lo ngại: ``Nhưng như thế liệu có ăn hết không?'' A-chan chỉnh lại cặp kính, rồi trừng mắt nhìn cậu: ``Cậu nhận ra không chỉ mỗi gia đình cậu ăn Tết chứ?''

Đến cả Pekora cũng nhướng mày: ``Chính bố anh vừa nói thế rồi mà, peko. Phải gió à?'' Sora thì gật đầu khẳng định: ``Còn có cả bọn em nữa chứ.''

Ryujin nhìn vào xe đẩy, cũng tỏ vẻ hơi ngạc nhiên: ``Nhưng như thế thì cũng nhiều quá! Thế các cháu mua hết bao nhiêu?''

Nàng Thần tượng rút tờ hóa đơn ra đưa cho ông ấy. Khi nhìn thấy tờ hóa đơn dài chạm đất, ông giật mình: ``Tận từng đấy cơ á!?''

Cậu con trai ngó vào, cũng phụ họa: ``Vãi\dots\ Tận triệu bạc chứ ít gì\dots\ Hôm nay đi chợ thiệt hai tận hơn sáu triệu đấy ạ\dots''

``Đúng như dự đoán của bố con mình luôn chứ?'' ông bố cười trừ. Ngay khi ông vừa định hỏi tiếp, Matsuri từ xa chạy trở lại, gương mặt giãn ra đầy nhẹ nhõm: ``Hồ\~{} Đi xong nhẹ cả người\~{}''

Ryujin bật cười vui vẻ: ``Thôi, ta cũng về thôi nhỉ? Xem mấy đứa kia làm thế nào.''

Cả đội năm người đẩy một xe đầy hàng hóa rời khỏi siêu thị, chuẩn bị cho ngày Tết cận kề. Trên đường đi, Matsuri vẫn lẩm bẩm: ``Em thấy mọi người mua nhiều đồ quá\dots\ Trông như đang tích trữ đồ trước thảm họa ấy. Sao không chờ đến mấy ngày nữa rồi mua?''

Đáp lại, Kuon chỉ nở một nụ cười đầy ẩn ý: ``Rồi em sẽ biết thôi.''

Trời đã dần tối, ánh đèn từ những dãy nhà và phố xá sáng rực lên, báo hiệu một ngày sắm sửa tất bật sắp kết thúc. Cả đội chậm rãi rời khỏi khu mua sắm, tiếng cười đùa râm ran suốt cả quãng đường.

Nhìn lại chiếc xe đẩy chất đầy hàng, Matsuri thở dài ngao ngán: ``Chắc mang chỗ này về, chúng ta phải ăn cả tháng mất!''

Cậu Tổng cười lớn: ``Vậy thì tốt thôi, ít nhất sẽ không có ai trong công ty ta bị đói. Chứ đừng như ai đó nhịn đói cả tháng mà không rõ lý do gì cả\footnote{Xem lại {\flaregothic ÉPISODE 33}, quyển số Ba}.''

Nàng Cổ động trề môi, nhưng ánh mắt lộ rõ vẻ háo hức. Dù mệt mỏi, tất cả đều trông đợi vào những ngày Tết sắp tới, một khoảng thời gian hiếm hoi để họ cùng nhau quây quần, sẻ chia niềm vui và những câu chuyện thú vị.

Họ nhanh chóng lên xe, để lại đằng sau khung cảnh náo nhiệt của siêu thị trong dịp cuối năm, chuẩn bị trở về nhà với một cái Tết đầu tiên mà không ai có thể quên.

\vspace{-40pt}
\subsubsection*{\phantomsection\label{ep-003.IVd}}
\addcontentsline{toc}{subsubsection}{Okayu và Korone}

\begin{center}
    {\color{T4} Nhóm 4 - Okayu và Korone}
\end{center}

Ngay sau khi được Kuon giao nhiệm vụ, Okayu và Korone liền bắt tay vào công việc sơn nhà, không chần chừ lấy một giây. Điều bất ngờ đầu tiên là các lọ sơn và chổi sơn mà họ dùng dường như xuất hiện từ\dots\ hư không. Đến cả Kuon lẫn các thành viên khác của \hll\ cũng chỉ biết tròn mắt ngạc nhiên mà không tài nào hiểu được cặp đôi này lấy dụng cụ từ đâu.

Cách thức ``làm việc'' của nàng Khuyển và nàng Miêu cũng khác biệt hoàn toàn với bất kỳ định nghĩa lao động truyền thống nào. Không chỉ đơn thuần cầm chổi quét sơn lên tường, họ còn nghĩ ra đủ trò quái lạ nhằm ``biến công việc thành niềm vui'' - một niềm vui có phần hơi\dots\ vượt ngoài tầm kiểm soát.

Ví dụ, khi làm ở tầng 5, nàng Khuyển như hóa thân thành một cơn lốc sơn màu. Cô nàng cầm thùng sơn chạy vòng quanh căn nhà, tốc độ nhanh đến mức vượt xa cả lực hút của Trái Đất, giúp cô dễ dàng bay lên trần nhà và sơn toàn bộ bề mặt phía trên. Trong khi đó, nàng Miêu theo sát phía sau với một chiếc máy sấy công suất cực lớn, làm khô lớp sơn ngay tức thì. Kết quả là bức tường và trần tầng 5 trông như vừa được phủ bởi một lớp áo mới sáng bóng, đều mịn đến từng chi tiết.

Tầng 4 cũng không kém phần thú vị. Khi Aki, Kaien, và Ryuuhou vừa mang quần áo lên phơi ngoài ban công thì Okayu và Korone đã nhanh chóng xuất hiện để tiếp tục công việc sơn lại căn phòng. Theo chỉ dẫn của Kuon, cặp đôi này sơn toàn bộ không gian thành màu xanh nước biển, đem lại cảm giác tươi mát như đang ở giữa đại dương. Điều đặc biệt là, bất chấp sự ``ngẫu hứng'' trong cách làm, mọi đồ đạc trong phòng, từ rèm cửa, kệ sách đến cả mấy chiếc ghế sofa, đều sạch sẽ không vấy một giọt sơn.

Nếu hỏi họ làm thế nào, câu trả lời hẳn sẽ khiến bất cứ ai cũng phải ngẩn người: họ dùng\dots\ bình phun sơn lấy ra từ ứng dụng vẽ trên máy tính. Cái bình nhỏ xíu, chỉ to cỡ hai gang tay, lại có khả năng phủ màu đều đặn đến mức khó tin. Và cái quan trọng nhất: không hề có dấu hiệu chúng sẽ bị đổ, văng tung tóe hay làm bẩn bất cứ đồ vật nào trong phòng.

``Thật kỳ lạ,'' đó là những lời lẩm bẩm trong sự kinh ngạc của Kuon khi nhìn đôi bạn vừa hoàn thành công việc mà như đang tham gia một chương trình tạp kỹ. Nhưng cuối cùng, cậu chỉ biết cười bất lực và lặng lẽ ghi nhận khả năng ``đặc biệt'' của cặp đôi này.

\begin{center}
    \uline{Tầng 3.}
\end{center}

Tuy nhiên, việc họ thực hiện công việc sơn tường ở tầng 3, lại đánh dấu một bước ngoặt khác biệt trong công việc của Okayu và Korone. Nếu ở các tầng trên, họ vẫn giữ được chút hình thức ``làm việc nghiêm túc'' (dù không giống cách làm thông thường), thì với tầng 3, mọi thứ lại chuyển hẳn sang một màn biểu diễn đầy sáng tạo và hài hước.

Vừa bước vào căn phòng rộng thênh thang, nàng Mèo không khỏi cảm thán: ``Căn phòng này rộng ghê ha.''

Nàng Chó, với bản tính tinh nghịch vốn có, ngay lập tức nảy ra một ý tưởng lạ lùng, cúi người thì thầm với bạn của mình: ``Thế này mà sơn bình thường thì chán lắm. Hay là thế này\dots''

Cô ghé sát tai, nói thì thào với người bạn của cô. Okayu lập tức phấn khích, đôi mắt sáng lên như tìm được đồng minh trong trò đùa mới: ``Ồ\~{} Giống như hồi ta stream hôm trước hả?''

Korone gật đầu: ``Đúng đó hen\~{}!''

Nói là làm, cả hai chạy ra ngoài để kiếm ``dụng cụ sơn nhà đặc biệt''. Không lâu sau, họ quay lại với cây chổi sơn dài, kèm theo vài khẩu súng bắn sơn mà chẳng ai biết họ nhặt được từ đâu - có vẻ như chúng cũng xuất hiện từ\dots\ không khí.

Và thế là, công việc sơn nhà bỗng hóa thành một trận chiến giả lập lấy cảm hứng từ Splat(oo)n, tựa game yêu thích của họ.

Nàng Khuyển chọn cho mình thùng sơn màu nâu, nhúng cây chổi dài vào trong và xoay tít như một chiến binh cầm giáo ra trận. Còn nàng Miêu, vốn thích cảm giác hành động tốc độ cao, nạp đạn sơn tím vào khẩu súng liên thanh, mắt ánh lên sự tinh quái.

\begin{figure}[H]
    \centering
    \begin{subfigure}[t]{0.45\textwidth}
        \centering
        \includegraphics[width=8cm]{../../../../Image/Book?/c1-s1.png}
    \end{subfigure}
    \quad
    \begin{subfigure}[t]{0.45\textwidth}
        \centering
        \includegraphics[width=8cm]{../../../../Image/Book?/c1-s2.png}
    \end{subfigure}
    \caption{Bức hình của cặp đôi OkaKoro, sử dụng cho Splatoon.}
\end{figure}

Họ biến căn phòng thành một đấu trường thực thụ. Những vật dụng như bàn kính, rèm cửa, và thậm chí cả nhà vệ sinh được tận dụng làm nơi ẩn náu lẫn chướng ngại vật. Không khí trở nên sôi động hơn bao giờ hết khi từng cú quất chổi và những loạt đạn sơn bắt đầu nhuộm bức tường đầy sắc màu.

Korone nhìn vũ khí của Okayu, tỏ vẻ bất bình: ``Ể\~{}? Okayu không công bằng! Bà có súng kìa!''

Okayu nở nụ cười, nhún vai đáp lại: ``Bà cũng có cả tá lựu đạn sơn đấy thôi? Với cả, cờ tay ai người ấy phất mà?''

Nàng Khuyển liếc xuống thắt lưng của mình, nơi treo vài quả lựu đạn sơn màu cam rực rỡ, bật cười: ``À, đúng rồi nhỉ\dots''

Nàng Miêu giơ súng lên cao, hô hào: ``Đây chỉ là một trận giao hữu thôi. Cứ thoải mái đi nha\~{}''

Cả căn phòng nhanh chóng trở thành chiến trường đầy ắp tiếng cười. Mục tiêu chính của họ rõ ràng không phải là sơn tường, mà là tận hưởng niềm vui từ trò đùa tinh nghịch này.

\ct

``Koro-san đi đâu rồi ta\~{}?'' Okayu vừa đi lại giữa nửa căn phòng bên trái vừa gọi với giọng chọc ghẹo, đôi mắt láo liên tìm kiếm. Ở phía bên kia, Korone đang nấp sau một cây cột lớn, cố gắng giữ im lặng nhưng không giấu nổi nụ cười tinh nghịch.

``Okayu tưởng thế mà chiến dữ\dots'' nàng Khuyển tròn mắt ngạc nhiên, vừa thì thầm vừa cố kiềm chế tiếng cười. Sở dĩ cô nói vậy là vì ngay khi trận đấu bắt đầu, nàng Miêu đã tung loạt tấn công phủ đầu khiến không chỉ bức tường mà cả căn phòng tràn ngập những vệt sơn tím. Đến cả Korone cũng bị dồn vào chân tường, không kịp trở tay.

Vũ khí trong tay Korone giờ đây có phần ``khiêm tốn'': một cây chổi sơn dài, vài quả lựu đạn sơn màu cam và\dots\ cô chợt nhớ ra, trên lưng mình còn giắt một cây kiếm nhựa, món đạo cụ còn sót lại từ lần diễn tập đóng vai thành các nhân vật của Splat**n hôm trước.

Nhân lúc nàng Mèo vẫn còn đang mải nhìn quanh tìm kiếm, cô tranh thủ lấy một quả lựu đạn sơn, kéo chốt rồi ném thẳng về phía đối thủ. Tiếng ``cục cục'' nhẹ vang lên khi quả lựu đạn lăn tới gần đối thủ. Nàng Miêu quay ngoắt lại, đôi tai vểnh lên cảnh giác, nhưng đã quá muộn. Một vụ nổ *BỤP* phát ra, và toàn thân Okayu ngay lập tức bị phủ đầy sơn màu cam rực rỡ.

Korone cười phá lên, không thể kiềm chế trước cảnh tượng của người bạn mình. Nhưng niềm vui của nàng Khuyển chẳng kéo dài lâu khi Okayu, trong lúc bị ``sơn toàn thân'', đã tranh thủ nạp đạn và bắn trả. Loạt đạn tím bắn trúng nàng Chó, khiến cô vừa bị bật ngã ra sau, vừa trở thành mục tiêu nhuộm tím hoàn hảo.

Và thế là, trận chiến sơn nhà giữa hai cô gái lại tiếp tục. Tiếng cười vang lên khắp căn phòng khi họ dùng tất cả những gì có thể: súng, lựu đạn, và thậm chí cả\dots\ thùng sơn. Mỗi cú ném, mỗi phát bắn đều khiến căn phòng tầng ba tràn ngập sắc nâu và tím.

Khi trận chiến kết thúc, cả hai người đều không thể phân biệt được ai thắng ai thua. Người nào người nấy cũng lấm lem từ đầu đến chân, quần áo chẳng khác gì một bảng pha màu di động. Bức tường căn phòng, vốn là mục tiêu ban đầu, giờ đây lại trở thành một kiệt tác trừu tượng với những vệt sơn đan xen như một tác phẩm nghệ thuật hiện đại.
Bức tường tầng 3 được khoác lên một diện mạo mới, không theo bất kỳ quy chuẩn nào, nhưng tràn đầy sức sống và dấu ấn của hai cô gái.

Korone ngó quanh một lượt, rồi ngơ ngác hỏi: ``À-rế? Chúng ta đã sơn xong căn phòng này từ khi nào vậy?'' Okayu đáp lại, vẫn giữ nụ cười điềm tĩnh nhưng không giấu được sự tinh quái: ``Tôi cũng không biết đâu nha\~{} Chưa gì chúng ta cũng đã hoàn thành nhiệm vụ rồi đấy.''

Cặp đôi Chó-Mèo đi xung quanh bãi chiến trường - cũng là kiệt tác của hai cô nàng vừa tạo ra một hồi lâu. Họ gật gù, tỏ vẻ tâm đắc với bức tranh đầy màu sắc bên trong văn phòng mà bọn họ vô tình tạo ra.

``Cơ mà\dots'' nàng Miêu bất chợt lên tiếng, ``nếu chỉ chiến trong phòng này thì cũng chán nhỉ?'' Korone phấn khích đáp: ``Đúng ha\~{} Hay là ta tiếp tục ở bên ngoài đi!''

Và thế là, không chút chần chừ, hai cô gái nhanh chóng thu dọn hiện trường. Những vết sơn trên bức tường thì được để nguyên như một bằng chứng cho ``thành quả lao động'' của họ, trong khi các dụng cụ sơn và đồ vật khác được sắp xếp lại gọn gàng. Sau đó, họ kéo nhau ra ngoài, sẵn sàng cho một ``cuộc chiến mới'' với những ý tưởng không ai đoán trước được.

\begin{center}
    \uline{Ngoài đường - Trước cửa nhà Hikawa.}
\end{center}

Chiều cận Tết, ngoài đường vẫn còn khá đông người qua lại.

Khi cặp Chó-Mèo tiếp tục vờn nhau bằng những trò chơi ``sơn bắn'', nhiều người đi đường nhìn vào cũng chỉ nghĩ rằng đây là hai cô gái trẻ cải trang thành thú đang chơi đùa vô hại. Một vài người thậm chí còn dừng lại một lát, cười và thầm nghĩ: ``Tuổi trẻ đúng là năng động thật!''

Nhưng chỉ chơi đùa với nhau thì đâu đủ, Okayu và Korone quyết định nâng cấp trò chơi lên một tầm cao mới: thi xem ai có thể sơn xong một tòa nhà trước. Và như thường lệ, họ không chọn tòa nhà nào khác ngoài nhà của cậu Tổng Kuon và cả nhà kế bên, vốn là một cửa hàng bán đồ điện tử.

Nàng Miêu nhiệm vụ sơn nhà của Kuon, trong khi bạn của cô đảm nhiệm căn nhà kế bên. Họ quyết định sử dụng chung một loại sơn màu xanh nước biển để đảm bảo ``tính công bằng''.

Okayu giơ súng sơn của mình lên và cười nhếch mép: ``Sẵn sàng chưa, Koro-san?'' Korone, đang cầm một quả bom sơn, gật đầu đáp: ``Rồi đó, Ogayu\~{}!''

Và thế là, cuộc chiến bắt đầu!

Nàng Miêu tận dụng kho ``vũ khí'' đa dạng của mình, từ súng tiểu liên sơn, súng phóng tên lửa, cho đến một khẩu đại bác bóng sơn khổng lồ. Trong khi đó, nàng Khuyển chọn cách chơi táo bạo hơn với hàng loạt lựu đạn và bom sơn.

Họ chia cuộc thi thành từng hiệp ngắn, mỗi hiệp kéo dài khoảng sáu mươi giây, và tận dụng lúc đường phố tạm vắng người để thực hiện ``công phá''. Với nhà của Kuon, mọi chuyện diễn ra khá suôn sẻ. Các thành viên \hll\ bên trong, dù đã quen với những trò vô lý của cặp đôi này, đều chọn cách phớt lờ. Hai nhân viên trực thì đã bị Rushia vô tình ``vô hiệu hóa'' từ trước, nên chẳng ai can thiệp.

Tuy nhiên, tình hình ở căn nhà kế bên lại hoàn toàn trái ngược.

Korone, trong một lần quá hăng say, đã vô tình ném lựu đạn quá mạnh, làm vỡ tan cửa sổ tầng ba của cửa hàng điện tử. Quả lựu đạn nổ tung, khiến sơn bắn tung tóe khắp căn phòng bên trong. Một người đàn ông, rõ ràng là chủ nhà, ló đầu ra với khuôn mặt đầy phẫn nộ: ``{\color{red} \textbf{Đứa nào làm bẩn căn phòng của ông thế này!?}}''

Nhưng tiếng hét của gã dường như rơi vào khoảng không, hoàn toàn bị ngó lơ bởi cả hai cô nàng \textit{GAMERS}.

``Lượt tiếp theo nào, Ogayu!''

``OK\~{}''

Korone đã nhanh chóng ném thêm một quả lựu đạn nữa, lần này lại\dots\ vỡ luôn cửa sổ còn lại. Đến đây, người đàn ông từ tòa nhà bên kia đã lên cơn tức tối: ``{\color{red}\textbf{Được lắm! Để đây ông trị tội bọn bây!}}''

Thế nhưng, vừa dứt lời, ông lại bị một quả lựu đạn khác đáp thẳng vào mặt. Sơn bắn tung tóe, phủ kín cả người ông bằng màu xanh. Và như thể chưa đủ, ông còn ``ăn'' thêm vài phát súng sơn từ Okayu, khiến ông ngã nhào ra phía sau, không thể thốt lên thêm lời nào.

Trận chiến ``Pê-U-Bê-Gờ phiên bản sơn nhà'' tiếp tục trong tiếng cười sảng khoái của cả hai. Những quả bom sơn, lựu đạn, và đạn sơn liên tục bay ra, biến hai tòa nhà thành những ``bức tranh hiện đại'' kỳ quái. Sau mười hiệp đấu, cả hai ngôi nhà đều được phủ kín màu xanh nước biển, nhưng đổi lại là hàng loạt ô cửa sổ bị vỡ, sơn vương vãi khắp nơi, thậm chí còn chui hẳn vào bên trong cả hai tòa nhà, và vài người đi đường bị dính đạn lạc phải lủi vội.

Nàng Mèo tím chống tay vào đầu gối, nhìn thành quả và mỉm cười: ``Koro-san khá đấy, sơn nhà nhanh thật đó!'' Nàng Khuyển nâu nghe cô bạn nói vậy, bật cười: ``Ogayu cũng giỏi mà! Mà bà kiếm đâu ra quả đại bác to thế này vậy?''

Okayu đáp: ``Lấy từ trong game bắn xe tăng ra đấy.'' Korone nghiêng đầu ngơ ngác: ``Ể\dots\ L-Là cái trò nào ấy nhỉ?''

``P*cket T*nk. Có lẽ bà không biết đâu. Cái trò đó từ thời Tống rồi.''

Sau một hồi cười đùa, cả hai cùng thu dọn bãi chiến trường (dĩ nhiên là trừ bức tường sơn) và quay vào trong nhà, sẵn sàng cho nhiệm vụ tiếp theo. Một buổi chiều trước Tết với cặp đôi Chó-Mèo chưa bao giờ ``bình thường'' đến thế.

Sau khi hoàn thành ``trận chiến'' kéo dài gần hết buổi chiều, Okayu và Korone trở lại căn nhà của cậu con trai. Cả hai trông không khác gì những bức tranh di động, với sơn tím sơn nâu bám đầy trên tóc, quần áo và cả khuôn mặt.

Okayu thả người xuống chiếc ghế gần cửa, tay cầm chai nước mát, thở phào nhẹ nhõm:
``Bà thấy sao, Koro-san?''

Korone đặt thùng dụng cụ xuống, ngồi bệt bên cạnh: ``Mệt thì mệt, nhưng mà vui quá xá! Lần đầu tiên tôi thấy sơn nhà mà giống chơi game vậy luôn!''

Nàng Miêu cười nhẹ: ``Đúng là bọn mình làm kiểu gì cũng ra 'game' được nhỉ. Mà lần sau chắc phải nghĩ trò gì đó ít gây rắc rối hơn\dots''

Korone bật cười: ``Rắc rối gì chớ? Chỉ là\dots\ có chút xíu tổn thất ngoài ý muốn thôi mà!''

Cả hai cô gái tai thú cười một trận sảng khoái, để lại phía sau là một buổi chiều đầy màu sắc, tiếng cười và những kỷ niệm đáng nhớ. Tầng ba được sơn xong, căn nhà kế bên cũng trông như mới, dù có vài ``tổn thất nhỏ''. Với cặp đôi Chó-Mèo, đây có lẽ là một buổi chiều cận Tết đặc biệt nhất mà họ từng trải qua.

\vspace{-40pt}
\subsubsection*{\phantomsection\label{ep-003.IVe}}
\addcontentsline{toc}{subsubsection}{Mio, Choco, Ayame}

\begin{center}
    {\color{T5} Nhóm 5 - Mio, Choco, Ayame}\\Phòng bếp - Phía của Kuon.
\end{center}

Mio, Ayame và Choco đi lên tầng trên, nơi được Kuon giao phó cho họ dọn dẹp. Khi mở cửa căn phòng, cả ba người đều không khỏi bất ngờ trước khung cảnh trước mắt: một căn bếp nhỏ hẹp, tường được lát bằng gạch gốm màu xám ngà đã cũ kỹ, loang lổ những vết bẩn của thời gian. Sàn nhà làm bằng gạch ngói dính đầy vết giày dép in hằn, còn bức tường trắng thì xuất hiện nhiều vết nứt nhỏ chạy dài như những đường chỉ.

Không gian này tuy đơn sơ nhưng chứa đựng không ít đồ đạc: một chiếc tủ lạnh cỡ vừa, một cái chạn treo đựng gia vị đã ngả màu vàng vì bụi bẩn, một lò chiên không dầu nằm kế bên bếp ga, tất cả được xếp trên một chiếc bàn sứ có chân sắt đã bị gỉ mọt. Phía dưới bàn là một cái lò nướng, lò vi sóng, và một nồi áp suất cũ kỹ. Chồng lên đó là vài cái chảo, cái nồi lẻ tẻ, và thêm cả những vật dụng lặt vặt như dụng cụ đánh trứng hay muôi múc canh.

Ở góc bếp là một bồn rửa chén bằng inox, trên thành bồn có hai chiếc thớt, một bằng gỗ và một bằng thủy tinh. Tuy nhiên, thứ khiến nàng Quỷ sừng phải nhướng mày kinh ngạc chính là\dots\ một chiếc máy giặt được đặt ngay cạnh bồn rửa, tạo cảm giác vô cùng lạc lõng.

Cô quay sang nhìn hai người bạn đồng hành, không giấu nổi sự tò mò: ``Ai đời lại để máy giặt trong bếp thế này? Nhà anh ấy chật đến mức không còn chỗ nào khác à?''

Choco nhún vai, nhẹ nhàng đáp: ``Cũng có thể lắm. Nhà anh ấy chỗ nào tôi cũng thấy đồ đạc ngổn ngang. Với cả, có vẻ đặt ở đây thì tiện hơn cho việc phơi quần áo ngoài ban công ấy chứ.''

Mio, lúc này đang ngó nghiêng quanh phòng, cũng gật đầu đồng tình: ``Đúng đấy. Tôi còn để ý thấy sạp phơi đồ ngoài ban công kìa. Đặt máy giặt gần bếp thế này chắc là để tiết kiệm không gian thôi.''

Ayame chẹp miệng, liếc quanh căn bếp chật hẹp: ``Công nhận, trông thế này chắc khó mà nấu ăn được tử tế. Chả trách nhà Hikawa-san không ăn uống gì ở đây, dazo.''

Nàng Y tá nhẹ nhàng giải thích: ``Kuon-sama có nói à cả nhà thường dùng phòng khách làm nơi ăn uống luôn mà. Còn bếp này chắc chỉ để nấu nướng cơ bản thôi.''

Nàng Sói đen mỉm cười, xắn tay áo của mình lên: ``Thôi nào, đừng thắc mắc nữa. Việc của chúng ta là dọn sạch căn bếp này, đúng không?''

Nàng Oni lên tiếng thắc mắc: ``Thì đúng là như thế, nhưng vấn đề là chúng ta sẽ dọn dẹp chỗ này như thế nào?''

Nàng Sói đen chỉ nhún vai, giọng thản nhiên: ``Anh ấy có nói gì đâu. Tôi nghĩ cứ thấy bẩn chỗ nào thì dọn chỗ đấy thôi.''

Nghe vậy, cả ba người cùng đeo bao tay cao su, mỗi người cầm lấy một chiếc giẻ lau từ trên bàn. Họ bắt đầu bằng việc lau sạch bề mặt các vật dụng, rồi từ từ tiến đến từng ngóc ngách của căn bếp nhỏ bé. Nhưng chẳng ai ngờ rằng, công cuộc dọn dẹp này sẽ không hề đơn giản như họ nghĩ\dots

\ct

Năm phút sau, công việc lau bề mặt dường như đã hoàn tất. Chỉ trong thời gian ngắn, cả ba người đã lau sạch sẽ từ bàn bếp, tủ lạnh cho đến chạn bát đũa treo trên cao. Các vết bẩn loang lổ trước đó giờ chỉ còn là ký ức, để lại một không gian sáng bóng đến bất ngờ.

Choco hài lòng đặt chiếc giẻ lau xuống bàn, tay vỗ nhẹ vào nhau như thể hoàn thành một nhiệm vụ lớn: ``Thế là xong rồi đấy. Nhanh gọn thật!''

Mio đứng thẳng dậy, đưa mắt nhìn quanh căn bếp đã gọn gàng hơn hẳn, nhưng vẫn không giấu nổi sự ngờ vực: ``Sao em thấy mọi thứ dễ dàng thế nhỉ? Cứ như kiểu\dots\ ta chưa làm được gì cho ra hồn ấy.''

Ayame, vẫn còn đôi chút bồn chồn, lên tiếng: ``Tôi cũng cảm giác vậy. Hay là mình làm thêm cái gì đó đi? Chứ thế này thì nhàn quá\dots''

Ngay sau đó, cô nàng đã nảy ra một ý tưởng táo bạo. Đôi mắt cô sáng lên, hào hứng đề xuất: ``Hay chúng ta rửa bát đũa nhà anh ấy đi? Giống như đợt dọn nhà cuối năm ấy, nhớ không?''

Nàng Y tá nghiêng đầu, đôi mày hơi cau lại tỏ vẻ nghi hoặc: ``Nhưng mà, Kuon-sama có nói là cần mình làm đâu nhỉ? Với lại, nhìn bát đĩa cũng đâu có bẩn lắm.''

Mio, vốn là người thích mọi thứ gọn gàng, lên tiếng ủng hộ ngay: ``Cứ làm một thể đi. Đằng nào mình cũng đang dọn dẹp giúp anh ấy mà. Coi như là thêm chút chu đáo.''

Không chần chừ thêm, cả ba người bắt đầu lấy từng chiếc bát, đĩa, cốc chén từ trong chạn xuống, chất thành một chồng lớn trên bồn rửa. Ayame nhanh nhẹn đeo găng tay vào trước tiên, miệng hào hứng: ``Lần này tôi phụ trách rửa nhé! Còn hai người kia lo tráng và lau khô cho nhanh.''

Choco nhún vai đồng ý, trong khi Mio chỉ cười, đi tìm chiếc khăn sạch để sẵn sàng cho công việc. Công việc tưởng chừng như đơn giản, nhưng họ đâu ngờ rằng phía sau những chiếc bát đĩa này lại ẩn chứa vô số ``bất ngờ'' mà cả ba không hề chuẩn bị trước\dots

\ct

Trong lúc cả ba đang tập trung làm sạch đống bát đĩa, Mio bất ngờ cảm nhận được điều gì đó không ổn. Một cảm giác lạnh ướt bất thường dưới chân khiến nàng Sói khẽ rùng mình. Cô cúi xuống nhìn và giật mình khi thấy một vũng nước lớn đang lan ra dưới sàn.

``Mọi người ơi, có chuyện rồi!'' cô kêu lên đầy cảnh giác. Nghe thấy tiếng gọi, Choco và Ayame đồng loạt quay đầu lại.

Nàng Quỷ sừng nhướng mày hỏi: ``Sao thế, Mio? Chuyện gì vậy?'' còn nàng Y tá cũng tò mò bước tới: ``Sao em kêu to thế? Có gì bất thường sao?''

Nhưng khi cả hai nhìn xuống sàn, họ ngay lập tức hiểu ra vấn đề. Vũng nước dưới chân Mio đang ngày càng lan rộng, tạo thành một tấm gương phản chiếu ánh sáng từ bóng đèn trần.

Nàng Sói đen nhanh chóng lấy lại sự bình tĩnh, rồi lẩm bẩm: ``Đường ống nước bị làm sao à?'' Cô cúi xuống kiểm tra phía dưới bồn rửa. Tuy nhiên, mọi thứ ở đó vẫn vận hành bình thường, không hề có dấu hiệu rò rỉ.

``Lạ thật đấy. Nếu không phải đường ống, thì\dots'' cô chưa kịp nói hết câu, Sensei đã ngắt lời: ``Có khi nào là rãnh nước nhà anh ấy không?''

Nàng tóc trắng nhanh nhẹn kiểm tra khu vực rãnh nước nhỏ dưới bàn bếp. Đúng như dự đoán, dòng nước bị chặn lại bởi một khối bùn to đùng và vài mảnh rác vụn mắc kẹt. ``Đúng là rãnh nước của anh ấy đang tắc này, yo.''

Choco thở dài, lắc đầu: ``Lại là vấn đề rãnh nước. Tôi biết ngay mà.'' Nàng Quỷ sừng nhoẻn miệng cười đắc ý: ``Chắc chúng ta phải dọn cả rãnh nước luôn rồi. Đằng nào thì cũng đang dọn bếp mà!''

Không để lãng phí thêm thời gian, cả ba người cùng nhau kê bàn bếp sang một bên để lộ rãnh nước. Tuy nhiên, ngay khi bắt tay vào dọn, họ mới nhận ra rằng đây không chỉ đơn thuần là một nhiệm vụ lau dọn thông thường.

Trong rãnh nước có đủ mọi thứ kỳ quặc: một đống bùn đất, vài bao bì cũ, và thậm chí cả tai nghe không dây. Ayame hí hoáy bới rác, vừa nhặt lên vừa liến thoắng: ``Nhìn này, ở đây còn có cả tai nghe này! Ai mà lại làm rơi ở chỗ thế này nhỉ?''

Nàng Y tá nhíu mày, nhặt lên một túi ni-lông nhăn nhúm: ``Sao bếp anh ấy lại nhiều rác thế này? Không thể hiểu nổi.''

Mio thì khỏi phải nói, khuôn mặt nàng Sói méo xệch khi cầm lên một thứ gì đó mềm mềm, ươn ướt: ``Éc! Đây là đồ ăn thừa hay cái gì thế này? Khiếp quá!''

Choco cười an ủi, cố giữ cho Mio bình tĩnh lại: ``Chắc là trong lúc nấu ăn rơi lọt xuống thôi. Không sao đâu, Mio-sama. Bình tĩnh nào.''

Nhưng sự bình tĩnh đó không kéo dài lâu. Ayame bỗng thét lên khi nhặt được một vật thể đáng sợ: ``\textbf{Ôi giời ơi, còn có cả \uline{gián} này! Khiếp quá!}''

Tiếng hét của nàng Quỷ sừng khiến cả Choco và Mio đứng hình trong vài giây. Khi họ nhận ra Ayame thực sự đang cầm một con gián, cả hai đồng loạt hét lên đầy hoảng hốt: ``\textbf{Sao bà/em lại cầm nó lên!?}''

Nàng Quỷ sừng bình thản đáp, mặt không chút biểu cảm: ``Còn có cả một tổ gián dưới này nữa.'' Nghe vậy, cả hai cô gái càng khiếp sợ hơn, la hét đầy hoảng loạn.

Mio ôm mặt thét: ``Éc! Giết nó đi! Giết nó ngay đi!'', còn Choco thì lùi lại, giọng run rẩy: ``Bỏ nó ra đi! Khiếp quá!''

Không chút do dự, Ayame rút hai thanh katana quen thuộc từ sau lưng, chém thẳng vào con ``tiểu cường'' đáng thương. Với một nhát dứt khoát, nó ``lên đường'' ngay lập tức. Cả Mio và Choco thở dài nhẹ nhõm.

Sự nhẹ nhõm của hai cô gái lại không thể kéo dài được, khi từ góc rãnh nước, một đàn gián lúc nhúc chui ra, khiến không khí trở nên hỗn loạn hơn bao giờ hết.

Choco hét lớn, nhảy lên nàng Sói đen: ``\textbf{Ayame-sama! Cứu bọn này với!!}'', còn nàng Sói đen thì nhảy thẳng lên nóc tủ lạnh, cũng không kìm được sự sợ hãi: ``\textbf{Tại sao bếp nhà anh ấy lại nhiều gián như vậy chứ!?}''

Ayame, vẫn giữ vẻ mặt điềm nhiên, gõ nhẹ vào chuôi kiếm và tuyên bố đầy quyết tâm: ``Cứ để đấy cho tôi, yo! Chuyện nhỏ!''

\ct

Sau một hồi vật lộn với đàn gián đáng sợ, cả ba cô gái thở hổn hển giữa chiến trường bừa bộn. May mắn thay, với sự giúp đỡ kịp thời của Miko và Mel từ phòng bên, họ đã ``quét sạch'' cả một ``đại gia đình'' gián ẩn náu dưới rãnh nước. Tuy nhiên, kết quả để lại là một khung cảnh không thể thảm hơn.

Những vết chém rải rác khắp bức tường, tạo thành các đường nứt dài như dấu vết của một trận chiến khốc liệt. Một vài chiếc bát đĩa vừa rửa xong thì bị vỡ tan tành. Cái chạn treo tường, vốn được thiết kế để chịu lực tốt, giờ cũng không thể chống đỡ nổi và đã đổ sụp xuống, kéo theo cả đống xoong nồi, dao kéo rơi vãi dưới sàn, và thậm chí là càng nhiều bát đĩa bị vỡ tung tóe hơn. Xác gián nằm rải rác khắp nơi, tạo nên khung cảnh chẳng khác gì bãi chiến trường.

Choco nhìn quanh, tay chống hông, thở dài: ``Cô nghĩ là em hơi quá tay rồi đó, Ayame-sama. Không cần phải làm tới mức này đâu!''

Mui cũng không giấu được sự bực bội: ``Tại em mà giờ thì bếp nhà anh ấy tan tành hết rồi! Chúng ta lại phải dọn từ đầu. Chẳng biết bao giờ mới xong được!''

Ayame, vẫn giữ nụ cười tinh nghịch, gãi đầu đáp ``Xin lỗi nha\~{} Nhưng lâu lắm rồi tôi mới được `xung kích' thế này. Cũng vui mà, đúng không?''

Miko, đứng nhìn cảnh tượng đổ nát, khẽ lắc đầu: ``Cũng không thể trách được. Bếp nhà anh ấy bẩn như thế, ai mà chịu nổi chứ ne.''

Mel nhíu mày, tỏ vẻ lo lắng: ``Nhưng giờ phải làm thế nào đây? Cứ để thế này thì không ổn đâu.''

Cuối cùng, cả đội quyết định bắt tay vào sửa chữa hậu quả. Nàng ``Docchi Docchi'' chịu trách nhiệm quét dọn sàn nhà, gom lại từng mảnh vỡ và đống rác rơi vãi. Mio và Choco tiếp tục công việc rửa số bát đĩa còn lại, vừa làm vừa lẩm bẩm đầy chán nản.

Mel và Miko, với sự khéo léo của mình, cùng nhau dựng lại cái chạn treo tường. Dù không thể sửa chữa hoàn hảo, họ tạm thời chống đỡ nó bằng vài viên gạch và đặt lại các đồ dùng vào vị trí cũ.

Nàng Sói đen thở dài sau khi lau xong vũng nước cuối cùng: ``Cuối cùng cũng tạm ổn. Nhưng thật sự, chị không muốn dọn dẹp thêm lần nào nữa đâu.''

Nàng Y tá gật đầu đồng tình: ``Ừ, lần sau cô sẽ để anh ấy tự xử lý bếp của mình. Mình giúp thế này là quá đủ rồi.''

Nàng Quỷ sừng, vẫn giữ vẻ mặt hào hứng, nhìn quanh và nói: ``Không sao đâu. Chúng ta đã làm tốt lắm rồi mà. Coi như luyện tập chút cho vui!''

Nàng Vu nữ nhíu mày, đẩy nhẹ vai nàng Quỷ sừng: ``Luyện tập gì mà làm tanh bành thế này? Nếu anh ấy mà thấy thì chắc ngất mất.''

Nàng Dơi cười khẽ, lắc đầu: ``Thôi, đừng nghĩ nữa. Xong việc là tốt rồi. Dọn dẹp cũng là một phần vui vẻ mà, đúng không?''

Cuối cùng, dù mệt mỏi, mọi người vẫn bật cười khi nhìn lại thành quả của mình. Một căn bếp sạch hơn, sáng sủa hơn, nhưng cũng đầy dấu vết của một trận chiến đầy ``kịch tính''.

Mel khẽ thở phảo, tay xách túi rác chứa đống xác gián: ``Thế là xong. Nhưng mà\dots\ nếu anh ấy hỏi thì ta trả lời thế nào nhỉ?''

Mio chau mày, đưa tay xoa cằm suy nghĩ: ``Nói thật thì chắc anh ấy sẽ nổi cáu mất. Nhưng nói dối thì\dots\ tôi không làm được.''

Miko nghiêng đầu, tỏ vẻ nghiêm túc hiếm thấy: ``Hay là ta cứ nói là vì bếp nhà anh ấy quá bẩn nên ta thay đổi diện mạo nhỉ? Như thế nghe hợp lý hơn đấy, ne.''

Ayame bật cười, giơ hai tay như đầu hàng: ``Thế thì chắc anh ấy sẽ bảo từ giờ khỏi cần giúp nữa. Nhưng mà tôi không thấy phiền đâu, làm thế này vui mà!''

Choco khoanh tay: ``Ayame-sama có vẻ tận hưởng nhất đấy nhỉ. Nhưng liệu Kuon-sama có đồng ý với chuyện này không?''

Mel bật cười, nhẹ nhàng nói: ``Thôi nào, dù sao cũng là giúp anh ấy đón Tết mà. Chỉ mong anh ấy không nhìn thấy cái chạn bị kê bằng gạch là được.'' % Dòng code đầu tiên được viết vào 2025!

Cả đội đứng lặng vài giây, nhìn lại đống hỗn độn dù đã được gói gọn trong những túi rác, rồi đồng loạt bật cười.

Nàng Sói đen khẽ thở dài, nói như tự nhủ: ``Chắc là anh ấy sẽ ngạc nhiên lắm đây. Hy vọng không đến mức `sốc văn hóa'\dots''

Nàng Y tá vỗ nhẹ vai nàng Sói: ``Thôi, làm xong rồi thì xuống tập trung nào. Không biết anh ấy sẽ bảo chúng ta làm gì tiếp theo đây.''

Cả đội rời khỏi phòng bếp, để lại không gian im lìm nhưng sáng sủa hơn hẳn so với lúc họ bắt đầu.

\vspace{-40pt}
\subsubsection*{\phantomsection\label{ep-003.IVf}}
\addcontentsline{toc}{subsubsection}{Miko và Mel}

\begin{center}
    {\color{T6} Nhóm 6 - Miko và Mel}\\Phòng thờ - Phía của Hijaki.
\end{center}

Cùng thời điểm với đội dọn bếp, Miko và Mel đảm nhận nhiệm vụ lau dọn phòng thờ của nhà anh trai Ryujin - Hikawa Hijaki. Khác với những đội khác đang bận rộn với các công việc nặng nhọc, công việc của họ dường như đơn giản và nhàn hạ hơn, khi công việc đơn giản chỉ là lau dọn chén nước và bàn thờ. Về cơ bản, gần giống như những gì mà nàng Vu nữ đã từng làm từ khi cô còn làm việc ở Ngôi đền ảo.

Miko mở cửa phòng thờ, để lộ một không gian trang nghiêm, trầm mặc. Cô khẽ thở dài, lấy ra cây chổi phất bụi quen thuộc được mang từ Ngôi đền ảo của mình. Miko mỉm cười, cảm giác như quay lại những ngày còn làm công việc của một trụ trì: ``Mọi thứ ở đây\dots\ cũng giống như hồi chị còn ở đền vậy, ne. Chỉ cần nhẹ nhàng lau dọn là xong.''

Mel đứng phía sau, tay cầm một chiếc khăn lau, tò mò ngắm nhìn căn phòng. Những tia nắng từ khung cửa sổ nhỏ hắt lên mái tóc vàng óng, khiến cô trông như một thiên thần đang lạc vào thế giới cổ kính.

``Thật thế ư? Chị có thể chờ em quét trần nhà rồi ta có thể dọn bàn thờ được không?'' nàng Ma cà rồng hỏi, trong khi đã chồm người lên một chiếc ghế để quét bụi trên cao.

Miko liếc nhìn, khẽ nhíu mày: ``Cũng được, nhưng chờ bụi bay hết xuống đã. Mà này, em không sợ bụi bay vào cốc chén trên bàn thờ à?''

Mel khựng lại, đôi mắt mở to như nhận ra vấn đề. Cô vội vàng nhảy xuống, nhìn quanh tìm kiếm một vật gì đó để che chắn. Sau một hồi, nàng lấy từ túi ra một chiếc khăn đen nhỏ, phủ lên bộ ấm chén một cách cẩn thận.

``Em nghĩ là được rồi đó!''

Nàng Vu nữ chống nạnh nhìn chiếc khăn bé xíu, nhăn mặt: ``Cái khăn này làm được gì ne? Bé tí thế này không che hết được đâu!''

Tuy nhiên, khi nàng Ma cà rồng đang định giải thích thì bất ngờ thay, từ phía dưới chiếc khăn trùm, một giọng nói trầm cất lên:
``Thả ta ra! Tối quá!''

Cả hai đứng sững, mặt biến sắc. Khi còn chưa kịp định thần, một giọng khác, the thé hơn, vang lên: ``Cái đám này đang làm trò gì thế hả!?''

Miko phản xạ tức thì, lập tức giơ cây chổi phất lên như vũ khí, hét to: ``\textbf{Ai đó!? Ra mau!}'' trong khi Mel run rẩy núp sau lưng nàng trông đền, giọng lí nhí cố trấn an đàn chị nóng máu: ``B-Bình tĩnh nào, Miko-chan\dots\ có thể là\dots\ là gió?''

Cái khăn trùm ngọ nguậy dữ dội, như thể bên dưới có ai đó đang cố gắng vùng vẫy thoát ra. Khi chiếc khăn rơi xuống, để lộ\dots\ chẳng có gì cả, ngoài bộ ấm chén vẫn nằm nguyên vị trí ban đầu.

Nàng Dơi há hốc miệng, đôi tay run rẩy chỉ vể phía bàn thờ: ``M-Ma ư!?'' Nàng Vu nữ quay lại, trố mắt nhìn cô, tỏ vẻ chút khó tin: ``Em là ma cà rồng mà còn sợ ma sao!?''

Cô nàng xấu hổ đỏ mặt, nhưng cũng không dám tiến lại gần bàn thờ nữa. Cả hai đứng nhìn nhau trong vài giây, trước khi Miko thở dài, buông chổi xuống: ``Chắc là tưởng tượng thôi\dots\ tiếp tục nào. Dọn nhanh cho xong rồi ra ngoài, ne!''

Cô khẽ gật đầu, nhưng vẫn không quên giữ khoảng cách an toàn với bàn thờ.

Họ vừa tiếp tục dọn dẹp bàn thờ, vừa tự nhẩm với nhau rằng có lẽ đó chỉ là tiếng gió rit thổi qua, không hơn không kém.

Thế nhưng, ngay khi họ đinh ninh như vậy, bất chợt từ bàn thờ vang lên tiếng quát lớn, dứt khoát đến mức cả hai giật mình:

\begin{dialogue}
    \speak{Ma-ông} Giới trẻ ngày nay đúng là hư thân mất nết!
    \speak{Ma-bà} Cái đám mặc quần áo lòe loẹt này là ai đây!?
    \speak{Ma-ông} Lại còn con kia nữa, mặc quần áo hở hang quá!
\end{dialogue}

Mel ngay lập tức cảm thấy như mình đang bị ``nhắm đến.'' Cô lí nhí thanh minh, mặt đỏ bừng: ``Này, tôi chỉ có mỗi bộ đồ này thôi mà!''

Chưa dừng lại ở đó, ``ma ông'' quay sang chỉ trích luôn nàng Vu nữ: ``Lại còn cái đứa nói `ne, ne' nữa! Ăn nói bình thường đi giùm tao cái!''

Miko cảm thấy bị xúc phạm nặng nề. Cô đập mạnh cây chổi phất xuống sàn, hét lớn: ``\textbf{Cách ăn nói của tôi thì liên quan gì đến ông hả, ne!?}''

Tình hình trở nên hỗn loạn. Cả bốn người - hai ``ma'' và hai người - bắt đầu cãi cọ om sòm. ``Ma bà'' chỉ trích phong cách ăn mặc của Mel, ``ma ông'' thì chê bai cách ăn nói của Miko, còn hai cô nàng thì ra sức phản bác, không chịu nhường nhịn chút nào.

Cuối cùng, nàng Vu nữ đã không thể chịu đựng thêm. Cô giương cao cây gậy phẩy - vốn cũng là cây gậy trừ tà - với ánh mắt đầy tức giận: ``Đã thế, tôi cho ông bà bay màu luôn, ne!''

Nàng Ma cà rồng hốt hoảng chạy đến ngăn cản: ``Miko-chan, bình tĩnh lại nào! Đừng làm quá như vậy chứ!''

Nhưng đã quá muộn. Miko-chi bắt đầu thực hiện một điệu múa trừ tà đầy uy lực. Những động tác mạnh mẽ cùng tiếng hét đầy quyết tâm của cô khiến không gian trở nên căng thẳng.

\begin{dialogue}
    \speak{Ma-bà} Nó đang thực hiện điệu múa gì thế?
    \speak{Ma-ông} Có chó mới biết được. Giới trẻ thế này là hỏng hết rồi-
\end{dialogue}

Miko hét lớn: ``\textbf{Trừ tà!!}''

Vừa dứt lời, căn phòng bỗng chốc trở nên im bặt. Không còn bất kỳ tiếng nói nào, ngoài tiếng thở hổn hển của Miko-chi. ``Chị\dots\ đuổi chúng\dots\ đi rồi\dots\ ne\dots''

Mel nhìn về phía bàn thờ, tỏ vẻ lo âu: ``Miko-chan, chị không sợ họ sẽ đến trả thù chứ?''

Nàng Vu nữ đứng dậy, nhìn cô nàng, nói với giọng chắc nịch: ``Họ bị biến khỏi trần gian rồi, chắc không sao đâu.'' Dù vậy, lời nói của cô không khiến nàng Dơi nhẹ nhõm hơn: ``Em cảm thấy ta vừa làm gì đó tội lỗi ấy\dots''

Miko phẩy tay, cầm lấy chiếc khăn và chiếc gậy phẩy bụi của mình: ``Kệ đi ne. Ta dọn dẹp tiếp nào.'' Nói rồi, tiền bối của cô tiếp tục lau dọn bàn thờ một cách bình thản.

Về phía Mel, cô chỉ có thể thở dài, không còn cách nào khác ngoài tiếp tục lau dọn bộ bàn ghế và tủ kệ trong phòng.

\ct

Khi hai cô nàng vừa hoàn thành công việc và chuẩn bị xuống tầng ba để hội quân, bất chợt có những tiếng hét thất thanh vọng lên từ phòng bếp bên cạnh. Miko và Mel giật mình nhìn nhau, rồi lập tức chạy sang kiểm tra.

\begin{dialogue}
    \speak{Miko} Có chuyện gì vậy, ne!?
    \speak{Mel} Ở đây có gì mà la toáng lên thế-
\end{dialogue}

Cảnh tượng trước mắt khiến cả hai bàng hoàng. Mio, Choco và Ayame đang trong tư thế ``tử chiến'' với một đàn gián hoành hành dưới sàn. Nàng Vu nữ ngay lập tức hét lên hoảng sợ, ``song phi'' lên tủ lạnh cùng với nàng Sói đen.

``C-Chuyện gì đang xảy ra vậy ne!?'' nàng Vu nữ hoảng loạn vung cây gậy trừ tà về phía chúng, trong khi vẫn đang cố thủ trên nóc tủ lạnh.

``Chúng tôi đang dọn rãnh nước ở bếp nhà anh ấy\dots'' Mio cất lời, để rồi Choco tiếp tục: ``\dots thì bọn cô gặp phải lũ gián đang ở dưới kia kìa!''

Trong khi đó, Mel cuống cuồng đi tìm cây chổi để tiêu diệt lũ gián, nhưng tay chân cô lóng ngóng đến mức đụng đổ cả một chồng bát đĩa đã được rửa sạch trước đó.

Những gì diễn ra sau đó là một màn hỗn loạn không thể tả.

\vspace{-40pt}
\subsubsection*{\phantomsection\label{ep-003.IVg}}
\addcontentsline{toc}{subsubsection}{Kaien, Ryuuhou, Aki}

\begin{center}
    {\color{T7} Nhóm 7 - Kaien, Ryuuhou, Aki}\\Phòng bếp - Phía của Hijaki.
\end{center}

Vì tất cả mọi người đang bận rộn lau dọn các tầng khác, nên Aki cùng với mẹ và em gái của Kuon đảm nhiệm công việc giặt quần áo ở tầng bốn. Kaien, tay xách một chậu quần áo cao đến gần đỉnh đầu, dẫn đầu cả nhóm: ``Ba cô cháu mình xuống tầng 4 giặt quần áo nào, Aki-chan!''

Nhìn chồng quần áo chất cao như núi, nàng Dị giới không khỏi thắc mắc: ``Sao nhà cô có nhiều quần áo phải giặt thế ạ?''

Người mẹ của nhà Hikawa khẽ bật cười, rồi nói với giọng tếu táo nhưng chua xót: ``Bố con nhà này (ám chỉ Ryujin và Kuon) có bao giờ chịu giặt đâu. Chồng với chả con, lười thế đấy!''

Aki bật cười nhẹ, nhưng vẫn tỏ vẻ tò mò: ``Nhà cô không có máy giặt sao ạ?''

Kaien đáp. nói với sự tự tin: ``Có chứ, nhưng máy giặt làm sao sạch bằng tay được! Với lại, giặt tay cũng là cách để kiểm tra quần áo kỹ hơn.''

Cô quay sang Ryuuhou và hỏi: ``Cô ấy nói thế, đúng không Ryuuhou-chan?'' Cô em gái của Kuon gãi đầu ngượng nghịu: ``Mẹ em nói thì em biết vậy thôi\dots\ Với cả ấy\dots''

Cô ấy ghé sát tai vào nàng Dị giới, nói một cách khẽ khàng: ``\hl{Bà ấy cũng là một người cuồng việc nhà đấy.}'' Aki ngơ ngác nhìn cô không hiểu ý, nhưng cũng không muốn hỏi gì thêm.

Ba người vác ba chậu quần áo xuống tầng, đặt trước cửa phòng tắm. Kaien bắt đầu tìm thêm một cái chậu để hứng nước, nhưng lục tìm mãi không thấy.

``{\color{red} Chắc phải để quần áo dưới sàn rồi giặt nhỉ\dots}'' bà lẩm bẩm. Định bảo cô con gái lên nhà lấy thêm chậu, thì Aki lên tiếng, với sự tự tin kỳ lạ: ``Cô để cháu.''

Cô nàng Dị giới liền nằm ngửa xuống sàn, bắt đầu rặn một cách nghiêm túc. Hai mẹ con nhà Hikawa nhìn cô nàng đang ú ớ như đang đi đẻ, rồi nhìn nhau khó hiểu.

``Cháu\dots\ đang làm gì vậy?'' người mẹ hỏi với sự khó hiểu. Nàng Tóc bay đáp: ``Cháu tạo ra một cái chậu!'' rồi tiếp tục rặn, mặt cô đỏ bừng.

Ryuuhou nhìn chị gái với sự ngờ vực: ``Tạo ra một cái chậu á? Làm gì có chuyện đó! Aki-onee-chan, đừng có đùa nữa!''

Đến cả Kaien cũng lắc đầu. Bà thở dài: ``Thôi, vớ vẩn. Ngồi dậy đi, để cô tự đi tìm chậu trên nhà-''

Chưa để bà ấy nói xong, Aki liền cắt lời, nở một nụ cười rạng rỡ: ``Xong rồi đây!''

Cả Kaien và Ryuuhou tròn mắt kinh ngạc khi thấy một cái chậu nhựa vàng cỡ bốn gang tay xuất hiện ngay trước mặt, dường như vừa ``chui ra'' từ đâu đó\footnote{Cơ chế này đã được thể hiện ở các tập {\flaregothic ÉPISODE 10} và {\flaregothic ÉPISODE 16} ở quyển số Một.}. Họ cùng nhau đồng thanh: ``Thế cũng được cơ à!?''

Aki nở một nụ cười rạng rỡ, vui vẻ đáp: ``Cháu nói thật mà! Cô thử chạm vào đi.''

Người đàn bà của nhà Hikawa nửa tin nửa ngờ, đưa tay chạm vào chiếc chậu. Tiếng nhựa cồm cộp vang lên dưới ngón tay bà, hoàn toàn giống như những chiếc chậu nhựa thông thường. Vẫn còn hơi bất ngờ, bà lẩm bẩm: ``Cái này\dots\ đúng là kỳ lạ thật\dots''

Ryuuhou, sau khi nhìn kỹ chiếc chậu, cầm lên thử và cảm nhận trọng lượng. Cô gật gù: ``Chuẩn chậu rồi mẹ. Khá nặng đấy!''

Dù vẫn thấy khó hiểu, Kaien quyết định tạm tin Aki. Bà dùng chiếc chậu mới tạo ra để bỏ quần áo, trong khi chiếc chậu lớn nhất lấy trên nhà được hứng nước.

Nàng Dị giới đứng dậy, phủi bụi trên quần áo: ``Xong! Cháu bảo mà, cháu giúp được!'' Kaien mỉm cười nhẹ: ``Giỏi lắm, Aki-chan. Nhưng lần sau nhớ giải thích trước đã nhé, kẻo người ta lại tưởng cháu bị gì.''

Trong khi đó, cô em gái của Kuon vẫn còn đang tỏ vẻ ngơ ngác, không hiểu chuyện gì vừa xảy ra: ``Mẹ ơi, con vẫn không hiểu sao chị ấy làm được\dots''

Cả ba bật cười nhẹ nhàng rồi tiếp tục công việc. Aki, với sự hài hước vốn có, đã khiến buổi giặt quần áo trở nên vui vẻ hơn rất nhiều.

Trong lúc chờ nước hứng đầy, mẹ của Kuon cẩn thận lấy gói bột giặt mà bà đã mang từ phòng tắm tầng trên xuống, pha đều vào chậu nước. Bà quay sang hỏi Aki: ``Cháu có biết giặt quần áo không?''

Cô nàng Dị giới nghiêng đầu, dù cô biết rõ khái niệm giặt quần áo là gì: ``Cháu ạ? Cháu chưa bao giờ giặt tay cả\dots\ Cháu chỉ toàn dùng máy giặt thôi ạ.''

Người mẹ mỉm cười: ``Vậy để cô dạy cho nhé.'' Cô em gái cũng hồ hởi: ``Để em phụ chị nữa!''

Bà lấy một chiếc áo thun trắng, đặt vào chậu nước đã hòa bột giặt, vừa làm vừa giảng giải: ``Nhìn cô làm này. Đầu tiên, thả quần áo vào nước, ngâm một lúc để bột giặt thấm vào sợi vải\dots''

``Vâng ạ!''

``Sau đó, dùng hai tay vò mạnh ở những chỗ dễ bẩn nhất như cổ, nách và tay áo. Đây là những nơi cần chú ý.''

Tới đây, nàng Tóc bay tỏ ra thắc mắc: ``Dạ\dots\ Nhưng tại sao mấy chỗ đó lại bẩn hơn ạ?'' Ryuuhou nhanh nhảu giải thích: ``Thì do mồ hôi và bụi bẩn tích tụ nhiều ở đó chứ sao nữa!''

Aki gật gù, mỉm cười: ``Ồ\dots\ Ra là vậy!'' Kaien tiếp tục từ tốn giải thích: ``Rồi ta lặp lại như thế vài lần, sau đó vắt sạch nước xà phòng\dots\ Thế là xong!''

Bà đặt chiếc áo vừa giặt xong vào một cái chậu khô khác rồi quay sang động viên cô: ``Cháu thử làm xem nào!'' Cô hồ hởi đáp: ``Vâng, để cháu thử!''

Bà lấy một cái chậu khác, đổ nước và bột giặt vào, sau đó đưa cho Aki. Với từng món quần áo, Kaien và Ryuuhou chỉ dẫn cô nàng Dị giới khá cặn kẽ. Aki chăm chú làm theo, từng bước một, và ``trộm vía em ấy học hỏi nhanh lắm'' (lời Kuon kể lại), nên không tốn quá nhiều thời gian để dạy cô.

``Đấy, cháu còn chăm hơn cả thằng Kuon nhà cô. Nó lười đến mức mà cô không chịu nổi,'' Kaien vui vẻ nói, không quên khoáy vào nỗi hổ thẹn của cậu Tổng.

Đáp lại, nàng Dị giới chỉ cười một cách ngại ngùng, thầm nghĩ: ``\textit{Anh ấy dường như lúc nào cũng bận mà cô ơi\dots}''

\ct

Nhờ sự phối hợp nhịp nhàng, chỉ trong thời gian ngắn, ba người đã giặt gần hết số quần áo. Mọi chuyện không có gì quá khó khăn, cho đến khi\dots

*ROẸT* Một âm thanh khẽ vang lên khiến cả ba người dừng tay. Aki bối rối giơ lên một chiếc áo dài tay màu xanh lá với vết rách rõ rệt ở đằng trước.

Nàng Dị giới hoảng loạn: ``T-Thôi chết! Chau lỡ làm rách áo của Ryujin-san rồi!''

``Để cô xem nào,'' Kaien chậm rãi lấy chiếc áo mà cô đã vô tình làm rách. Bà cầm chiếc áo lên, xem xét kỹ lưỡng rồi mỉm cười nhẹ nhàng: ``Không sao đâu, để tí nữa cô khâu lại là được.''

Nàng Tóc bay ngạc nhiên: ``Khâu lại được ạ? Áo rách thế này liệu có ổn không?''

Kaien cười xòa: ``Có gì mà không làm được? Cứ yên tâm, lát nữa giặt xong rồi đem đi phơi, cô sẽ khâu cho.'' Đến cả cô em gái cũng phụ họa: ``Chị cứ yên tâm đi. Bà ấy làm việc nhà gì cũng thạo mà.''

Nghe vậy, Aki cúi đầu cảm kích: ``Cháu cảm ơn cô ạ!''

Cả ba tiếp tục công việc, và chẳng mấy chốc, tất cả quần áo đã được giặt sạch sẽ. Người mẹ của nhà Hikawa đứng dậy, phủi tay: ``Rồi, mang đống này đi phơi nào!''

Nàng Dị giới hăng hái bê chậu quần áo của mình - chiếc chậu đỏ do cô tự tạo ra - lên tầng sáu, theo sau là Kaien và Ryuuhou.

Việc phơi quần áo không quá khó khăn. Aki nhanh chóng quen với cách treo quần áo lên sạp, những chiếc còn lại được sử dụng mắc phơi và treo lên dây chung với nhà Hijaki.

``Vậy là xong! Ngày mai quần áo sẽ được phơi khô thôi,'' Ryuuhou nói trong sự mừng rỡ.

``Nhờ cháu mà chúng ta đã giặt xong đống quần áo nhanh hơn hẳn đấy,'' Kaien bộc bạch, ``chứ không á, có khi bây giờ cô mới chỉ xong một phần ba núi quần áo thôi!''

Cả ba cô gái đều bật cười với lời bình của bà. Khi mọi việc hoàn thành, cả ba cùng nhau xuống tầng để chuẩn bị cho nhiệm vụ tiếp theo, không ngó ngàng gì tới căn bếp đã có sự thay đổi đôi chút.

\vspace{-40pt}
\subsubsection*{\phantomsection\label{ep-003.IVh}}
\addcontentsline{toc}{subsubsection}{Fubuki, Rushia, Aqua}

\begin{center}
    {\color{T8} Nhóm 8 - Fubuki, Rushia, Aqua}\\Tầng 3.
\end{center}

Fubuki, Aqua, và Rushia tập trung tại tầng ba nhà của Kuon để bắt đầu công việc dọn dẹp. Gọi là dọn dẹp, nhưng thực ra ở đây cũng chẳng có mấy việc để cho họ làm, do nó đã được dọn dẹp từ rất lâu.

Tầng này vốn không được sử dụng trong nhiều năm, ngoại trừ lần gần nhất, đó chính là cặp đôi Korone - Okayu, những người vừa sơn lại bức tường của căn phòng màu trắng này thành một ``kiệt tác nghệ thuật.''

Dẫu vậy, bụi bặm và mạng nhện cùng với sự xuống cấp đôi chút của thời gian vẫn len lỏi khắp nơi, tạo nên một cảm giác tuy có trong sáng nhưng lại có phần ảm đạm, thậm chí còn lạnh lẽo.

Ban đầu, cả ba nghĩ rằng việc dọn dẹp sẽ rất nhẹ nhàng, chỉ cần quét bụi, lau sàn, và xử lý vài mạng nhện. Nhưng hiện thực lại không hề dê dàng đến vậy.

Nàng Cáo trắng nằm ườn ra trên một chiếc ghế bành cũ kỹ phủ vải bạt, đôi tai cáo của cô cụp xuống như thể thể hiện sự buồn chán cực độ, giống như hồi mà cô và Matsuri đã làm hồi sáng nay: ``Chán quá\~{} Chẳng có gì để làm cả\dots''

Aqua và Rushia đứng sững lại, chưng hửng nhìn thái độ lười biếng của tiền bối họ.

\begin{dialogue}
    \speak{Rushia} Fubuki-senpai, đừng có lười thế\dots\ Ta vẫn có việc để làm mà-nanodesu!
    \speak{Aqua} Nhưng mà, chị ấy nói cũng đúng\dots\ Chị thấy tầng này hầu như chẳng có gì nhiều để cho ta làm cả\dots
\end{dialogue}

Nàng Pháp sư ôm đầu, đôi mắt liếc qua nàng Hầu gái đầy bất mãn: ``Sao lại không? Lau lại sàn, quét trần, rồi còn dọn nhà vệ sinh hai bên nữa!''

Aqua nghiêng đầu: ``Thế sao? Vậy để chị đi quét trần với quét nhà nha. Mấy cái này hầu gái như chị làm suốt mà.''

Nàng Cáo trắn nhướng mày, nở nụ cười nham hiểm như muốn bắt bẻ: ``À-rế\~{}? Chị tưởng em nói mình không phải hầu gái cơ mà nhỉ\~{}?'' Nàng Cô đồng cũng gật gù đồng tình: ``Em cũng thấy nghi lắm đó-nanodesu\~{}''

Nàng ``Hành tím'' đỏ mặt, vung cây chổi cầm sẵn trong tay chỉ thẳng vào cả hai: ``Em \emph{là} hầu gái, hiểu chưa? Với cả đừng có đứng đó mà châm chọc nữa, đi phụ em dọn lẹ đi chứ!''

Fubuki thở dài đầy uể oải, cuối cùng cũng chịu đứng dậy lấy cây chổi quét sàn và phụ \sout{B}Aqua: ``Rồi, rồi\dots\ Chị phụ một tay đây. Nhưng đừng hối chị làm nhanh nhé\~{}''

Trong khi đó, Rushia dường như lại có cách tiếp cận công việc theo kiểu riêng của mình. Đứng giữa phòng, cô nhắm mắt, bắt đầu thì thầm những câu chú ngữ phức tạp, đôi tay giơ lên cao. Không khí xung quanh bỗng trở nên lạnh hơn, khiến cả Aqua và Fubuki ngừng tay quay lại nhìn.

Từ dưới nền gạch, những bàn tay xương xẩu đen ngòm bắt đầu xuất hiện, bám vào mặt đất và từ từ trồi lên. Chỉ trong nháy mắt, một binh đoàn những bộ xương người với đôi mắt rực sáng màu xanh lục tập hợp phía sau Rushia, sẵn sàng chờ lệnh.

``Cái gì thế này, Rushia-chan?'' nàng Cáo trắng tỏ vẻ tò mò với những bộ xương mà nàng Pháp sư vừa triệu lên.

Nàng Pháp sư trả lời: ``Fandead của em đấy! Các thuộc hạ của ta, lên nào!'' Aqua nhìn vậy, cô giật mình: ``S-Sợ thế!!''

Cô đồng tiếp tục phẩy tay, chỉ đạo đám Fandead và phân chia nhiệm vụ: ``Nhóm này dọn nhà vệ sinh, nhóm kia lau sàn, nhóm còn lại quét trần. Mau lên, nhanh gọn vào!''

Nhìn thấy đám thuộc hạ cần mẫn thực hiện nhiệm vụ, nàng Cáo trắng nhanh chóng trở về trạng thái\dots\ lười biếng: ``Có họ làm rồi, vậy chị còn phải làm gì nữa\~{}?'' Trong khi đó, nàng Hầu gái thì vẫn đang run lẩy bẩy: ``C-Chúng\dots\ Chúng biết cử động kìa!!''

Cô nàng Pháp sư phải quay sang trấn an cô: ``Chúng không làm hại chị đâu. Với lại, mấy đứa này chỉ tồn tại trong một khoảng thời gian ngắn thôi. Xong việc là chúng lại quay về mặt đất.''

Cô không quên quay sang người tiền bối đang nằm thảnh thơi trên ghế bành và nói vẻ nghiêm túc: ``Điều đó có nghĩa là ta vẫn phải dọn đó, Fubuki-senpai.''

Đáp lại Rushia là hai tiếng thở dài - một từ Aqua đang cố trấn tĩnh, một từ Fubuki đầy chán nản:

\begin{dialogue}
    \speak{Aqua} R-Rồi, chị làm ngay đây\dots
    \speak{Fubuki} Rồi\~{}\dots
\end{dialogue}

\ct

Aqua tiếp tục làm công việc của mình, lau chùi từng góc nhỏ của căn phòng. Nàng Cáo trắng, với thái độ miễn cưỡng, cuối cùng cũng chịu làm việc và phụ Aqua lau các đồ nội thất trong phòng. Trong khi đó, nàng Cô đồng tóc xanh vẫn đứng ra điều khiển các thuộc hạ xương xẩu của mình dọn dẹp một cách bài bản.

Mọi thứ tưởng chừng như êm đềm, nhưng sự yên tĩnh ấy không kéo dài lâu\dots

\textbf{*XOẢNG!*}

Tiếng kính vỡ bất ngờ vang lên. Nàng Hầu gái vừa lau dọn những chiếc cốc trên kệ thì bất cẩn trượt tay làm rơi một chiếc xuống đất. Cô hoảng hốt quỳ xuống thu nhặt những mảnh vỡ: ``Thôi chết, mình làm vỡ cốc rồi!''

Tiếng động cũng làm cho hai người kia giật mình, vội chạy đến hỏi han. Nàng Pháp sư lên tiếng: ``Chị có sao không!?''

Aqua đáp, vẫn vội vàng gom lại những vết nứt vỡ: ``Chị ổn\dots\ Nhưng cái cốc này bị vỡ có sao không nhỉ?''

Fubuki đang lau một góc khác bèn liếc mắt nhìn qua rồi nói với giọng tỉnh bơ: ``Không sao đâu\~{} Dù gì nó có phải cốc nhà anh ấy đâu.''

``\textit{Ý chị là sao chứ\dots}'' Aqua thầm nghĩ, rồi tiếp tục lau dọn phòng. Thế nhưng, ngay khi cô vừa mới quay mặt đi, một tiếng \textbf{*XOẢNG*} nữa vang lên, lần này tới từ phía mấy bộ xương. Hai ``nhân viên'' bộ xương của Rushia trong lúc khiêng ghế thì bị trượt chân do sàn quá trơn - hậu quả từ việc lau sàn chưa khô của Fubuki trước đó - khiến chúng làm vỡ cả cái bàn thủy tinh gần đó.

Aqua hoảng hốt: ``Sao phòng này toàn đồ dễ vỡ thế!?'' Nàng Cáo trắng giả ngơ, tay vẫn miệt mài lau sàn: ``Chị hơm có biế đâu à nha\~{}''

Nhưng đó chỉ mới là khởi đầu. Sàn nhà trơn bóng vô tình trở thành "cạm bẫy" cho các bộ xương khác. Hết đứa này trượt ngã làm rơi cốc trên kệ, đứa khác va vào cửa kính nhà vệ sinh. Thậm chí, một bộ xương còn ``phá cách'' hơn khi\dots\ đâm vỡ cả cửa sổ.

Lúc này, mặt nàng Hầu gái tái mét, còn nàng Pháp sư thì đổ mồ hôi hột.

\begin{dialogue}
    \speak{Aqua} N-Nguy to rồi!
    \speak{Rushia} Phòng này có nhiều kính quá!
    \speak{Fubuki} Cũng chẳng trách được, phòng này đã lâu rồi chưa sử dụng mà\dots
\end{dialogue}

Mặc cho Fubuki đáp lại bằng giọng hờ hững như thể cô không can tâm gì tới chuyện này, cả Aqua và Rushia đều méo mặt khi càng lúc càng có nhiều tiếng đổ vỡ hơn.

``Kính để lâu đâu có dễ vỡ thế này được!'' nàng ``Hành tím'' thốt lên, theo sau là nàng Pháp sư cũng tỏ vẻ băn khoăn: ``Ừ nhỉ\dots\ Chuyện này có gì đó sai sai\dots''

Đúng lúc đó, từ tầng hai vang lên tiếng hét thất thanh. Cả ba tá hỏa chạy xuống. Dưới đó là hai nhân viên của cửa hàng viễn thông đang hoảng sợ, và nguyên nhân là do\dots\ hai bộ xương của Rushia đang cần mẫn vận chuyển đồ đạc từ chỗ này sang chỗ khác.

Nàng Hầu gái trố mắt ngỡ ngàng: ``Sao bộ xương của em lại ở dưới đó!?'' trong khi chủ nhân của hai bộ xương kia cố gắng điều khiển chúng: ``Fandead, dừng lại! Các ngươi không phải dọn ở đây đâu!''

Hai nhân viên nữ mặc đồng phục cửa hàng viễn thông la hét thảm thiết:

{\color{red}
    \begin{dialogue}
        \speak{Emp-A} Cái gì thế này!? Ai đó đuổi chúng đi hộ tôi với!!
        \speak{Emp-B} Éc!! Đừng có lại gần đây!!
    \end{dialogue}}

Rushia lập tức niệm chú. Không khí căn phòng đột ngột chuyển đỏ rực. Những dòng ký tự ma thuật hiện lên, và cô hét lớn: ``\textbf{Xuống mồ!}''

Ngay lập tức, hai bộ xương chui tọt xuống đất, để lại hai lỗ nhỏ trên sàn gạch. Hai nhân viên vì quá sợ hãi đã ngất lịm ngay tại chỗ.

Fubuki ngơ ngác nhìn hai nữ nhân viên bất tỉnh sau sự cố đó, đổ mồ hôi hột: ``Họ đúng là nhát thật đấy\dots''

Aqua thì hoảng loạn hơn: ``T-T-Ta xử lý họ như thế nào đây!?'', còn Rushia thì trầm ngâm một lúc lâu, rồi đáp với vẻ lưỡng lự: ``Chắc là\dots\ để họ tự tỉnh thôi nhỉ?''

``C-Chị nghĩ như vậy là tốt nhất\dots\ Chứ động vào người họ có khi còn nguy hiểm hơn\dots'' nàng Hầu giá cất giọng đồng tình, dù còn run rẩy vì những gì họ đã chứng kiến.

``Chị cũng thế\~{}'' Fubuki hùa theo. Thế là, họ quyết định quay lại tầng ba, mặc kệ sự hỗn loạn ở tầng dưới. Trở lại phòng cũ, Aqua và Rushia cố gắng hoàn thành công việc, trong khi Fubuki lại bắt đầu nằm ườn ra như trước.

Gọi là ``hoàn thành công việc'' chứ thực chất họ đang cố gắng phi tang hết hiện trường bừa bộn mà đám thuộc hạ của Rushia đã làm, trong khi nàng Cáo trắng thì vẫn bình chân như vại, ngáp ngắn ngáp dài nhìn hai người họ cật lực làm.

May mắn thay, Shion - người đã hoàn thành công việc từ rất sớm - xuất hiện. Cô nhìn nàng Pháp sư và nàng Hấu gái hấp tấp lau dọn, bèn hỏi: ``Có chuyện gì mà ba người ồn ào vậy?''

Dù đang cố tỏ ra nghiêm nghị, nhưng khi thấy khung cảnh bừa bộn, cô không nhịn được mà bật cười: ``Có chuyện gì thế? Ba người đang giúp gia đình anh ấy hay là đang rủ nhau đập phá vậy!?''

Aqua lên tiếng: ``C-Cũng tại sàn trơn chứ bộ!'' Rushia chi tay về phía cô nàng đang nằm thảnh thơi trên ghế bành: ``Tại Fubuki-senpai lau nhà chưa khô nên mới vậy đó! Với cả em cũng đâu có muốn làm vỡ đồ cơ chứ!''

Nàng Cáo trắng bật dậy, lên tiếng phản ứng dữ dội: ``N-Này! Chính Aku-tan bảo bọn chị phải lau dọn nha! Chị đã cố hết sức rồi đấy!''

Nàng Pháp sư không chịu được, cũng phải lên tiếng: ``Có ai lau nhà trơn trượt như chị không!? Lau như chị thì thà bọn em không nhờ còn hơn!'' Nàng Hầu gái cũng góp vào: ``Chị cứ lười chảy thây thế, ai mà chịu cho nổi!?''

Trong lúc ba cô nàng còn đang tranh cãi kịch liệt, nàng Phù thủy cẩn thận bước vào, tránh những mảnh kính vỡ dưới sàn. Cô tiến tới, nói với ba người: ``Thôi, thôi được rồi, mọi người đừng om sòm lên nữa. Có vẻ như các chị cần đến phép thuật của một thiên tài như em đây.''

Ngay lập tức, cả ba cô gái dừng lại, nhìn về phía cô với những ánh mắt ngơ ngác. Cô mỉm cười tự tin, xắn tay áo như thể sắp thực hiện một màn trình diễn hoành tráng.

Chưa kịp để ai thốt lên một lời nào, Shion lùi lại vài bước, tạo ra một vòng tròn phép thuật của mình. Cô hít sâu, nhắm mắt tập trung. Bầu không khí trong phòng đột nhiên trở nên khác lạ, như thể một luồng gió vô hình đang cuộn quanh.

Cô niệm một câu thần chú, nhẩm đọc một cách liến thoắng, rồi hét lớn: ``\textbf{Khôi phục!}''

Từng mảnh kính vỡ từ dưới sàn bỗng chốc bay lên không trung, xoay vòng như một điệu múa. Chúng hợp lại, khôi phục hình dáng cũ của những chiếc cốc, cửa sổ, và cả bàn thủy tinh. Chỉ trong chớp mắt, mọi thứ đã trở lại như chưa từng bị hỏng.

Aqua há hốc miệng: ``Oa\dots\ Shion-chan làm được thật này!'' Fubuki thì vỗ tay nhẹ: ``Không tệ đâu, Shion-chan à\~{}'', thậm chí cả nàng Pháp sư cũng phải công nhận tài năng phép thuật của cô: ``Đúng là Shion-senpai có khác.''

Được cả ba cô gái khen một cách nức nở đầy hiếm thấy, Shion vênh mặt lên đầy tự hào: ``Thấy chưa? Đừng bao giờ nghi ngờ kỹ năng phép thuật của Shion đại tài này.''

Về cơ bản, căn phòng đã được hoàn tác về trạng thái cũ, với cốc chén còn nguyên vẹn, cửa sổ không còn lấy một nứt vỡ, bức tranh mà cặp đôi OkaKoro vẫn còn đó\dots

\dots Ngoại trừ việc, tất cả bụi bẩn vẫn còn nguyên xi ở đó, như thể cả ba người vẫn chưa hề đụng tay đụng chân vào.

``A-A-nô\dots'' Aqua lên tiếng, ``Bà có thể giúp bọn này nốt phần việc còn lại không? Kuon-senpai về mà thấy chúng ta chưa làm gì thì sẽ khó chịu đấy.''

Nàng Phù thủy liếc qua đống việc còn dang dở, khẽ nhăn mặt: ``Hầy,  đúng là giúp người khác cũng mệt thật. Thôi được rồi, em làm nốt cho.''

Cô lại tạo cho mình một vòng tròn ma pháp khác, niệm một câu thần chú khác nữa. Sau khi đọc xong, cô búng tay một cái (giống như cách mà cô làm với phòng của Hijaki), lập tức từng mảng bụi bẩn trong phòng tự động biến mất, đồ đạc được sắp xếp ngay ngắn về vị trí ban đầu. Khi hoàn thành, căn phòng trở nên sáng bóng và gọn gàng như mới.

Sau khi mọi việc đâu vào đó, cả ba đều nằm bệt xuống sàn, thở phào nhẹ nhõm. Nàng Phù thủy cũng ngồi phịch xuống vì đã sử dụng kha khá ma pháp trong ngày hôm nay.

Aqua than vãn: ``Mệt chết mất\dots\ Dọn dẹp đúng là không đơn giản chút nào\dots'' Nàng Cáo trắng thì tỏ vẻ nhởn nhơ: ``Ừ, nhưng vui mà\~{}''

Rushia thì nhăn mặt, phản ứng lại ngay với tiền bối cô: ``Chị gần như chẳng giúp bọn em một tí nào cả!'' còn Shion thì quệt mồ hôi đổ trên trán: ``Vậy là xong rồi đấy! Mọi người nợ tôi một bữa ăn đấy.''

Cả ba cô gái đều bật cười với yêu cầu của cô. Không phải là những tiếng cười khinh bỉ, mà là những nụ cười của sự chân thành. Cả bốn người bắt đầu bàn nhau xem họ sẽ ăn gì tại thành phố Kawachi này, trong lúc chờ những người còn lại tụ tập ngay tại đây. Tiếng cười nói rôm rả vang lên, chủ yếu xoay quanh các món đặc sản địa phương mà Kuon gợi ý, nhưng chẳng ai đi đến thống nhất được vì mỗi người lại thích một món khác nhau.

\vspace{-40pt}
\subsubsection*{\phantomsection\label{ep-003.IVi}}
\addcontentsline{toc}{subsubsection}{Tất cả mọi người}

\begin{center}
    {\color{T9} Master - Tất cả mọi người}\\Tầng 3.
\end{center}

Giữa lúc cuộc tranh luận về đồ ăn vẫn chưa ngã ngũ, từng bóng dáng thân quen lần lượt bước vào căn phòng tầng ba. Căn phòng giờ đây khoác lên mình một diện mạo hoàn toàn mới: bàn ghế được lau sạch bóng, bụi bặm biến mất, mạng nhện không còn dấu vết, trần nhà sáng sủa, sàn nhà không có một vết bẩn.

Đặc biệt nhất là những bức tường được sơn xen kẽ hai màu tím và nâu nhạt bởi sự phối hợp nhịp nhàng của Korone và Okayu. Trông như một khung cảnh ấm áp, vừa trang nhã lại có chút hoài niệm, đủ để người ta ngắm nghía mãi không thôi.

Dù vậy, căn phòng vẫn phảng phất mùi sơn mới chưa tan hết. Mọi người phải mở toang các cửa sổ để đón gió, hi vọng mùi khó chịu sớm bay đi.

``Bọn anh mua đồ về rồi đây!'' Kuon thốt lên, tay xách vài túi đồ mua về từ siêu thị. Kế sau đó là Sora và Ryujin, những người cũng xách những túi đồ nặng không kém gì cậu.

\begin{dialogue}
    \speak{Sora} Mọi người vất vả rồi!
    \speak{Ryujin} Các cháu dọn nhà đến đâu rồi? Báo cáo cho chú nghe xem nào.
\end{dialogue}

Lập tức, tất cả bảy đội còn lại bắt đầu trình bày công việc và những gì họ bắt gặp trong quá trình làm việc. Bắt đầu với đội đầu tiên.

Noel hồ hởi tiến tới: ``Phòng của bố con anh đã được dọn hoàn tất! Sạch sẽ không còn một hạt bụi!'' Roboco cũng tiếp lời: ``Bọn em còn dọn cả đằng sau kệ tủ rồi.'

Cậu Tổng gật gù, tỏ vẻ như muốn nói: ``Tốt, như thế đạt trên cả yêu cầu rồi đấy.'' Thế nhưng, điều mà nàng Người máy sau đó thốt ra làm cho cậu nhướng mày: ``Nhưng mà ở đằng sau kệ tủ có nhiều bụi với rác quá.''

Kuon nghe vậy chỉ biết cười khì, trong khi Kaien nhìn hai bố con với vẻ mặt không bằng lòng: ``Biết ngay là bố con anh sẽ lại bừa bộn mà.''

Thế rồi, Haato chiêm vào, như thể muốn mách lẻo cả gia đình họ: ``Chưa kể hai tên này còn vô tình phá hỏng đồ đạc trong phòng của hai chú cháu nữa cơ.'' Nghe đến đây, cả gia đình nhà Hikawa đều nhìn họ với con mắt ngạc nhiên. Như một thói quen, cậu Tổng bất giác lớn giọng: ``\textbf{Hai đứa phá cái gì cơ!?}''

``Cái kệ ở trong nhà tắm và cái ổ điện ạ,' nàng Song tính đáp, không muốn dây dưa với hai cô gái được gọi là ``Ngốc nghếch'' trong đội.

Cậu Tổng nhíu mày, giọng pha chút bực bội: ``Hai đứa đại PON này đi phá văn phòng công ty chưa đủ sao mà lại đi phá nhà anh thế? Hả!?''

Đáp lại là một cái giọng lí nhí của Roboco: ``Bọn em đã sửa lại cho anh rồi mà!'' trong khi nàng Hiệp sĩ đứng kế bên cô chỉ biết cười trừ.

Người bố đứng bên cạnh cậu cũng phải vỗ vai bảo cậu bình tĩnh lại: ``Thôi, đừng nóng con. Chắc là họ không cố ý đâu.'' Kuon chỉ có thể ôm trán, lắc đầu: ``Bố thử lên công ty con làm một ngày thì bố sẽ thấy hai tên này quậy như thế nào.''

Để giúp cậu bớt nóng hơn, Haachama đã khéo léo nói thêm: ``Em bắt hai tên ngốc đó phải tự bỏ tiền túi ra để sửa đồ nhà anh rồi đó, Kuon-senpai. Ryujin-san, chú không phải lo đâu ạ.''

Hai cô nàng hậu đậu cũng chỉ nhìn cậu với những con mắt đầy thông cảm: ``Tha cho bọn em đi mà!'' và ``Làm ơn đừng phạt bọn em!''

Cậu Tổng sau khi nghe ba người nói vậy, liền làm nguội cái đầu của mình lại, rồi nói trong tiếng thở dài: ``Ít nhất thì hai đứa cũng có ý tự giác đấy. Đền bù là được rồi. Với cả, đằng nào nhà anh cũng định đi thay cả hai thứ đó mà.''

Hai cô nàng nhìn cậu với những ánh mắt ngạc nhiên, như thể đây là điều mà họ không ngờ tới. ``Thật ạ?''

Cậu đáp: ``Ừ. Anh nhớ cái kệ nhà tắm bị lỏng vít, tấm ván của nó cũng quá cũ rồi, còn ổ điện thì mấy tháng nay cứ chập chờn lúc có lúc không. Hai đứa vừa tạo cái cớ để đi thay đồ mới rồi đấy.''

Cả Roboco và Noel đều thở phào nhẹ nhõm với câu nói này của cậu. Haato vẫn không giấu nối sự ngạc nhiên: ``Vậy\dots\ bọn em vừa thay hộ cho nhà anh ạ?''

Người bố của Kuon cười hiền từ, đáp: ``Có thể nói như vậy. Được rồi, hết bao nhiêu chú sẽ trả. Cảm ơn các cháu nhiều nha.''

Nhìn vẻ mặt tươi vui của ông cùng với khuôn mặt có phần mãn nguyện của cậu, cả ba cô nàng cùng nhau đập tay ăn mừng vì đã hoàn thành nhiệm vụ xuất sắc.

``Được rồi, đội tiếp theo. Subaru-chan và Shion-chan?'' Nàng Phù thủy hào hứng bước tới, theo sau là nàng Vịt với vẻ mặt căng thẳng.

``Em đã dọn sạch nhà bác của anh chỉ trong nháy mắt!'' Shion nói với vẻ rất đỗi tự hào, đến độ có thể cảm nhận được hàng ngàn ngôi sao sáng đang tỏa ra xung quanh cô.

Không mất quá nhiều thời gian để cậu Tổng đọc vị được cô: ``Để anh đoán nha: em dùng ma thuật để dọn dẹp, đúng chứ? Cái tính lười động tay động chân như em, anh còn lạ gì nữa.''

Subaru nhìn cô bạn ưỡn ngực kia, thở dài một cách bực bội: ``Đúng như anh nói đó. Cậu ta dọn sạch sành sanh, không còn một cái gì trong phòng cả!''

Cả ba người trong gia đình lặng một người một lúc lâu. Kuon lên tiếng, giọng hơi run, lông mày giật giật: ``\dots Ý-Ý em là toàn bộ đồ đạc trong phòng bác ấy bị bốc hơi hết, đ-đúng chứ!? Vì nếu thế thì anh chuẩn bị sẵn cả bài `cải lương' để `hát' cho các em đây\dots''

Thấy rằng Kuon lại chuẩn bị nổi giận, nàng Phù thủy vội lên tiếng: ``Bà điêu vừa thôi! Tôi đã đưa hết đồ đạc về như cũ mà!''

Cậu Tổng nghiêm nghị quay sang nhìn Subaru, ánh mắt đầy dò xét: ``Có đúng như vậy không, Subaru-chan?''

Nàng Vịt lúng túng, tiếp tục trả lời cậu bằng một giọng e dè: ``Dạ\dots\ đúng là như vậy ạ. Căn phòng đó sạch sẽ, không có hạt bụi, và tất nhiên đồ đạc còn nguyên vẹn ạ.''

``Đúng như cậu ta nói đấy,'' nàng Phù thủy tiếp lời.

Kuon khẽ thở dài nhẹ nhõm, khuôn mặt giãn ra một chút: ``Thế là may rồi\dots\ Không thì anh chả biết phải ăn nói thế nào với các bác nữa.''

Cậu lắc đầu, nhưng không quên gật nhẹ để cảm ơn đội hai người họ, vì ``dù gì họ cũng đã hoàn thành nhiệm vụ.'' Subaru cảm ơn rối rít, còn Shion nhìn cô với sự khó hiểu.

Ngay khi định chuyển sang dội tiếp theo, đội dọn bếp lập tức bước lên, vẻ mặt họ có chút căng thẳng. Choco là người xung phong lên nói, mặc dù cô tỏ vẻ e ngại: ``Bếp nhà anh cũng đã được dọn rồi đó, Kuon-sama, cơ mà\dots''

Mio rít lên, cả người cô khẽ run rẩy khi bắt phải nhớ lại trải nghiệm kinh hoàng: ``Nhà anh nhiều gián quá!''

Cậu Tổng khẽ rùng mình, không phải vì bất ngờ mà vì xấu hổ, như đã tiên đoán được trước đó: ``\dots Ừ. Xin lỗi vì đã bắt mấy đứa phải làm chuyện này.''

Kaien thậm chí còn cất giọng xóc xỉa tới hai bố con: ``Thế mà bảo bếp nhà mình làm sạch sẽ lắm cơ! Nguyên cả một ổ gián còn không thèm động đến kia\dots''

Ayame đứng ra, tiếp tục báo cáo về việc họ đã làm: ``Thế nên ba đứa bọn em với Mel-senpai đã diệt hộ anh rồi. Nhưng\dots''

Thấy cô ngập ngừng, cậu khẽ nhíu mày, khoanh tay hỏi: ``Mấy đứa lại phá đồ phá đạc gì nữa sao?''

Nàng Y tá lên tiếng: ``Không hẳn là phá đồ đâu, Kuon-sama, chỉ là\dots'' Nàng Sói đen tiếp lời: ``Bọn em vô tình làm vỡ mấy cái bát đĩa trong khi đuổi bắt gián ạ.''

``Và cả chạn treo bát đĩa nữa!'' Ayame nói trong sự khẩn thiết, như thể muốn xin được gia đình cậu rộng lượng.

Người đàn ông trong gia đình sau khi nghe ba cô gái nói, chỉ biết lắc đầu cười gượng. Bình thường, ông sẽ không ngần ngại mà ``cà khịa'' hai đứa con của mình, thậm chí còn lớn tiếng với họ, nhưng lần này ông lại nhẹ nhàng hơn: ``Thôi không sao, đằng nào bác cũng định dỡ nó ra rồi.''

Cả ba cô gái nghiêng đầu. Ông tiếp tục giải thích: ``Lần trước cái chạn bát đũa đó cũng rơi đổ sập xuống giữa đêm, chắc là vì quá cũ, vít bị gãy, cả đống bát đĩa đổ xuống. May mà không ai bị thương. Bác đã định thay nó từ hồi đó mà bận quá.''

Cậu nhớ lại hồi mà cả nhà chạy tán loạn khắp phòng đêm hôm đó: ``Lúc đó con còn tưởng là cái gì nổ cơ\dots'' Mẹ của cậu cũng lên tiếng: ``Em còn tưởng là trộm đột nhập vào\dots''

Mio, Choco và Ayame bật cười với sự kiện khiến cả một gia đình hoang mang như thế. Nàng Y tá cười nhẹ, nói với gia chủ: ``Vậy chắc là nhà cháu vừa làm hộ nhà chú làm một việc rồi nhỉ?''

Cậu bật cười, khẽ thở dài và nhún vai: ``Có thể nói là như vậy. May cho ba đứa nha.''

Mọi người cùng nhau cười nghiêng ngả, xua đi không khí căng thẳng ban đầu. Kuon chỉ lắc đầu, một thói quen không thể bỏ được mỗi khi nhà mình có những câu chuyện ``dở khóc dở cười.''

Chốt hạ lại, cậu hỏi: ``Tóm lại là cái bếp nhà anh dọn sạch hoàn toàn rồi đúng không?'' Cả ba người đồng đoạt gật đầu. Dù thế, cả Sói đen và nàng Y tá đều bày tỏ sự bất bình.

\begin{dialogue}
    \speak{Choco} Kuon-sama, lần sau bọn anh tự dọn lại căn bếp của mình đi.
    \speak{Mio} Bọn em sẽ không giúp anh làm công việc này nữa đâu!
\end{dialogue}

Riêng Ayame thì tỏ vẻ lém lỉnh, tỏ vẻ muốn đề bạt ý kiến: ``Nếu cần thì em có thể trở thành dững sĩ diệt gián cho, da zo!''

Kuon phì cười, khẽ gật đầu: ``Thôi được rồi, khi nào nhà anh cần thì sẽ gọi. Còn về cái chạn nhà bếp thì bọn anh xử lý sau.''

Cậu gọi đội tiếp theo lên, bao gồm Aki và Ryuuhou. Họ nhanh nhảu báo cáo tất cả những gì họ đã làm với đống quần áo trong nhà. Cũng không có gì đáng bàn, khi mà người nhà của cậu, cùng với sự giúp sức của nàng Dị giới, đã hoàn thành công việc khá nhanh.

Ngoại trừ một chi tiết khiến cậu con trai phải cau mày, thậm chí còn khiến mọi người chỉ muốn thốt lên ``ối dồi ôi.''

Kuon đặt chậu quần áo màu xanh lá xuống sàn, mắt nhìn Aki đầy thắc mắc: ``Khoan đã\dots Em tạo ra cái chậu giặt từ cái \emph{góc ấy} sao?''

Nàng Tóc bay có vẻ hơi ngượng ngùng nhưng vẫn đáp lại: ``Nhà anh thiếu chậu mà, em đành phải chế tạm một cái thôi.''

Cậu nhướng mày, giọng đầy chất vấn: ``Em có biết là chỉ cần đi qua bên tầng năm, ở phòng bác Hijaki-san, lấy chậu dự phòng là xong không? Cần gì phải mất công làm mấy cái thao tác `vi diệu' như thế?''

Nghe vậy, Aki mở to mắt, rõ ràng là không hề biết điều này: ``Ể? Thật á? Sao không ai nói với em vậy?'' Cậu thở dài, nhìn cô em gái đang buột miệng cười: ``Em không nhắc gì cho em ấy à, Ryuuhou?''

Cô em gái vẫn đang cười xòa, đáp: ``Em định bảo chị ấy đấy, nhưng mà chị ấy bất ngờ làm \emph{chuyện đó}, làm sao cản kịp được anh!''

Kaien lúc này xen vào, giọng điệu nửa đùa nửa trách: ``Mẹ cũng đã bảo rồi, nhưng chắc em ấy mải nghĩ gì đó mà không nghe. Thôi, cái chậu đó để lại làm kỷ niệm cũng được.''

Mọi người bật cười, và dường như chuyện này nhanh chóng được bỏ qua. Tuy nhiên, có một việc khác không dễ dàng lờ đi. Kuon nhận ra một chiếc áo chiếc áo dài tay thuộc về bố mình bị xé một đường khá dài. Cậu tiếp tục nhướng mày, hỏi: ``Aki-chan, em làm rách cái này sao?''

Nàng Dị giới lúng túng: ``E-Em xin lỗi! Em không cố ý đâu, chỉ là khi vắt nó mạnh tay quá\dots''

Cậu liếc sanng Ryujin, người cũng chỉ cười trừ: ``Lần đầu cháu giặt tay, đúng chứ? Thôi không sao, cái áo này cũ lắm rồi. Chú cũng định bỏ nó đi mà chưa có dịp. Lần này coi như là lý do hợp lý vậy.''

Kaien nhăn mày: ``Anh này, đúng là\dots Cái áo này vẫn còn tốt mà, khâu lại vài đường là xong, anh cứ bỏ\dots''

Ông quay sang người vợ, khẽ lắc đầu: ``Thôi, nó cũng thực sự cũ lắm rồi. Em muốn giữ thì cứ đưa cho thằng Kuon nó mặc.'' Cậu con trai đáp: ``Thế cũng được bố. Đằng nào con cũng đang thiếu áo dài tay.''

Người đàn bà chỉ biết thở dài, chấp nhận lý do có phần hơi khiên cưỡng này từ ông, còn mọi người đều gật gù, cười trừ trước sự dễ dãi đến bất ngờ của ông.

Sau khi đội giặt đồ đã xong, Miko đi tới cậu, cất lên: ``Mà này ne, anh không định hỏi bọn em dọn phòng thờ thế nào à?''

Cậu con trai khẽ cười, tay vắt chéo trước ngực: ``Dọn phòng thờ thì có gì để kể đâu. Chắc cũng chỉ phủi bụi, trải khăn, lau bộ ấm chén chứ gì?''

Tuy nhiên, lời nói tiếp theo của nàng Ma cà rồng khiến không khí thay đổi ngay lập tức: ``Không đâu! C-C-Có ma đấy, Kuon-senpai!'' Cậu nhìn Mel với sự hoài nghi xen lẫn đôi chút ngạc nhiên: ``Ma á? Em chắc là không đùa chứ?''

Nàng Vu nữ lập tức hưởng ứng: ``Thật đó, ne! Lúc bọn em đang trải khăn lên bộ ấm chén, tự dưng thấy hai con ma hiện lên. Chúng nó còn chửi bọn em nữa cơ!''

Cậu lập tức tỏ vẻ nghi ngờ, liền quay sang hỏi ông bố: ``{\color{red} Bố mẹ có biết nhà mình đang thờ ai không ạ?}''

Ryujin và Kaien ngẫm nghĩ một lúc lâu, rồi cùng đồng loạt lắc đầu: ``Không./Không con ạ.'' Người bố còn giải thích thêm: ``Bố nhớ là các cụ nhà mình thờ ở quê mà nhỉ?''

``Nói chung là nhà mình không có thờ ai đâu,'' người mẹ nói, ``chắc lại là mấy đứa chán quá lại bày trò gì đây mà.''

Sau khi nghe phán quyết của hai người, Mel và Miko sững sờ khi họ không tin câu chuyện của họ. Họ định tiếp tục đôi co, nhưng rồi cậu xen vào hạ hỏa: ``Thôi, thôi được rồi. Dọn dẹp sạch sẽ là OK rồi. Quan trọng là để các bác còn thắp hương nữa.''

Cặp đôi cũng chỉ đành tiu nghỉu lùi xuống, người thì phồng má lên, người thì vỗ về an ủi, vì không ai trong ba người họ tin vào điều họ nói cả.

Đội cuối cùng cũng lên báo cáo sau đó, với Rushia là người mở đầu: ``Tầng ba cũng đã được bọn em dọn sạch hoàn toàn!'' Fubuki cũng chiêm vào phụ họa: ``Cả ba bọn em đã sạch không dấu vết luôn\~{}''

Tuy nhiên, Aqua lập tức chen vào, giọng đầy bức xúc: ``Không đúng! Chính Shion-chan là người dọn lại mọi thứ đó, trong khi bọn em thì phải lo xử lý cả đống cốc vỡ và cửa sổ do đám Fandead gây ra á!''

Kuon nhướng mày ngạc nhiên: ``Anh không biết là em còn có thể triệu hồi cả người của mình lên đấy, Ru-chan ạ.'' Lập tức, Aqua tiếp tục trút lên sự bức xúc (và đôi phần sợ hãi) mà tiếp tục khai nhận: ``Vấn đề là đám đó cũng vô dụng lắm luôn đó!''

Nàng Pháp sư ngỡ ngàng: ``Ơ, nhưng mà em triệu mọi người lên để giúp chúng ta mà\dots Aqua-senpai\dots'' Nàng Hầu gái không để cô nói hết, tiếp tục kể lể: ``Còn Fubuki-senpai thì chỉ lau mỗi sàn nhà thôi, mà lau kiểu gì còn làm nó trơn trượt nữa chứ! Khiến bọn em lao đao cả nửa ngày trời.''

Cậu vừa nghe cô thở dài, vừa trông vẻ mặt đáng thương như bị mắng của nàng Pháp sư, cùng với vẻ mặt không can của nàng Cáo trắng. Cậu cố gắng giữ bình tĩnh cô: ``Thôi, thôi được rồi. Đừng nóng nữa, đừng nóng nữa. Anh phải thừa nhận là cái tầng đó có hơi nhiều cốc đĩa với cả đồ thủy tinh thật, nên chắc mấy đứa cũng hơi bất ngờ.''

Thậm chí người bố của Kuon cũng phải lên tiếng thừa nhận: ``Phòng đó cũng lâu rồi không được sử dụng nữa. Chắc là đồ thủy tinh dưới đó cũng cũ lắm rồi. Xin lỗi các cháu nhé.''

Cả ba cô gái khi được nghe hai bố con nhà Hikawa bày tỏ sự cảm thông, đều thở dài nhẹ nhõm.

Fubuki tiếp tục lanh lợi báo cáo: ``Với cả bọn em cũng đã dọn cả tầng hai rồi đó\~{}'' Kuon ngạc nhiên, hỏi lại: ``Ờm, thì\dots\ Tầng ba xong thì cứ coi như xong đi, sao phải dọn cả tầng hai nữa làm gì?''

Cô trả lời đầy hồn nhiên, như đang muốn đổ tội cho em ấy: ``Tại Rushia-chan cũng muốn lau thử tầng dưới thôi mà\~{}''

Lời nói vô tư của cô khiến Aqua nổi đóa, còn Kuon thì khẽ nghiêng đầu khó hiểu:

\begin{dialogue}
    \speak{Aqua} \textbf{Chị đang nói cái gì chứ!?}
    \speak{Kuon} Ờm\dots\ Hả?
\end{dialogue}

Rushia cúi đầu, lí nhí thú nhận: ``Với lại\dots\ em lỡ tay làm hai nhân viên ở tầng hai bất tỉnh rồi\dots''

Câu nói của cô khiến cậu ngạc nhiên: ``Lạ nhỉ\dots\ Anh tưởng giờ này họ đã phải về nhà rồi chứ? Cũng giáp Tết rồi mà nhỉ\dots'' Cậu liếc nhìn ba người họ, rồi khẽ lác đầu: ``Thế này thì anh phải giải thích với họ kiểu gì đây\dots''

Sau một hồi cằn nhằn với bản thân, cậu ũng chỉ biết bó tay, thở dài kết luận: ``Lần sau cứ làm đúng những gì anh giao thôi, được chưa?''

Ba người đồng thanh đáp lớn, mặt của họ giãn ra: ``Rồi ạ\~{}!''

Cuối cùng, đến lượt đội sơn tường lên báo cáo công việc của họ, với căn phòng tầng ba này làm minh chứng. Kuon cùng với các thành viên trong gia đình mình đi lại giữa phòng, tay chống hông, tỏ vẻ hứng khởi những cũng không giấu nổi sự tò mò: ``Anh thấy mấy đứa làm cũng tốt phết đấy nhỉ, nhất là ở tầng ba này. Chắc là mất nhiều công sức lắm hả?''

Korone mỉm cười, lắc đầu đầy hứng khởi: ``Đâu có! Em với Okayu nghĩ ra vài cách thú vị để sơn lại bức tường này mà.'' Cậu lại nhướng mày, lờ mờ đoán được ý mà họ sẽ nói tiếp theo. Cậu hỏi, nửa đùa nửa thật: ``Ồ, thế à? Để anh đoán xem. Có phải mấy đứa tạt sơn ngẫu nhiên lên không?''

Okayu nhếch mép, nửa tự hào nửa trêu chọc: ``Không đâu, bọn em giả lập chơi Splat**n đó!'' Câu trả lời của cô khiến Ryuuhou chỉ biết im lặng vài giây trước khi lên tiếng: ``Có
khác gì tạt sơn ngẫu hứng đâu các chị\dots''

Ryujin bật cười ha hả, vỗ vai nàng Chó: ``Chả trách mà nhìn bức tường trông `nghệ thuật' thế này. À không, phải nói là tạt sơn một cách có tổ chức chứ nhỉ?'' Đến Kuon cũng phải công nhận: ``tạt sơn mà không làm bẩn rèm thế này cũng khá đấy.''

Nàng Chó gật đầu đầy tự hào: ``Bọn em còn sơn lại cả tòa nhà nữa cơ!'' Nghe đến đây, cậu Tổng tròn mắt, sửng sốt: ``Ờ\dots\ Hả? Cả tòa nhà này luôn sao?''

Ông bố cũng bày tỏ sự ngạc nhiên: ``Thảo nào lúc chú về nhà cửa trông mới thế.'' Nàng Mèo cười khúc khích, bổ sung: ``Bọn em còn sơn lại cả tòa bên cạnh nữa\~{}''

Sự kinh ngạc của cậu con trai giờ chuyển thành bối rối: ``Anh có bảo mấy đứa đi sơn lại tòa nhà đâu! Anh chỉ bảo sơn lại phòng thôi mà. Làm thế để làm gì cho khổ? Lại còn tự dưng sơn lại cả nhà bên cạnh nữa chứ!''

Korone hồn nhiên đáp, không hề tỏ vẻ áy náy một chút nào: ``Bọn em chỉ mất khoảng 20 phút là xong mà.'' Cô em gái của Kuon khoanh tay, nhướng mày nhìn cặp đôi Chó-Mèo: ``Chắc lại dùng cách tạt sơn như vừa nãy đúng không ạ?''

Okayu cười tít mắt: ``Đúng luôn\~{}'', đồng thời giơ ngón cái lên như muốn khẳng định cho điều họ nói.

Những câu nói ngây ngô của cả đội khiến cậu con trai không biết nên vui hay buồn. Cậu chép miệng, cố giữ bình tĩnh, phán một câu đơn giản: ``Thôi, xong là xong. Anh khôngnói nhiều nữa.''

Mọi người cùng cười rộ lên, nhưng ánh mắt của Kuon lại hướng ra cửa sổ, nơi những bứctường sáng sủa mới tinh\dots\ với những vệt sơn đủ màu sắc trông chẳng khác nào một tác phẩm trừu tượng. Cậu chỉ biết thở dài, nghĩ thầm: ``\textit{Nghệ thuật thật đấy, nhưng chắc nhà mình sắp nổi tiếng cmn luôn rồi\dots}''

\ct

Cuối cùng, khi các nhóm đã báo cáo xong phần việc được giao, Kuon đứng giữa sảnh lớn, vỗ tay để tập hợp mọi người lại. ``Vậy để anh tổng kết lại: nhà mình đã xong mọi thứ-'' nhưng chưa kịp dứt lời, cả nhóm đã reo hò vang dội: ``Hoan hô!''

Cậu phì cười, giơ tay ra hiệu cho họ bình tĩnh: ``Ngoại trừ cái cầu thang.'' Ryuuhou đứng cạnh cậu, bổ sung: ``Và cả đống kho trên cầu thang nữa ạ.''

Tiếng ``Ể\~{}!?'' vang lên đồng thanh từ các cô gái, tràn đầy sự ngỡ ngàng lẫn mệt mỏi. Trong đó, Fubuki phản ứng mạnh mẽ nhất, giậm chân thở dài: ``Còn nữa hả?'' như thể cô chỉ muốn nghỉ ngơi ngay bây giờ.

Mio liếc nhìn cô bạn, thở dài: ``Tôi nghe nói bà còn chưa làm được cái trò sói gì đâu đấy\dots'' Aqua và Rushia nhìn cô nàng Cáo trắng bằng những ánh mắt tỏ vẻ không bằng lòng.

Kuon nhún vai, cười trừ: ``Mấy đứa quên mất à? Ta còn phải tổng vệ sinh cho toàn bộ nhà nữa. Với cả xong cái đó là ta có thể nghỉ được rồi.''

Cô em gái của Kuon lên tiếng động viên: ``Em biết các chị giờ cũng thấm mệt, nhưng ta cùng cố gắng thêm chút nữa đi ạ!''

Tất cả mọi người nhìn nhau, rồi cũng đành gồng mình lên cố gắng hoàn thành nốt nhiệm vụ cuối cùng trong ngày cận Tết.

Nàng Hầu gái lên tiếng, vẻ đầy nghi hoặc: ``Vậy chúng ta sẽ làm thế nào?''

Cậu lắc đầu, đáp với giọng lười biếng: ``Bình thường anh sẽ cầm tay chỉ việc bảo mấy đứa làm cái nọ cái kia như ban nãy, nhưng\dots\ thôi. Cứ để mọi người tự xử.''

Haato cau mày: ``Nhưng cụ thể là những công việc nào?''

Ryujin giải thích ngắn gọn: ``Các cháu cứ quét cầu thang, lau thành cầu, đổ rác ở trên các tầng, rồi dọn dẹp mấy cái kho trên cầu thang thôi.''

Cậu Tổng thêm vào: ``Nói chung là làm thế nào thì làm, anh không quản nữa. Miễn là đừng phá nhà của anh như Robo-PON với Noel-PON là được.''

Kaien khích lệ thêm, giọng điệu pha chút hài hước: ``Với cả nhiều người thế này, mỗi người một tí, thế là xong ngay mà. Các cháu đồng ý không?''

Tất cả các cô gái của \hll\ đồng thanh rất lớn: ``\textbf{Đồng ý!!}''

Rồi nhanh chóng phân chia công việc: người thì nhặt rác, người quét cầu thang, người lau thành lan can.

Tuy nhiên, như thường lệ, những tình huống hài hước lại xảy ra.

\ct

Noel đang cẩn thận quét dọn cầu thang từ tầng sáu xuống thì để ý thấy một mảng tường gần đó bị mục. Cô quay sang hỏi nàng Thỏ đang đi ngang qua: ``Mình có thể gạt vữa tường được không?''

Pekora gật đầu, đáp không mấy để tâm: ``Tường trông cũng bị mục rồi đó, peko. Cứ cạo ra đi.''

Không nói không rằng, nàng Hiệp sĩ liền rút ra cây chùy quen thuộc của mình và đập một cú thật mạnh. *RẦM!* Một mảng tường lớn nứt toác, gạch vụn văng tung tóe. Nàng Thỏ hoảng hốt, vội ngăn lại: ``Cạo như thế là quá rồi đó, peko! Bà đang định san ủi luôn nhà anh ấy sao!?''

Nàng Hiệp sĩ nhìn mảng tường vừa rơi, ngượng nghịu: ``Ừm\dots\ tôi hơi quá tay một chút\dots''

\ct

Hay, ở khu vực cầu thang tầng ba, Mio, Fubuki và Mel đang cặm cụi lau thành lan can. Nàng Cáo trắng thì luôn miệng than thở: ``Đến bao giờ thì ta mới có thể lau xong thành cầu thang nhà anh ấy đây\~{}'' còn bạn của cô thì bày tỏ khó chịu: ``Từ nãy tới giờ bà gần như đã làm gì đâu. Đừng có than vãn nữa và lau nhanh lên!''

Đến cả nàng Dơi cũng bày tỏ sự băn khoăn: ``Cầu thang nhà anh ấy đúng là dài thật đấy. Giá như ta có cách nào đó tiện lợi hơn nhỉ\dots''

Công việc có vẻ chậm chạp và vất vả, khiến Roboco nảy ra một ý tưởng: ``Sao không đổ nước thẳng lên đây luôn nhỉ? Sạch nhanh lắm!''

Nàng Sói đen đã lắc đầu phản đối ngay: ``Như thế làm sao sạch được! Còn phải kỳ cọ nữa chứ!'' Nhưng, nàng Cáo trắng với bản tính lười biếng của mình, lại hưởng ứng nhiệt liệt: ``Ý hay đấy! Ta làm liền!''

Mặc cho Mio và Mel cố ngăn cản, Roboco và Fubuki đã nhanh tay lấy vòi nước áp suất cao (từ bàn tay người máy) xịt thẳng vào lan can. Nước bắn tung tóe khắp nơi, khiến Mel và Mio bị ướt sũng từ đầu đến chân. Bản thân họ ít nhiều cũng bị ``dính đạn'' nước từ chúng.

Nàng Ma cà rồng kêu lên, tỏ vẻ bất mãn: ``Dừng lại đi! Ướt hết cả rồi!'', trong khi nàng Sói đen đang run cầm cập vì lạnh: ``Hai người bị ấm đầu rồi hay sao!?''

Dù thế, nàng Người máy vẫn cố gắng biện minh: ``Nhưng mà dùng vòi nước áp suất cao cũng sạch mà-''

Bỗng, *RẮC!* một tiếng. Tiếng động chát chúa vang lên, kéo theo ánh mắt kinh hãi của cả nhóm. Một vết nứt dài xuất hiện trên thành cầu thang. Roboco ngượng ngùng tắt vòi nước, cười gượng: ``Chắc là\dots\ mạnh tay quá rồi\dots\''

Mio giận dữ hét lên: ``Mạnh cái con khỉ! Mau tìm cách che vết nứt đi, kẻo anh ấy phát hiện ra bây giờ!'' Fubuki thì phá lên cười, càng khiến cho bạn của cô tức giận hơn: ``Còn cái bà này nữa!''

Nàng Người máy chỉ biết gãi đầu xin lỗi, còn nàng Dơi chỉ đứng đó, mặt đỏ bừng, cố lau nước trên người.

\ct

Rushia và Shion được giao nhiệm vụ thông đường ống nước nóng. Sau khi bàn bạc, cả hai quyết định phân công: Shion sẽ đứng trong kho cầu thang để dùng ma pháp ``tăng áp lực nước'', còn Shion sẽ kiểm tra tình trạng đường ống từ phòng tắm của nhà Kuon.

Nàng Pháp sư hét lớn từ xa: ``Chị chắc là làm được không vậy, Shion-senpai? Nhớ cẩn thận nhé, đừng có quá tay!''

Nàng Phù thủy đáp: ``Rồi! Cứ để chị! Shion-sama đây là phù thủy đại tài mà!'' một cách chắc nịch. ``Giờ thì, xem đây!''

Cô bắt đầu niệm thần chú, một làn sáng xanh phát ra từ tay cô, len lỏi vào hệ thống ống nước. Từ bên phòng tắm, Rushia thấy vòi nước rung lên bần bật, nước bắt đầu chảy, ban đầu chỉ nhỏ giọt.

``Được rồi! Thêm chút nữa đi!'' Cô đồng la lên, vẻ hứng khởi. Đối phương nghe thấy, liền tăng cường độ phép thuật: ``\textbf{Lên nước mạnh!}''

*TÁCH!* Tiếng van nước trong kho cầu thang vang lên một cách đáng ngờ. Rồi đột nhiên, nước từ đường ống phun mạnh ra như suối, bắn tung tóe khắp nơi. Shion khựng lại, hoảng hốt: ``Ể-Ể!? Chuyện gì đang xảy ra vậy!?''

Bên trong phòng tắm cũng không khá khẩm hơn, khi Rushia thấy vòi nước phun ra như mưa lũ. Cô hét toáng lên: ``\textbf{S-Shion-senpai!! Chị làm cái gì thế!? Nước chảy lên láng hết rồi này!!}''

Shion cuống cuồng đáp: ``C-Chị chỉ làm như em bảo thôi mà! Này\dots\ nước trong kho tràn hết ra rồi! Chị không bịt nổi nữa đâu!''

Rushia đáp lại: ``Thế thì khóa van đi chứ!!'' Nàng Phù thủy bối rối, nhìn cái van nước đang rung lắc điên cuồng, rồi hét lại: ``Nhưng chị không biết khóa kiểu gì! Nó quay ngược lại rồi!''

Cô cố hết sức xoay van, nhưng càng xoay, nước lại càng phun mạnh hơn. Kho nhanh chóng ngập nước, biến thành một ``bể bơi'' mini. Nàng Pháp sư vội vàng chạy sang, vừa cười vừa thở dài: ``Trời ơi chị ơi! Chị vừa làm cái gì thế-nanodesu!?''

Nhìn kho nước ngập gần tới đầu gối, Shion buông một câu đầy châm biếm: ``Có khi nhà anh ấy đổi nghề mở hồ bơi luôn đi cho rồi!'' Còn Rushia chỉ biết ngượng đầu, không biết phải giải quyết chuyện này ra sao.

Cả hai loay hoay thêm một hồi mới dừng được dòng nước, nhưng hậu quả là một phần kho cầu thang đã bị ``tái định nghĩa'' thành vùng đất ngập lụt, đồng thời báo hại luôn cho các thành viên đang miệt mài lau dọn cho cầu thang phía nhà cậu Tổng.

\ct

Miko, Subaru và Matsuri nhận nhiệm vụ cắt tỉa cây ở ban công tầng sáu.

Nàng Vu nữ đang chăm chú cắt từng nhánh cây, Subaru thì gom cành lá rụng vào bao tải, còn Matsuri đứng đó vừa làm vừa hát nghêu ngao.

Miko nhắc nhở, tay vẫn thoăn thoắt cắt tỉa: ``Subaru-chan, nhớ gom cẩn thận ne, đừng để rơi lá ra sàn!'' Nàng Vịt đáp lại, tỏ vẻ nghiêm túc: ``Thì tôi đang làm này! Đừng có càu nhàu nữa!''

Nàng Cổ động nghe vậy liền cười lớn: ``Có gì phải sồn sồn thế? Chỉ là nhặt lá rơi thôi mà.''

Đúng lúc đó, Miko chợt nhận thấy một chậu cây lớn trên lan can. Cô nghiêng đầu thắc mắc: ``Ủa? Cái chậu này sao lại ở đây nhỉ? Hay là cắt bớt cành đi cho gọn nhỉ?''

Subaru nghe vậy, liền thốt lên cảnh báo: ``K-Khoan đã! Đừng đụng vào nó trước khi-''

Nhưng đã quá muộn. Với độ PON bất bại, nàng Vu nữ cắt một nhát gọn lỏn vào cành cây to đang đỡ chậu cây. Chậu mất điểm tựa, đổ nghiêng và rơi xuống, vỡ toang ngay dưới sàn.

Nàng Cổ động nhìn cảnh đó, thốt ên: ``C-Chết rồi, cái gì đây-'' Từ đống vỡ vụn của chậu cây, một đàn ong đông nghịt bay vút ra, vo ve dữ dội. Cả ba người hoảng hồn, đứng như trời trồng.

Nàng Vịt la lớn, vừa chạy vừa vung tay loạn xạ: ``\textbf{Ối! Là tổ ong kìa!}'' Matsuri hét lớn: ``Chạy mau! Đừng để chúng đốt!'' rồi phóng xuống nhà, kéo theo nàng Vu nữ vẫn còn đang ngơ ngác không hiểu chuyện gì đang xảy ra.

Đàn ong, như thể bị chọc giận, bám sát ba người họ không rời. Cả nhóm vừa la hét vừa chạy quanh tầng sáu, làm náo loạn cả khu vực.

Đang khi tuyệt vọng không biết trốn đi đâu, nàng Tomboy bỗng phát hiện dòng nước đang chảy lênh láng từ kho dẫn xuống tầng dưới. Không chần chừ, cô hét lớn: ``Nhảy xuống nước mau! Ong ghét nước!''

Cả ba người lao xuống khu vực nước ngập, vừa té xuống, vừa lặn xuống. àn ong bay quanh một lúc rồi bỏ đi, để lại khung cảnh hỗn loạn với nước chảy tràn lan.

Họ ngoi lên mặt nước, rồi thở phào nhẹ nhõm. ``C-Có vẻ như đàn ong đã đi rồi đấy\dots''

Matsuri, người vừa gạt nước khỏi tóc, vừa than vãn: ``Nhờ dòng nước này mà thoát, nhưng cũng mất sạch hình tượng idol rồi!'' Subaru nhíu mày: ``Chúng ta có bao giờ là idol đâu mà\dots''

Miko thì ngỡ ngàng, như thể cô vẫn đang ``treo máy'' sao biến cố vừa rồi. Một hồi sau, cô tỉnh lại, bây giờ mới hiểu được mọi chuyện. ``C-Chị không biết trong chậu đó có tổ ong ne\dots''

Subaru, nhìn cả đám ướt nhẹp, thở dài: ``Miko-chi, lần sau làm gì thì nhớ để ý chút nhé\dots Cơ mà, ướt thế này cũng không tệ lắm đâu!''

\ct

Okayu và Korone nhận nhiệm vụ dọn dẹp tủ quần áo trong phòng khách nhà Kuon.

Ngay khi hai người họ mở tủ ra, họ đã được chiêm ngưỡng tủ quần áo của nhà cậu đang\dots chất thành núi, với áo này áo nọ được gấp khá cẩu thả, áo không xếp ra áo, quần cũng không xếp ra quần.

Nàng Miêu nhẹ nhàng cất lời: ``Koro-san, đồ của gia đình anh ấy có vẻ khá bừa bộn nhỉ?'' Nàng Khuyển gật gù: ``Ừ\dots\ Tôi cảm thấy có đồ đã không được mặc mấy năm nay rồi đấy. Tôi còn ngửi thấy có mùi, chắc là cần dọn kỹ hơn.''

Nói rồi, Korone đưa tay kéo ra vài cái áo cũ. Bất ngờ, từ bên trong tủ, mấy con gián bò ra, lượn lờ khắp nơi.

``Gyaaah!!! Gián!!'' cô hét lên, vứt đống quần áo vừa cầm xuống đất. Okayu lùi lại, tỏ vẻ hơi tái mét: ``Cái gì thế này!? Sao trong tủ lại có gián chứ!?''

Hai người hoảng loạn, chạy vòng quanh phòng khách, cố gắng tránh xa những vị khách không mời mà đến. Korone vừa la hét, vừa chỉ tay tới đống gián đang bò: ``O-Ogayu!! Làm gì đi chứ! Đuổi chúng ra đi!''

Okayu cũng đang hoảng hốt không kém: ``Sao lại là tôi!? Bà làm đi chứ!''

``Không đời nào! Bà mạnh mẽ hơn tôi mà!!'' nàng Khuyển la lớn, tỏ ý phản đối.

Đang lúc hai người tranh cãi, một con gián bất ngờ bay vọt lên, khiến cả hai hét toáng lên và nhảy tọt lên trên nóc của tủ quần áo.

Đúng lúc đó, con chó Mi - cô chó kiểng màu trắng nhỏ nhắn của nhà Kuon - từ đâu chạy vào, sủa một cách loạn xạ. Nó nghiêng đầu ngước nhìn cả hai cô nàng với vẻ thắc mắc, rồi ngó thấy mấy con gián đang bò trên sàn.

Korone lập tức vỗ tay reo lên: ``Mi-chan! Cứu bọn chị với!'' Okayu liền hùa theo: ``Phải đó, làm ơn đuổi hết lũ gián này đi!''

Như hiểu được tình hình, con ``cờ hó'' nhà Kuon sủa vài tiếng (điều mà nó rất hiếm khi làm), rồi bắt đầu truy đuổi những con gián. Chỉ trong vòng vài phút, chú chó nhỏ đã khiến lũ gián bỏ chạy tán loạn, không còn con nào dám bén mảng lại gần.

Sau khi không còn bóng dáng của một con ``tiểu cường'' nào nữa, cặp đôi mới thở phào nhẹ nhõm, từ từ xuống khỏi nóc tủ.

Korone vỗ tay khen ngợi: ``Mi-chan giỏi quá! Đúng là siêu anh hùng của ngày hôm nay!'' Okayu xoa đầu Mi, cười trừ: ``Nếu không có Mi-chan, chắc chúng ta còn ngồi đây đến tối mất á\~{}''

Đáp lại, con chó kiểng chỉ nằm dài trên sàn, thở phì phò như vừa hoàn thành một nhiệm vụ khó khăn.

\ct

Choco, Ayame, Sora và Aki cùng nhận nhiệm vụ thu gom rác từ các phòng để mang xuống dưới tầng đem đi đổ.

``Rồi, mỗi người phụ mình một tay nhé!'' Sora vừa cười vừa xắn tay áo, bắt đầu gom túi rác từ phòng bếp. Aki thì lúi húi buộc gọn túi rác từ phòng khách: ``Chỗ này cũng gần xong rồi. Choco-sensei, cô gom rác ở phòng bên kia xong chưa ạ?''

``Xong rồi đây, rồi đây\~{}'' Choco bước ra từ phòng tắm với một túi rác lớn, khuôn mặt thoáng vẻ khó chịu: ``Mùi kinh quá đi mất\dots\ Cái này chắc để lâu rồi.''

Ayame bước tới với dáng vẻ tự tin, kéo theo một túi rác từ phòng ở tầng dưới: ``Ha! Chỉ là rác thôi mà. Không có gì phải sợ cả, đúng không?''

Nhưng ngay lúc đó, Aki kéo hộc tủ đựng rác trong bếp ra, bất ngờ một bầy gián nhỏ từ đâu lao ra, bò lổm ngổm khắp sàn.

Nàng Dị giới la hét lên, lùi lại phía sau. Nàng Thần tượng cũng giật bắn mình, ôm chặt túi rác (khô) vào người: ``\textbf{G-Gián!! Nhiều gián quá!!}''

Tình hình của nàng Y tá thì còn thảm hơn, khi cô lại phải trải nghiệm cảm giác kinh hoàng này lần thứ hai. Cô núp ngay sau lưng nàng Quỷ sừng, người vừa tự phong là ``dũng sĩ diệt gián'' ban nãy: ``Ayame-chan, cứu bọn cô với!''

Nhưng khác với mọi người, nàng Oni vẫn bình thản. Cô hắng giọng, rút từ trong túi ra\dots\ một chiếc dép, thay vì là hai thanh katana dắt trên lưng! Cô hét lên: ``Mấy con gián kia, hãy chuẩn bị đối mặt với Ayame-sama!!''

Không chờ thêm giây nào, cô lao thẳng vào bầy gián, dùng dép quật xuống từng con một với động tác dứt khoát. Những con gián không bị đập trúng cũng hoảng loạn bò lung tung, khiến tình hình càng hỗn loạn hơn.

Nàng Quỷ sừng hùng hồn tuyên bố: ``Tránh ra! Ta sẽ quét sạch lũ gián này!!'' và tiếp tục đập bọn chúng một cách không ngừng nghỉ.

Nàng Dị giới đứng từ xa, vừa run rẩy, vừa vỗ tay tán thưởng: ``A-Ayame-chan giỏi thật đấy\dots\ Cứ như là anh hùng vậy!''

Trong khi đó, Choco vẫn đang nấp sau lưng Sora, hai người bắt đầu khẽ khóc: ``T-Tại sao chuyện này lại xảy ra như vậy chứ\dots''

Cuối cùng, sau một hồi vật lộn, Ayame cũng đập được con gián cuối cùng. Cô đứng thở hổn hển, nhưng nụ cười rạng rỡ hiện lên trên gương mặt: ``Xong rồi! Không còn con nào nữa!''

Sora bước tới, tỏ vẻ ngưỡng mộ nàng Oni ấy: ``Ayame-chan đúng là dũng sĩ diệt gián! Cảm ơn em nhiều!'' Choco thì rón rén bước tới, nhìn đống xác gián trên sàn, rồi quay sang nàng Quỷ sừng: ``Em\dots\ đúng là đáng sợ theo cách của riêng mình đấy. Nhưng cô vẫn cảm ơn.''

Ayame cười khì, gãi đầu: ``Không có gì đâu! Đây là chuyện nhỏ với Ojou-san này thôi mà.''

Trong lúc mọi người còn đang hân hoan, từ một góc tủ, một con gián nữa bất ngờ bò ra. Aki hét lên: ``\textbf{L-Lại nữa kìa!!}''

Ayame nhíu mày, nhưng trước khi cô kịp ra tay, con gián đã nhanh chóng bò vào túi rác mở của Choco mà cô đang xách trên tay.

``Gyaaaah!!! Gián chui vào túi rác rồi!!'' nàng Y tá hét lớn, quăng luôn túi rác ra hành lang cầu thang. Đến đây, nàng Oni vỗ trán, thở dài: ``Thôi được rồi, lần này ta lại ra tay vậy\dots''

Mọi người đứng từ xa, nhìn Ayame tiếp tục thể hiện vai trò ``dũng sĩ diệt gián'' của mình, lòng thầm cảm phục nhưng cũng không khỏi rùng mình.

Đến bản thân Ayame cũng là một người không chịu được khi phải nhìn thứ sinh vật gớm ghiếc này, nhưng có lẽ vì cô muốn ra oai với mọi người, nên cô phải gồng mình làm vậy. Thực tế ra mà nói, bản thân Ojou-san cũng là một người\dots\ khá sợ ``tiểu cường''.

\ct

Dần dần, mọi người đã hoàn thành xong công việc. Tiếng cười nói, tiếng đồ vật được sắp xếp ngăn nắp, cùng những sự việc ``giời ơi đất hỡi'' đã diễn ta và cả âm thanh vui vẻ của những bước chân đi lại vang lên khắp nhà. Đúng như kỳ vọng, căn nhà rộng lớn giờ đây đã trở nên sáng sủa, sạch sẽ, chuẩn bị sẵn sàng để đón Tết.

Họ cùng nhau tụ họp tại tầng ba để nghỉ ngơi. Những khuôn mặt đỏ bừng vì mệt, nhưng ánh lên sự hân hoan. Ryuuhou và Kuon mang đến một thùng nước ngọt, vừa khệ nệ đặt xuống vừa nói: ``Như vậy tập thể chúng ta đã làm xong rồi đấy. Mọi người uống chút nước đi.'' ``Anh bao.''

Họ phát cho mỗi người một lon nước lạnh, được lấy ra từ chiếc thùng mà cậu đã mua về trước đó. Tất cả nhận lấy, mở nắp và đồng thanh nói lời cảm ơn.

Sora thở hắt ra, dựa lưng vào ghế: ``Ăn Tết đúng là vất vả thật\dots'' Cậu Tổng mỉm cười, nói: ``Nhưng như thế mới là Tết, đúng chứ? Khổ trước sướng sau thôi.''

Ryujin gật gù nhìn các cô gái: ``Đó mới chỉ là dọn nhà thôi. Các cháu vẫn còn phải học tập ở đây nhiều.''

Subaru, người đang nằm dang tay trên sàn, nhổm dậy với đôi mắt lấp lánh đầy mãn nguyện: ``Nhưng mà làm việc tập thể cũng vui mà!''

Subaru, người đang nằm dang tay trên sàn, nhổm dậy với đôi mắt lấp lánh đầy mãn nguyện: ``Nhưng mà làm việc tập thể cũng vui mà!''

Cả phòng bật cười trước màn đấu khẩu hài hước của hai cô gái. Kaien đứng bên cửa sổ ngắm nhìn khoảng sân trước nhà, nơi những ánh nắng chiều nhẹ nhàng chiếu xuống, nói với giọng trầm ấm: ``Cô nghĩ cả nhà ta sẵn sàng đón một cái Tết rồi nhỉ?''

Mọi người đồng thanh, đầy hào hứng: ``Dạ đúng ạ!'' Cậu con trai đứng dậy, nâng cao lon nước như thể đang dẫn dắt cả một đội quân: ``Đây mới chỉ là khởi đầu thôi! Việc nặng cũng đã xong rồi, đã đến lúc ta hưởng trái ngọt!''

Choco bật cười, nheo mắt trêu đùa: ``Anh nói như thể chúng ta vừa mới hoàn thành một nhiệm vụ siêu khó đấy, Kuon-sama ạ\dots'' Haato, đang ngồi xếp bằng, gật đầu hùa theo: ``Nghe giống triết lý ghê ha\dots''

Kuon giả vờ không để ý đến những lời trêu chọc, tiếp tục với vẻ mặt nghiêm túc: ``Anh hy vọng mọi người sẽ có một cái Tết đầm ấm nhất trong cuộc đời!'' A-chan, tay cầm lon nước, cười nhẹ nhàng nhưng không quên nhắc nhở: ``Và cả dữ kiện làm phim nữa. Đừng quên chúng ta đến đây với mục đích gì!''

Cậu con trai bật cười, quay sang nhìn cô đồng nghiệp của mình: ``À đấy, nói đến phim thì chị đã quay được gì chưa?''

Nàng Đeo kính nháy mắt, tỏ vẻ đầy tự hào: ``Tôi đã sử dụng vài máy quay siêu mini để quay lại hành động của mọi người rồi. Phát minh của cậu cả đấy, Hikawa-san à.''

Cậu Tổng cười, vỗ nhẹ tay lên đầu mình, nói đầy ẩn ý: ``Đó mới chỉ là một phần thôi chị. Vẫn còn nhiều thứ nữa chị sẽ cần phải quay đấy.'' Đồng nghiệp của cậu liếc nhìn cậu, giọng hóm hỉnh nhưng không kém phần quyết tâm, thốt ra một câu đùa mà xứ Etsunan hay nói: ``Đương nhiên là thế rồi. Không thì chún ta móm à!''

Cả căn phòng bật cười lớn trước câu nói đùa dí dỏm của cô. Không khí ấm áp, thân tình bao trùm khắp nơi, khiến mọi người cảm thấy như đã thực sự gắn bó với nhau hơn bao giờ hết.

Buổi chiều 29 Tết, một buổi chiều đặc biệt với gia đình Hikawa và những cô gái nhà Tam Giác Xanh, không chỉ đơn thuần là chuẩn bị đón Tết Nguyên đán truyền thống mà còn là dấu mốc cho một cái Tết rất đặc biệt. Một cái Tết đầy tiếng cười, sự sẻ chia, và có lẽ, sẽ khó lòng lặp lại lần thứ hai trong đời họ.

\end{document}